% Motion in a circle with uniform speed — Pure Mathematics with Mechanics Upper Sixth — Cours signé OnBuch+
% Official MINESEC syllabus. Compiler avec : tectonic PMM12-motion-in-a-circle-with-uniform-speed.tex
\newcommand{\DOCMATIERE}{Pure Mathematics with Mechanics}
\newcommand{\DOCNIVEAU}{Upper Sixth}
\newcommand{\DOCMODULE}{Upper Sixth — Pure mathematics with mechanics}
\newcommand{\DOCLECON}{12}
\newcommand{\DOCTITRE}{Motion in a circle with uniform speed}
% OnBuch+ — préambule « cours signés » ENRICHI (compilé avec tectonic / XeLaTeX)
% Système visuel « Pop » d'OnBuch+ : fond crème (jamais blanc), bordures encre
% épaisses, ombres « dures » décalées, titres Archivo Black en capitales.
% Version MATHÉMATIQUES : graphes (pgfplots/tikz), boîtes pédagogiques
% (définition, propriété, méthode, attention, le savais-tu, exercice résolu,
% exercice, TP, bilan), page de garde et signature OnBuch+ ancrée en bas.
%
% Macros attendues dans main.tex AVANT \input{preamble} :
%   \DOCMATIERE  \DOCNIVEAU  \DOCMODULE  \DOCLECON (numéro)  \DOCTITRE
\documentclass[11pt]{article}

\usepackage{fontspec}
\usepackage[english]{babel}
\usepackage[a4paper,margin=2cm,top=3.4cm,bottom=2.8cm,headheight=1.4cm,footskip=1.3cm]{geometry}
\usepackage{amsmath,amssymb,amsfonts}
\usepackage{mathtools} % \xmapsto, \coloneqq... souvent utilisés par les modèles
\usepackage{cancel} % simplifications barrées (bilans de forces)
% \abs{z} (module, valeur absolue) et \norm{u} (norme), utilisés par les modèles
\providecommand{\abs}{}\let\abs\relax\DeclarePairedDelimiter{\abs}{\lvert}{\rvert}
\providecommand{\norm}{}\let\norm\relax\DeclarePairedDelimiter{\norm}{\lVert}{\rVert}
\usepackage[table]{xcolor} % [table] : fournit aussi \rowcolors (alternance de lignes)
\usepackage{pagecolor}
\usepackage{tikz}
\usetikzlibrary{arrows.meta,positioning,calc,shapes.geometric,shapes.misc,shapes.arrows,decorations.pathmorphing,decorations.pathreplacing,decorations.markings,patterns,shadows,mindmap,trees,fit,backgrounds,matrix,intersections,angles,quotes}
\usepackage{pgfplots}
\usepackage[version=4]{mhchem} % équations nucléaires \ce{^{235}_{92}U}
\usepackage[european,siunitx]{circuitikz} % schémas électriques
\pgfplotsset{compat=1.18}
\usepgfplotslibrary{fillbetween,groupplots} % aires entre courbes (calcul intégral)
\usepackage{enumitem}
\usepackage{fancyhdr}
% [most] charge toutes les bibliothèques tcolorbox (listings, minted, raster,
% documentation...) et gonfle inutilement la mémoire TeX sur un document long
% (~300 boîtes) : on ne charge que ce qui est réellement utilisé.
\usepackage[skins,breakable]{tcolorbox}
\usepackage{graphicx}
\usepackage{colortbl}
\usepackage{booktabs}
\usepackage{tabularx}
\usepackage{diagbox} % case d'en-tête diagonale des tableaux à double entrée
% Colonnes extensibles centrée / gauche / droite (C, L, R), souvent utilisées par les modèles
\newcolumntype{C}{>{\centering\arraybackslash}X}
\newcolumntype{L}{>{\raggedright\arraybackslash}X}
\newcolumntype{R}{>{\raggedleft\arraybackslash}X}
\usepackage{array}
\usepackage{multicol}
\usepackage{siunitx}
% parse-numbers=false : accepte aussi \SI{2,5 \times 10^{-3}}{...}
\sisetup{locale=UK,output-decimal-marker={.},group-minimum-digits=4,parse-numbers=false}
\DeclareSIUnit\ohm{\ensuremath{\symup{\Omega}}} % U+2126 absent de la police
\DeclareSIUnit\division{div} % réglages d'oscilloscope (ms/div, V/div)
\DeclareSIUnit\tr{tr} % tours (vitesse de rotation, ex. \SI{1500}{\tr\per\minute})
\usepackage{qrcode}
% \degree : commande fréquemment utilisée par les modèles pour un angle ou une
% température en dehors de \SI{}{} (non fournie par les paquets chargés).
\providecommand{\degree}{^{\circ}}
% \qed (fin de démonstration), utilisé par les modèles sans amsthm
\providecommand{\qed}{\hfill$\square$}
% \barre : double barre des valeurs interdites dans un tableau de variations (tkz-tab)
\providecommand{\barre}{\ensuremath{\|}}
% environnement proof (amsthm non chargé) utilisé par les modèles
\ifdefined\proof\else\newenvironment{proof}[1][Proof]{\par\noindent\textit{#1.}\ }{\qed\par}\fi
\usepackage{lastpage}
\usepackage{hyperref}
\hypersetup{hidelinks}

% ============================================================
% Palette « Pop » OnBuch+ — mêmes hex que la classe P (plus_ui.dart)
% ============================================================
\definecolor{poporange}{HTML}{F59321}
\definecolor{popdark}{HTML}{E07A0C}
\definecolor{popcream}{HTML}{FFF6E8}
\definecolor{popink}{HTML}{1C1714}
\definecolor{poppurple}{HTML}{7A5AE0}
\definecolor{popgreen}{HTML}{2E9E6A}
\definecolor{poppink}{HTML}{E0457B}
\definecolor{popblue}{HTML}{2F7DE1}
\definecolor{popgold}{HTML}{C98A12}
\definecolor{popmuted}{HTML}{8A7F76}
\definecolor{popline}{HTML}{EADFD2}
\definecolor{popgray}{HTML}{8A7F76}
\colorlet{popgrey}{popgray}
% Alias tolérants (le modèle nomme parfois les couleurs différemment)
\colorlet{popwhite}{white}
\colorlet{popblack}{popink}
\colorlet{popred}{poppink}
\colorlet{popyellow}{popgold}
% Teintes claires pour fonds de tableaux / zones
\colorlet{popblueL}{popblue!10!white}
\colorlet{poporangeL}{poporange!14!white}
\colorlet{popgreenL}{popgreen!12!white}
\colorlet{poppurpleL}{poppurple!12!white}
\colorlet{poppinkL}{poppink!12!white}
\colorlet{popgoldL}{popgold!14!white}

\pagecolor{popcream}

% ============================================================
% Typographie
% ============================================================
\defaultfontfeatures{Ligatures=TeX}
\setmainfont{PlusJakartaSans-Medium.ttf}[Path=../../../fonts/,BoldFont=PlusJakartaSans-Bold.ttf]
% Maths en sans-serif (Fira Math) avec chiffres et lettres droites de Plus
% Jakarta, pour que les formules se fondent dans le texte.
\usepackage[math-style=ISO,bold-style=ISO]{unicode-math}
\setmathfont{FiraMath-Regular.otf}[Path=../../../fonts/]
\setmathfont{PlusJakartaSans-Medium.ttf}[Path=../../../fonts/,range={up/num,up/{Latin,latin},\mathup}]
% (icomma retiré : décimales avec point en anglais)
% Caractères mathématiques Unicode absents des polices de texte (Plus Jakarta,
% Archivo Black) : rendus par leur équivalent mathématique, en texte comme en
% formule — sinon ils s'affichent en carré vide (ex. « Arithmétique dans ℤ »).
\usepackage{newunicodechar}
\newunicodechar{ℝ}{\ensuremath{\mathbb{R}}}
\newunicodechar{ℤ}{\ensuremath{\mathbb{Z}}}
\newunicodechar{ℂ}{\ensuremath{\mathbb{C}}}
\newunicodechar{ℕ}{\ensuremath{\mathbb{N}}}
\newunicodechar{ℚ}{\ensuremath{\mathbb{Q}}}
\newunicodechar{𝔻}{\ensuremath{\mathbb{D}}}
\newunicodechar{α}{\ensuremath{\alpha}}
\newunicodechar{β}{\ensuremath{\beta}}
\newunicodechar{γ}{\ensuremath{\gamma}}
\newunicodechar{δ}{\ensuremath{\delta}}
\newunicodechar{ε}{\ensuremath{\varepsilon}}
\newunicodechar{θ}{\ensuremath{\theta}}
\newunicodechar{λ}{\ensuremath{\lambda}}
\newunicodechar{μ}{\ensuremath{\mu}}
\newunicodechar{σ}{\ensuremath{\sigma}}
\newunicodechar{τ}{\ensuremath{\tau}}
\newunicodechar{φ}{\ensuremath{\varphi}}
\newunicodechar{ψ}{\ensuremath{\psi}}
\newunicodechar{ω}{\ensuremath{\omega}}
\newunicodechar{Σ}{\ensuremath{\Sigma}}
% majuscules produites par \MakeUppercase dans les titres (α-aminés -> Α)
\newunicodechar{Α}{\ensuremath{\alpha}}
\newunicodechar{Β}{\ensuremath{\beta}}
\newunicodechar{Γ}{\ensuremath{\Gamma}}
\newunicodechar{Δ}{\ensuremath{\Delta}}
\newunicodechar{Ε}{\ensuremath{\varepsilon}}
\newunicodechar{Θ}{\ensuremath{\Theta}}
\newunicodechar{Λ}{\ensuremath{\Lambda}}
\newunicodechar{Μ}{\ensuremath{\mu}}
\newunicodechar{Τ}{\ensuremath{\tau}}
\newunicodechar{Φ}{\ensuremath{\Phi}}
\newunicodechar{Ψ}{\ensuremath{\Psi}}
\newunicodechar{Ω}{\ensuremath{\Omega}}
\newunicodechar{↦}{\ensuremath{\mapsto}}
\newunicodechar{∈}{\ensuremath{\in}}
\newunicodechar{∉}{\ensuremath{\notin}}
\newunicodechar{∧}{\ensuremath{\wedge}}
\newunicodechar{⇔}{\ensuremath{\Leftrightarrow}}
\newunicodechar{⟺}{\ensuremath{\Longleftrightarrow}}
\newunicodechar{⇒}{\ensuremath{\Rightarrow}}
\newunicodechar{⟹}{\ensuremath{\Longrightarrow}}
\newunicodechar{∘}{\ensuremath{\circ}}
\newunicodechar{∪}{\ensuremath{\cup}}
\newunicodechar{∩}{\ensuremath{\cap}}
\newunicodechar{≡}{\ensuremath{\equiv}}
\newunicodechar{⊂}{\ensuremath{\subset}}
\newunicodechar{⊥}{\ensuremath{\perp}}
\newunicodechar{∃}{\ensuremath{\exists}}
\newunicodechar{∀}{\ensuremath{\forall}}
\newunicodechar{∛}{\ensuremath{\sqrt[3]{\,}}}
\newunicodechar{∜}{\ensuremath{\sqrt[4]{\,}}}
\newunicodechar{⃗}{\ensuremath{{}^{\to}}}
\newunicodechar{ⁿ}{\ifmmode^{n}\else\textsuperscript{\ensuremath{n}}\fi}
\newunicodechar{ˣ}{\ifmmode^{x}\else\textsuperscript{\ensuremath{x}}\fi}
\newunicodechar{ᵇ}{\ifmmode^{b}\else\textsuperscript{\ensuremath{b}}\fi}
\newunicodechar{ʳ}{\ifmmode^{r}\else\textsuperscript{\ensuremath{r}}\fi}
\newunicodechar{ᵘ}{\ifmmode^{u}\else\textsuperscript{\ensuremath{u}}\fi}
\newunicodechar{ʸ}{\ifmmode^{y}\else\textsuperscript{\ensuremath{y}}\fi}
\newunicodechar{ᵃ}{\ifmmode^{a}\else\textsuperscript{\ensuremath{a}}\fi}
\newunicodechar{ᵉ}{\ifmmode^{e}\else\textsuperscript{\ensuremath{e}}\fi}
\newunicodechar{ᵏ}{\ifmmode^{k}\else\textsuperscript{\ensuremath{k}}\fi}
\newunicodechar{ᵖ}{\ifmmode^{p}\else\textsuperscript{\ensuremath{p}}\fi}
\newunicodechar{ᴺ}{\ifmmode^{N}\else\textsuperscript{\ensuremath{N}}\fi}
\newunicodechar{ᶜ}{\ifmmode^{c}\else\textsuperscript{\ensuremath{c}}\fi}
\newunicodechar{ʰ}{\ifmmode^{h}\else\textsuperscript{\ensuremath{h}}\fi}
\newunicodechar{ᵠ}{\ifmmode^{q}\else\textsuperscript{\ensuremath{q}}\fi}
\newunicodechar{ₙ}{\ifmmode_{n}\else\textsubscript{\ensuremath{n}}\fi}
\newunicodechar{ₐ}{\ifmmode_{a}\else\textsubscript{\ensuremath{a}}\fi}
\newunicodechar{ᵢ}{\ifmmode_{i}\else\textsubscript{\ensuremath{i}}\fi}
\newunicodechar{ᵣ}{\ifmmode_{r}\else\textsubscript{\ensuremath{r}}\fi}
\newunicodechar{ₖ}{\ifmmode_{k}\else\textsubscript{\ensuremath{k}}\fi}
\newunicodechar{ᵦ}{\ifmmode_{\beta}\else\textsubscript{\ensuremath{\beta}}\fi}
\newunicodechar{ₓ}{\ifmmode_{x}\else\textsubscript{\ensuremath{x}}\fi}
\newfontfamily\popdisplay{ArchivoBlack-Regular.ttf}[Path=../../../fonts/]
\newfontfamily\poplogo{SpaceGrotesk-Bold.ttf}[Path=../../../fonts/]
\newfontfamily\popsemi{PlusJakartaSans-SemiBold.ttf}[Path=../../../fonts/]
\newfontfamily\popxbold{PlusJakartaSans-ExtraBold.ttf}[Path=../../../fonts/]
\newfontfamily\popxboldls{PlusJakartaSans-ExtraBold.ttf}[Path=../../../fonts/,LetterSpace=3.0]

\setlist[enumerate]{leftmargin=1.7em,itemsep=5pt,topsep=4pt}
\setlist[itemize]{leftmargin=1.7em,itemsep=4pt,topsep=4pt,label={\color{poporange}\textbullet}}
\setlength{\parindent}{0pt}
\setlength{\parskip}{5pt}
\renewcommand{\arraystretch}{1.25}

% pgfplots : style Pop par défaut
\pgfplotsset{
  every axis/.append style={axis line style={popink,line width=1pt},tick style={popink},
    grid style={popline,line width=0.5pt},label style={font=\small},tick label style={font=\footnotesize}},
  popcurve/.style={poporange,line width=1.8pt,no markers,smooth},
}
% Même style disponible dans un \draw TikZ simple (hors pgfplots)
\tikzset{popcurve/.style={poporange,line width=1.8pt}}
\tikzset{popgrid/.style={popline,very thin}}
\colorlet{popcurve}{poporange} % le nom du style est parfois utilisé comme couleur

% ============================================================
% Pill / squiggle
% ============================================================
\newcommand{\poppill}[3]{%
  \tikz[baseline=-0.55ex]\node[draw=popink,line width=1.2pt,rounded corners=2.6pt,%
    fill=#2,inner xsep=6.5pt,inner ysep=3pt]{%
    \popxboldls\fontsize{7}{7}\selectfont\color{#3}\MakeUppercase{#1}};%
}
\newcommand{\popsquiggle}[1]{%
  \tikz[baseline=-0.4ex]{
    \draw[#1,line width=2.6pt,line cap=round]
      plot [smooth] coordinates {(0,0.09) (0.34,0.01) (0.68,0.21) (1.02,0.01) (1.36,0.10)};
  }%
}
% Mot-clé surligné (marqueur orange sous le texte)
\newcommand{\cle}[1]{{\popxbold\color{popdark}#1}}

% ============================================================
% En-tête / pied
% ============================================================
\pagestyle{fancy}
\fancyhf{}
\renewcommand{\headrulewidth}{0pt}
\renewcommand{\footrulewidth}{0pt}
\lhead{\raisebox{-0.15ex}{\poplogo\fontsize{13}{13}\selectfont\color{popink}OnBuch\color{poporange}+}}
\rhead{\poppill{\DOCMATIERE\ \textbullet\ \DOCNIVEAU\ \textbullet\ Chapter \DOCLECON}{white}{popink}}
\lfoot{\popxbold\fontsize{7.6}{7.6}\selectfont\color{popmuted}\MakeUppercase{Signed course OnBuch+}\ \textbullet\ {\popsemi\DOCTITRE}}
\rfoot{\tikz[baseline=-0.6ex]\node[rounded corners=8.5pt,fill=popink,minimum height=17pt,minimum width=17pt,inner xsep=5pt,inner sep=0]{\popxbold\fontsize{8}{8}\selectfont\color{popcream}\thepage\,/\,\pageref*{LastPage}};}

% ============================================================
% Titres
% ============================================================
\newcommand{\chaptitre}[2]{%
  \begin{tcolorbox}[enhanced,colback=poporange,colframe=popink,boxrule=2.6pt,arc=11pt,
      drop shadow=popink,left=15pt,right=15pt,top=12pt,bottom=11pt,before skip=3pt,after skip=18pt]
    \poppill{#2}{popink}{popcream}\par\vspace{9pt}
    {\raggedright\hyphenpenalty=10000\exhyphenpenalty=10000\popdisplay\fontsize{22}{24}\selectfont\color{popcream}\MakeUppercase{#1}\par}
  \end{tcolorbox}
}
% Section numérotée : pastille orange avec le numéro + titre Archivo
\newcounter{popsec}
\newcommand{\coursec}[1]{%
  \stepcounter{popsec}\setcounter{popsub}{0}%
  \par\vspace{16pt}\noindent
  \begin{minipage}{\linewidth}
  \tikz[baseline=-0.7ex]\node[draw=popink,line width=1.6pt,fill=poporange,rounded corners=4pt,minimum size=22pt,inner sep=0]{\popdisplay\fontsize{12}{12}\selectfont\color{popcream}\thepopsec};%
  \hspace{8pt}{\popdisplay\fontsize{15}{17}\selectfont\color{popink}\MakeUppercase{#1}}\par
  \vspace{4pt}\noindent{\color{poporange}\rule{36pt}{3.6pt}}
  \end{minipage}\par\nopagebreak\vspace{8pt}
}
\newcounter{popsub}
\newcommand{\courssub}[1]{%
  \stepcounter{popsub}%
  \par\vspace{9pt}\noindent{\popxbold\fontsize{11.5}{13}\selectfont\color{popink}{\color{poporange}\thepopsec.\thepopsub}\hspace{6pt}#1}\par\nopagebreak\vspace{4pt}
}

% ============================================================
% Boîtes pédagogiques — même recette : fond blanc, bordure épaisse colorée,
% ombre dure encre, badge plein. Titre optionnel [#1].
% ============================================================
\tcbset{popbase/.style={breakable,enhanced,colback=white,boxrule=2.3pt,arc=9pt,
  shadow={3pt}{-3pt}{0pt}{popink},left=14pt,right=14pt,top=11pt,bottom=11pt,before skip=11pt,after skip=15pt,
  fontupper=\popsemi\fontsize{10.6}{14.5}\selectfont\color{popink},
  fontlower=\popsemi\fontsize{10.6}{14.5}\selectfont\color{popink}}}
% Coche dessinée (le glyphe ✓ n'existe pas dans les polices du document)
\newcommand{\popcheck}[1]{\tikz[baseline=-0.5ex]\draw[#1,line width=1.6pt,line cap=round,line join=round](-0.13,0.0)--(-0.04,-0.09)--(0.13,0.1);}
\providecommand{\coche}{\popcheck{popgreen}}
\providecommand{\resultat}[1]{\par\smallskip\noindent\fcolorbox{popgreen}{popgreenL}{\parbox{\dimexpr\linewidth-2\fboxsep-2\fboxrule\relax}{\textbf{Result.} #1}}\par\smallskip}
\newcommand{\popboxhead}[3]{\poppill{#1}{#2}{white}\ifx\relax#3\relax\else\hspace{6pt}{\popxbold\fontsize{10.6}{12}\selectfont\color{popink}#3}\fi\par\vspace{7pt}}

\newtcolorbox{aretenir}[1][]{popbase,colframe=popgreen,before upper={\popboxhead{Key points}{popgreen}{#1}}}
\newtcolorbox{exemplebox}[1][]{popbase,colframe=poporange,before upper={\popboxhead{Exemple}{poporange}{#1}}}
\newtcolorbox{definition}[1][]{popbase,colframe=popblue,before upper={\popboxhead{Definition}{popblue}{#1}}}
\newtcolorbox{propriete}[1][]{popbase,colframe=poppurple,before upper={\popboxhead{Property}{poppurple}{#1}}}
\newtcolorbox{methode}[1][]{popbase,colframe=popgold,before upper={\popboxhead{Method}{popgold}{#1}}}
\newtcolorbox{attention}[1][]{popbase,colframe=poppink,before upper={\popboxhead{Warning}{poppink}{#1}}}
\newtcolorbox{experience}[1][]{popbase,colframe=popblue,colback=popblueL,before upper={\popboxhead{Experiment}{popblue}{#1}}}
% « Le savais-tu ? » : carte violette pleine (contexte, culture, Cameroun)
\newtcolorbox{savaistu}[1][]{popbase,colback=poppurple,colframe=popink,
  fontupper=\popsemi\fontsize{10.4}{14.2}\selectfont\color{white},
  before upper={\poppill{Did you know?}{white}{poppurple}\ifx\relax#1\relax\else\hspace{6pt}{\popxbold\color{white}#1}\fi\par\vspace{7pt}}}
% Exercice résolu : énoncé en haut, \tcblower puis la solution sur fond crème
\newcounter{popexo}\newcounter{popexr}
\newtcolorbox{exoresolu}[1][]{popbase,colframe=poporange,colbacklower=popcream,
  segmentation style={popink,line width=1.2pt,dashed},
  before upper={\stepcounter{popexr}\popboxhead{Worked example \thepopexr}{poporange}{#1}},
  before lower={\poppill{Solution}{popink}{popcream}\par\vspace{6pt}}}
% Exercice (non résolu en place) — la correction est dans la partie Corrigés
\newtcolorbox{exercice}[1][]{popbase,colframe=popink,
  before upper={\stepcounter{popexo}\popboxhead{Exercise \thepopexo}{popink}{#1}}}
% Corrigé
\newtcolorbox{corrige}[1][]{popbase,colframe=popgreen,colback=popgreenL,
  before upper={\popboxhead{Solution}{popgreen}{#1}}}
% Objectifs / prérequis / bilan
\newtcolorbox{objectifs}{popbase,colframe=popink,colback=poporangeL,
  before upper={\popboxhead{Chapter objectives}{popink}{}}}
\newtcolorbox{prerequis}{popbase,colframe=popink,colback=popblueL,
  before upper={\popboxhead{Prerequisites}{popblue}{}}}
\newtcolorbox{bilan}{popbase,colframe=popink,colback=white,boxrule=2.8pt,
  before upper={\popboxhead{Fiche bilan}{poporange}{L'essentiel à connaître pour l'examen}}}
% Figure encadrée avec légende
\newtcolorbox{popfigure}[1][]{enhanced,breakable=false,colback=white,colframe=popline,boxrule=1.4pt,arc=8pt,
  left=10pt,right=10pt,top=10pt,bottom=8pt,before skip=10pt,after skip=12pt,halign=center,
  lower separated=false,halign lower=center,
  fontlower=\popsemi\fontsize{9}{11.5}\selectfont\color{popmuted},
  #1}
\newcommand{\legende}[1]{\par\vspace{6pt}{\popsemi\fontsize{9}{11.5}\selectfont\color{popmuted}#1}}

% ============================================================
% Signature OnBuch+
% ============================================================
\newcommand{\poplogoinline}[1]{{\poplogo\fontsize{#1}{#1}\selectfont\color{popink}OnBuch\color{poporange}+}}
\newcommand{\signatureonbuch}{%
  \begin{tcolorbox}[enhanced,colback=popink,colframe=popink,boxrule=2.2pt,arc=10pt,
      drop shadow=poporange,left=14pt,right=14pt,top=10pt,bottom=10pt,before skip=0pt,after skip=0pt]
    \tikz[baseline=-0.7ex]\node[circle,fill=popgreen,draw=popcream,line width=1.3pt,minimum size=22pt,inner sep=0]{\popcheck{white}};%
    \hspace{9pt}\begin{minipage}[c]{0.62\linewidth}
      {\popxbold\fontsize{12}{13}\selectfont\color{popcream}Signed\ \textbullet\ {\poplogo OnBuch\color{poporange}+}}\par\vspace{2pt}
      {\raggedright\popsemi\fontsize{8}{10.5}\selectfont\color{popline}\MakeUppercase{OnBuch+ teaching team}\\\MakeUppercase{Follows the official MINESEC syllabus}\par}
    \end{minipage}\hfill
    {\poplogo\fontsize{22}{22}\selectfont\color{popcream}OnBuch\color{poporange}+}
  \end{tcolorbox}%
}

% ============================================================
% Page de garde
% #1 titre · #2 matière · #3 niveau · #4 module · #5 n° leçon · #6 durée ·
% #7 description / objectifs (texte ou itemize)
% ============================================================
\newcommand{\pagegarde}[7]{%
  \thispagestyle{empty}
  \newgeometry{a4paper,margin=1.7cm,top=1.6cm,bottom=1.4cm}%
  \begin{tikzpicture}[remember picture,overlay]
    \fill[poporange!18] ([xshift=-3.2cm,yshift=-3.6cm]current page.north east) circle (4.6cm);
    \fill[poppurple!14] ([xshift=2.4cm,yshift=7.6cm]current page.south west) circle (3.8cm);
  \end{tikzpicture}%
  {\poplogo\fontsize{30}{30}\selectfont\color{popink}OnBuch\color{poporange}+}\hfill
  \raisebox{0.6ex}{\poppill{Signed course \textbullet\ Premium}{popink}{popcream}}\par
  \vspace{1pt}\popsquiggle{poppurple}\par
  \vspace{8pt}
  \begin{tcolorbox}[enhanced,colback=poporange,colframe=popink,boxrule=2.8pt,arc=13pt,
      drop shadow=popink,left=18pt,right=18pt,top=12pt,bottom=13pt,after skip=10pt]
    \poppill{#2 \textbullet\ #3}{popink}{popcream}\hspace{5pt}\poppill{#4}{white}{popink}\par\vspace{8pt}
    {\popdisplay\fontsize{14}{14}\selectfont\color{popink}CHAPTER #5}\par\vspace{4pt}
    {\raggedright\hyphenpenalty=10000\exhyphenpenalty=10000\popdisplay\fontsize{25}{27}\selectfont\color{popcream}\MakeUppercase{#1}\par}
    \vspace{8pt}
    {\popxbold\fontsize{9.5}{9.5}\selectfont\color{popink}\MakeUppercase{Suggested time: #6}}
  \end{tcolorbox}
  \begin{tcolorbox}[enhanced,colback=white,colframe=popink,boxrule=2.2pt,arc=10pt,
      drop shadow=popink,left=16pt,right=16pt,top=10pt,bottom=10pt,after skip=8pt]
    \poppill{In this chapter}{poppurple}{white}\par\vspace{5pt}
    \popsemi\fontsize{9.3}{12.6}\selectfont\color{popink}#7
  \end{tcolorbox}
  \noindent\begin{minipage}[t]{0.487\linewidth}
    \begin{tcolorbox}[enhanced,colback=popgreen,colframe=popink,boxrule=2.2pt,arc=10pt,
        drop shadow=popink,left=11pt,right=11pt,top=10pt,bottom=10pt,height=3.1cm,valign=center]
      \tcbox[colback=white,colframe=popink,boxrule=1.3pt,arc=5pt,boxsep=0pt,nobeforeafter,
        left=3pt,right=3pt,top=3pt,bottom=3pt,valign=center]{\qrcode[height=1.9cm]{https://play.google.com/store/apps/details?id=com.luvvix.onbuch&hl=en}}%
      \hspace{8pt}\begin{minipage}[c]{0.5\linewidth}
        {\popdisplay\fontsize{11}{12.5}\selectfont\color{white}\MakeUppercase{Download OnBuch}}\par\vspace{4pt}
        {\raggedright\popsemi\fontsize{8.4}{11}\selectfont\color{white}Free \textbullet\ Google Play\\AI tutor, past papers, results\par}
      \end{minipage}
    \end{tcolorbox}
  \end{minipage}\hfill
  \begin{minipage}[t]{0.487\linewidth}
    \begin{tcolorbox}[enhanced,colback=poppurple,colframe=popink,boxrule=2.2pt,arc=10pt,
        drop shadow=popink,left=11pt,right=11pt,top=10pt,bottom=10pt,height=3.1cm,valign=center]
      \tikz[baseline=-0.7ex]\node[circle,fill=white,draw=popink,line width=1.3pt,minimum size=1.6cm,inner sep=0]{\popdisplay\fontsize{24}{24}\selectfont\color{poppurple}+};%
      \hspace{8pt}\begin{minipage}[c]{0.6\linewidth}
        {\popdisplay\fontsize{11}{12.5}\selectfont\color{white}\MakeUppercase{Go OnBuch+}}\par\vspace{4pt}
        {\raggedright\popsemi\fontsize{8.4}{11}\selectfont\color{white}All signed courses, unlimited AI tutor, PDF sheets — OnBuch+ tab in the app\par}
      \end{minipage}
    \end{tcolorbox}
  \end{minipage}
  \vspace{8pt}
  \popcontenu
  \vfill
  \signatureonbuch
  \restoregeometry
  \newpage
}

% Rangée « Ce cours contient » (page de garde)
\newcommand{\poptile}[3]{% #1 couleur #2 gros texte #3 légende
  \begin{tcolorbox}[enhanced,colback=#1,colframe=popink,boxrule=2pt,arc=9pt,shadow={2.5pt}{-2.5pt}{0pt}{popink},
      left=6pt,right=6pt,top=9pt,bottom=9pt,halign=center,valign=center,height=1.75cm,width=\linewidth,nobeforeafter]
    {\popdisplay\fontsize{12.5}{13}\selectfont\color{white}\MakeUppercase{#2}}\par\vspace{3pt}
    {\centering\popsemi\fontsize{7.8}{9.5}\selectfont\color{white}#3\par}
  \end{tcolorbox}}
\newcommand{\popcontenu}{%
  \par\noindent{\popxboldls\fontsize{7.5}{7.5}\selectfont\color{popmuted}THIS SIGNED COURSE CONTAINS}\par\vspace{7pt}
  \noindent\begin{minipage}[t]{0.235\linewidth}\poptile{popblue}{Course}{complete, illustrated, step by step}\end{minipage}\hfill
  \begin{minipage}[t]{0.235\linewidth}\poptile{poporange}{Methods}{and worked examples}\end{minipage}\hfill
  \begin{minipage}[t]{0.235\linewidth}\poptile{poppink}{Exercises}{exam style + solutions}\end{minipage}\hfill
  \begin{minipage}[t]{0.235\linewidth}\poptile{popgreen}{Summary sheet}{the essentials to remember}\end{minipage}\par}

% Sommaire
\newtcolorbox{sommaire}{enhanced,colback=white,colframe=popink,boxrule=2.2pt,arc=10pt,
  drop shadow=popink,left=18pt,right=18pt,top=13pt,bottom=13pt,before skip=6pt,after skip=20pt,
  before upper={\poppill{Contents}{poppurple}{white}\par\vspace{9pt}\popsemi\fontsize{10.6}{14.5}\selectfont\color{popink}}}

% Signature de fin de cours — poussée tout en bas de la dernière page
\newcommand{\findecours}{%
  \par\vfill\nopagebreak
  \begin{center}\popsquiggle{poporange}\end{center}\vspace{4pt}
  \signatureonbuch
}
\AtBeginDocument{\def\llbracket{\mathopen{[\mkern-3.5mu[}}\def\rrbracket{\mathclose{]\mkern-3.5mu]}}} % intervalles d'entiers (glyphe ⟦ absent de la police)
% Glyphes absents de Fira Math : redéfinis à partir de glyphes disponibles
\AtBeginDocument{%
  \def\setminus{\mathbin{\backslash}}\let\smallsetminus\setminus
  \def\vdots{\mathord{\vbox{\baselineskip3.2pt\lineskiplimit0pt\kern4pt\hbox{.}\hbox{.}\hbox{.}}}}%
  \def\ddots{\mathinner{\mkern1mu\raise6.5pt\hbox{.}\mkern2mu\raise3.6pt\hbox{.}\mkern2mu\raise0.7pt\hbox{.}\mkern1mu}}%
  \def\triangle{\mathord{\scalebox{0.9}{$\symup{\Delta}$}}}%
  \def\nabla{\mathord{\raisebox{\depth}{\scalebox{1}[-1]{$\symup{\Delta}$}}}}%
  \def\checkmark{\popcheck{popdark}}%
  \def\star{\mathbin{\ast}}\def\bigstar{\mathord{\ast}}%
  \def\diamond{\mathbin{\scalebox{0.6}{$\Diamond$}}}%
  \def\bigcup{\mathop{\mathchoice{\raisebox{-0.3em}{\scalebox{1.7}{$\cup$}}}{\scalebox{1.25}{$\cup$}}{\cup}{\cup}}\displaylimits}%
  \def\bigcap{\mathop{\mathchoice{\raisebox{-0.3em}{\scalebox{1.7}{$\cap$}}}{\scalebox{1.25}{$\cap$}}{\cap}{\cap}}\displaylimits}%
  \def\lesseqqgtr{\mathrel{\lessgtr}}\def\bigoplus{\mathop{\oplus}\displaylimits}%
}
\providecommand{\ssi}{if and only if} % abréviation fréquente
\AtBeginDocument{\providecommand{\wideparen}{\overparen}} % arc de cercle

%% FIGFIX-BEGIN v5 — correctifs globaux des figures (docs/onbuch-generation/tools/figfix)
\usepackage{adjustbox}\usepackage{etoolbox}\tcbuselibrary{raster}
% toute figure TikZ/pgfplots plus large que la ligne est réduite à la largeur disponible
\BeforeBeginEnvironment{tikzpicture}{\begin{adjustbox}{max width=\linewidth}}
\AfterEndEnvironment{tikzpicture}{\end{adjustbox}}
% légendes pgfplots lisibles par-dessus les courbes
\pgfplotsset{every axis/.append style={legend style={fill=white,fill opacity=0.92,text opacity=1,draw=black!15,inner sep=3pt}}}
% étiquettes d'arêtes (syntaxe edge["..."]) avec fond blanc
\tikzset{every edge quotes/.append style={fill=white,inner sep=1.2pt}}
%% FIGFIX-END
\begin{document}

\pagegarde{Motion in a circle with uniform speed}{Pure Mathematics with Mechanics}{Upper Sixth}{Upper Sixth — Pure mathematics with mechanics}{12}{13 periods}{This chapter studies the kinematics and dynamics of uniform circular motion: angular quantities, centripetal acceleration and force, and the link between linear and angular speed. It then applies these ideas to the conical pendulum, deriving its period formula and solving GCE-style problems.
\vspace{4pt}
\begin{multicols}{2}\small
\begin{itemize}
\item By the end of this chapter I can define angular displacement, angular speed and period, and use their units correctly (rad, rad/s, s, Hz).
\item By the end of this chapter I can explain why a body in uniform circular motion is accelerating even though its speed is constant.
\item By the end of this chapter I can calculate centripetal acceleration and centripetal force using a = v²/r = rω² and F = mv²/r = mrω².
\item By the end of this chapter I can convert fluently between linear speed and angular speed using v = rω in both directions.
\item By the end of this chapter I can solve applied problems on vehicles in curves, rotating wheels, satellites and fairground rides.
\item By the end of this chapter I can describe a conical pendulum and identify all forces acting on the bob.
\end{itemize}\end{multicols}}

\begin{sommaire}
\begin{multicols}{2}
\begin{enumerate}
\item Circular Motion and Angular Quantities
\item Centripetal Acceleration and Centripetal Force
\item Linear Speed, Angular Speed and Their Relationship
\item Applications of Linear and Angular Speed
\item The Conical Pendulum: Concept and Forces
\item Period of the Conical Pendulum: Derivation and Applications
\item Angular and Linear Speed in the Conical Pendulum; Synthesis and Exam Practice
\item Integration activity
\item Methods and worked examples
\item Exercises
\item Solutions to the exercises
\item Summary sheet
\end{enumerate}
\end{multicols}
\end{sommaire}

\begin{objectifs}
\begin{itemize}
\item By the end of this chapter I can define angular displacement, angular speed and period, and use their units correctly (rad, rad/s, s, Hz).
\item By the end of this chapter I can explain why a body in uniform circular motion is accelerating even though its speed is constant.
\item By the end of this chapter I can calculate centripetal acceleration and centripetal force using a = v²/r = rω² and F = mv²/r = mrω².
\item By the end of this chapter I can convert fluently between linear speed and angular speed using v = rω in both directions.
\item By the end of this chapter I can solve applied problems on vehicles in curves, rotating wheels, satellites and fairground rides.
\item By the end of this chapter I can describe a conical pendulum and identify all forces acting on the bob.
\item By the end of this chapter I can derive the period formula T = 2π√(h/g) for a conical pendulum from Newton's second law.
\item By the end of this chapter I can relate angular and linear speed for a conical pendulum and use tan θ = v²/(rg) = rω²/g.
\item By the end of this chapter I can solve GCE Advanced Level problems on circular motion and conical pendulums in a clear, structured way.
\end{itemize}
\end{objectifs}

\begin{prerequis}
\begin{itemize}
\item Newton's laws of motion and resolution of forces into components (Lower Sixth Mechanics).
\item Kinematics: speed, velocity, acceleration as rates of change (Lower Sixth).
\item Radian measure and the relation arc length s = rθ (Pure Mathematics, Lower Sixth).
\item Trigonometric ratios in right-angled triangles (sin, cos, tan) and Pythagoras' theorem.
\item Algebraic manipulation: rearranging formulae, solving simultaneous equations.
\end{itemize}
\end{prerequis}

\newpage

\coursec{Circular Motion and Angular Quantities}

At the busy Carrefour Warda roundabout in Yaoundé, a moto-taxi rider leans into the curve while a passenger clings on, feeling pushed outward. Why does the rider not fall, and why does the load on the carrier seem to fly off if the turn is taken too fast? On the playground of Lycée de Nkolbisson, a child swings a stone tied to a string in a horizontal circle above her head: the string makes a constant angle with the vertical, and the stone spins faster when the string is shortened. These everyday Cameroonian scenes all obey the same mathematics: \cle{motion in a circle with uniform speed}. This chapter gives you the tools to explain, predict and calculate them.

Before we can talk about forces (why the rider leans, why the string tilts), we must first learn the \emph{language} of circular motion: how to describe position, rotation and rate of rotation. That is the work of this first section.

\courssub{What is circular motion? Everyday observations}

Look around you and circular motion is everywhere:

\begin{itemize}
\item the wheels of a \emph{bendskin} (bicycle taxi) rolling through Bamenda;
\item the blades of a ceiling fan turning in a classroom in Buea;
\item a merry-go-round at a Christmas fairground;
\item the hands of the clock on the wall of your examination hall;
\item the Earth travelling around the Sun, and the Moon around the Earth;
\item a stone whirled on the end of a string.
\end{itemize}

In each case, some point of the object moves along a \cle{circle} (or very nearly one). In this chapter we idealise the moving object as a \emph{particle} $P$ travelling on a circle of fixed centre $O$ and fixed radius $r$.

\begin{definition}[Circular motion and uniform circular motion]
A particle is in \cle{circular motion} when its path is a circle of centre $O$ and radius $r$. The motion is called \cle{uniform circular motion} when the \textbf{speed} of the particle along the circle is \emph{constant}.
\end{definition}

\begin{attention}[Constant speed does not mean constant velocity]
Speed is a scalar; velocity is a vector. Even when the speed $v$ is constant, the \emph{direction} of the velocity changes continuously, because the velocity vector is always tangent to the circle and the tangent direction turns as the particle moves. Since velocity changes, the particle \textbf{is accelerating} — this surprising fact is the key to the whole chapter and will be made precise in the next section (centripetal acceleration). For now, keep the slogan: \emph{uniform circular motion = constant speed, changing velocity.}
\end{attention}

\courssub{Angular displacement and the radian}

To locate a particle on a circle, distances along a straight axis are clumsy. The natural coordinate is an \textbf{angle}.

\begin{definition}[Angular displacement]
Let a particle move on a circle of centre $O$ from position $A$ to position $B$. The \cle{angular displacement} is the angle $\theta = \angle AOB$ swept out by the radius $OP$. It is measured in \cle{radians} (symbol rad).
\end{definition}

Why invent a new unit, the radian, when degrees exist? Because the radian makes the geometry of the circle collapse into one beautiful formula. Measure the arc length $s$ from $A$ to $B$ along the circle. If you double the radius while keeping the same angle, the arc doubles: arc length is proportional to radius. The radian is \emph{defined} so that the constant of proportionality is exactly $1$.

\begin{definition}[The radian]
One \cle{radian} is the angle subtended at the centre of a circle by an arc whose length equals the radius. Consequently, for any angle $\theta$ in radians,
\[
\boxed{\,s = r\theta\,}
\]
where $s$ is the arc length and $r$ the radius. A full revolution corresponds to $\theta = 2\pi$ rad, since the full circumference is $s = 2\pi r$. The radian is a dimensionless unit (it is a ratio of two lengths), but we keep the symbol ``rad'' for clarity.
\end{definition}

\begin{popfigure}
\begin{center}
\begin{tikzpicture}[scale=2.2,>=stealth]
\draw[popline] (0,0) circle (1.5);
\fill (0,0) circle (0.02) node[below left=2pt] {$O$};
\coordinate (A) at (1.5,0);
\coordinate (B) at (60:1.5);
\draw[->,thick,popblue] (0,0) -- (A) node[midway,below] {$r$};
\draw[->,thick,popblue] (0,0) -- (B) node[midway,above left=1pt] {$r$};
\fill[poporange] (A) circle (0.035) node[right=2pt] {$A$};
\fill[poporange] (B) circle (0.035) node[above right=1pt] {$B$};
\draw[->,thick,popgreen] (0.55,0) arc (0:60:0.55) node[midway,right=4pt] {$\theta$};
\draw[->,very thick,poppurple] ([shift={(0.06,0)}]1.5,0) arc (0:60:1.5) node[midway,right=6pt] {$s=r\theta$};
\draw[->,dashed,popmuted] (B) arc (60:340:1.5);
\end{tikzpicture}
\end{center}
\legende{A particle moves from $A$ to $B$ on a circle of radius $r$. The angular displacement is $\theta$ (in radians) and the arc travelled is $s = r\theta$.}
\end{popfigure}

\begin{propriete}[Conversion between degrees and radians]
Since $360^{\circ}$ corresponds to $2\pi$ rad,
\[
180^{\circ} = \pi\ \text{rad}, \qquad 1^{\circ} = \frac{\pi}{180}\ \text{rad}, \qquad 1\ \text{rad} = \frac{180}{\pi}\ \text{degrees} \approx 57.3^{\circ}.
\]
To convert degrees to radians, multiply by $\dfrac{\pi}{180}$; to convert radians to degrees, multiply by $\dfrac{180}{\pi}$.
\end{propriete}

\begin{exemplebox}[Converting angles]
$60^{\circ} = 60\times\dfrac{\pi}{180} = \dfrac{\pi}{3}\ \text{rad}$; \quad $135^{\circ} = \dfrac{3\pi}{4}\ \text{rad}$; \quad $\dfrac{5\pi}{6}\ \text{rad} = \dfrac{5\pi}{6}\times\dfrac{180}{\pi} = 150^{\circ}$.
\end{exemplebox}

\courssub{Angular speed}

A degree or a radian tells us \emph{how far} the radius has turned. We now ask \emph{how fast} it turns.

\begin{definition}[Angular speed]
If the radius $OP$ sweeps an angular displacement $\Delta\theta$ in a time interval $\Delta t$, the \cle{angular speed} is
\[
\boxed{\;\omega = \frac{\Delta\theta}{\Delta t}\;}
\]
with $\Delta\theta$ in radians and $\Delta t$ in seconds. The SI unit of $\omega$ is the \textbf{radian per second}, written $\mathrm{rad\,s^{-1}}$ (or rad/s). For uniform circular motion, $\omega$ is constant: equal angles are swept in equal times. (One can also define an \emph{angular velocity} vector $\vec{\omega}$ pointing along the axis of rotation; at GCE Advanced Level we only need its magnitude $\omega$.)
\end{definition}

\courssub{Period and frequency}

Two more everyday words describe rotation, and examiners love converting between them.

\begin{definition}[Period and frequency]
The \cle{period} $T$ is the time for \textbf{one complete revolution}; its SI unit is the second (s). The \cle{frequency} $f$ is the \textbf{number of complete revolutions per second}; its SI unit is the hertz (Hz), with $1\ \mathrm{Hz} = 1\ \mathrm{s^{-1}}$. Clearly
\[
\boxed{\;f = \frac{1}{T} \qquad\text{and}\qquad T = \frac{1}{f}.\;}
\]
\end{definition}

Now connect these to $\omega$. In one period $T$, the radius sweeps exactly one full turn, i.e.\ $\Delta\theta = 2\pi$ rad. Hence:

\begin{propriete}[Fundamental relations for uniform circular motion]
\[
\boxed{\;\omega = \frac{2\pi}{T} = 2\pi f\;}
\]
This follows directly from $\omega = \Delta\theta/\Delta t$ with one revolution: $\omega = \dfrac{2\pi}{T}$, and substituting $T = 1/f$ gives $\omega = 2\pi f$. \emph{(The derivation is examinable and takes two lines — learn it.)}
\end{propriete}

\begin{exoresolu}[The second hand of a clock]
The second hand of the clock in your examination hall makes one revolution per minute. Find (a) its period, (b) its frequency, (c) its angular speed.
\tcblower
\textbf{Solution.}
\begin{enumerate}
\item One revolution takes $60\ \mathrm{s}$, so $T = 60\ \mathrm{s}$.
\item $f = \dfrac{1}{T} = \dfrac{1}{60}\ \mathrm{Hz} \approx 0.0167\ \mathrm{Hz}$.
\item $\omega = \dfrac{2\pi}{T} = \dfrac{2\pi}{60} = \dfrac{\pi}{30}\ \mathrm{rad\,s^{-1}} \approx 0.105\ \mathrm{rad\,s^{-1}}$.
\end{enumerate}
\emph{Check:} in $15\ \mathrm{s}$ the hand sweeps $\theta = \omega t = \frac{\pi}{30}\times 15 = \frac{\pi}{2}$ rad $= 90^{\circ}$ — exactly a quarter turn, as expected.
\end{exoresolu}

\begin{popfigure}
\begin{center}
\begin{tikzpicture}[scale=1.6,>=stealth]
\draw[thick,popline] (0,0) circle (1.6);
\foreach \a in {0,30,...,330} \draw[popmuted] (\a:1.45) -- (\a:1.6);
\fill (0,0) circle (0.03);
\draw[->,very thick,popblue] (0,0) -- (90:1.35) node[pos=0.75,left=1pt] {};
\draw[->,dashed,popmuted] (0,0) -- (0:1.35);
\draw[->,thick,poporange] (0.7,0) arc (0:90:0.7) node[midway,below right=1pt] {$\theta=\omega t$};
\node[popdark] at (0,-2.0) {$T = 60\ \mathrm{s}$,\quad $\omega = \dfrac{2\pi}{T} = \dfrac{\pi}{30}\ \mathrm{rad\,s^{-1}}$};
\end{tikzpicture}
\end{center}
\legende{The second hand of a clock: period $T = 60$ s, angular speed $\omega = 2\pi/T$. In time $t$ it sweeps $\theta = \omega t$.}
\end{popfigure}

\courssub{Revolutions per minute (rpm)}

Machines rarely quote $\omega$ in $\mathrm{rad\,s^{-1}}$. Grinding wheels, fans, engines and DJ turntables are rated in \cle{revolutions per minute} (rpm). If a wheel makes $N$ revolutions in one minute, then each revolution is $2\pi$ rad and one minute is $60$ s, so
\[
\omega = \frac{N \times 2\pi\ \text{rad}}{60\ \text{s}} = \frac{2\pi N}{60}\ \mathrm{rad\,s^{-1}}.
\]

\begin{methode}[Converting rpm to rad/s]
To convert a rotation rate of $N$ rpm to an angular speed in $\mathrm{rad\,s^{-1}}$:
\begin{enumerate}
\item Multiply $N$ by $2\pi$ to get the angle swept per minute, in radians.
\item Divide by $60$ to convert ``per minute'' to ``per second''.
\item Write the result with the unit $\mathrm{rad\,s^{-1}}$, keeping $\pi$ symbolic unless a decimal is requested.
\end{enumerate}
In short: $\omega = \dfrac{2\pi N}{60} = \dfrac{\pi N}{30}\ \mathrm{rad\,s^{-1}}$.
\end{methode}

\begin{exoresolu}[The grinding wheel in Douala]
A grinding wheel in a workshop in Douala spins at $1200$ rpm. Find (a) its angular speed in $\mathrm{rad\,s^{-1}}$, (b) its frequency in Hz, (c) its period.
\tcblower
\textbf{Solution.}
\begin{enumerate}
\item $\omega = \dfrac{2\pi N}{60} = \dfrac{2\pi \times 1200}{60} = 40\pi\ \mathrm{rad\,s^{-1}} \approx 125.7\ \mathrm{rad\,s^{-1}}$.
\item $f = \dfrac{\omega}{2\pi} = \dfrac{40\pi}{2\pi} = 20\ \mathrm{Hz}$ (equivalently, $1200/60 = 20$ revolutions per second).
\item $T = \dfrac{1}{f} = \dfrac{1}{20} = 0.05\ \mathrm{s}$.
\end{enumerate}
So the wheel completes a full turn every fiftieth of a second — no wonder a slipping workpiece is dangerous!
\end{exoresolu}

\begin{attention}[Never mix degrees and radians]
The formulas $s = r\theta$, $\omega = \Delta\theta/\Delta t$ and $\omega = 2\pi/T$ are valid \textbf{only when angles are in radians}. If a problem gives an angle in degrees, convert first: multiply by $\pi/180$. A candidate who writes ``$\omega = 360/60 = 6\ \mathrm{rad\,s^{-1}}$'' for the second hand of a clock has mixed degrees with radians and loses the marks — the correct value is $\pi/30\ \mathrm{rad\,s^{-1}}$.
\end{attention}

\begin{exemplebox}[A fairground merry-go-round]
A merry-go-round completes one revolution every $8.0\ \mathrm{s}$. Its angular speed is
\[
\omega = \frac{2\pi}{T} = \frac{2\pi}{8.0} = \frac{\pi}{4}\ \mathrm{rad\,s^{-1}} \approx 0.785\ \mathrm{rad\,s^{-1}},
\]
and its frequency is $f = 1/8.0 = 0.125\ \mathrm{Hz}$. In $20\ \mathrm{s}$ a child on it sweeps $\theta = \omega t = \frac{\pi}{4}\times 20 = 5\pi$ rad, i.e.\ two and a half turns.
\end{exemplebox}

\begin{aretenir}[Key points of this section]
\begin{itemize}
\item Uniform circular motion: constant \emph{speed}, continuously changing \emph{velocity} (direction).
\item Angular displacement $\theta$ in \textbf{radians}; $1$ rad is the angle for which arc $=$ radius; $s = r\theta$; $360^{\circ} = 2\pi$ rad.
\item Angular speed $\omega = \dfrac{\Delta\theta}{\Delta t}$, unit $\mathrm{rad\,s^{-1}}$.
\item Period $T$ (s per revolution), frequency $f = \dfrac{1}{T}$ (Hz).
\item $\omega = \dfrac{2\pi}{T} = 2\pi f$; from $N$ rpm: $\omega = \dfrac{2\pi N}{60}$.
\end{itemize}
\end{aretenir}

\begin{methode}[Exam tip: units earn marks]
GCE Advanced Level examiners expect the unit written at \emph{every} final answer: $\mathrm{rad\,s^{-1}}$ for $\omega$, Hz for $f$, s for $T$, rad for $\theta$. A correct numerical answer with a missing or wrong unit routinely loses the final mark. Make it a habit now: no unit, no full credit.
\end{methode}

\begin{exercice}[Consolidation]
\begin{enumerate}
\item Convert to radians: $45^{\circ}$, $210^{\circ}$, $300^{\circ}$. Convert to degrees: $\dfrac{\pi}{6}$ rad, $\dfrac{7\pi}{4}$ rad.
\item The minute hand of a clock: find its period, frequency and angular speed.
\item A ceiling fan rotates at $300$ rpm. Find $\omega$ in $\mathrm{rad\,s^{-1}}$ and the angle swept in $0.5\ \mathrm{s}$.
\item A wheel has angular speed $10\pi\ \mathrm{rad\,s^{-1}}$. How many revolutions does it make in one minute?
\end{enumerate}
\end{exercice}

\begin{corrige}[Exercise 1]
\begin{enumerate}
\item $45^{\circ} = \frac{\pi}{4}$ rad; $210^{\circ} = \frac{7\pi}{6}$ rad; $300^{\circ} = \frac{5\pi}{3}$ rad. $\frac{\pi}{6}\ \text{rad} = 30^{\circ}$; $\frac{7\pi}{4}\ \text{rad} = 315^{\circ}$.
\item $T = 3600\ \mathrm{s}$; $f = \frac{1}{3600}\ \mathrm{Hz} \approx 2.78\times 10^{-4}\ \mathrm{Hz}$; $\omega = \frac{2\pi}{3600} = \frac{\pi}{1800}\ \mathrm{rad\,s^{-1}} \approx 1.75\times 10^{-3}\ \mathrm{rad\,s^{-1}}$.
\item $\omega = \frac{2\pi\times 300}{60} = 10\pi\ \mathrm{rad\,s^{-1}} \approx 31.4\ \mathrm{rad\,s^{-1}}$; $\theta = \omega t = 10\pi \times 0.5 = 5\pi\ \text{rad}$ (two and a half turns).
\item $f = \frac{\omega}{2\pi} = 5\ \mathrm{Hz}$, so in $60\ \mathrm{s}$ the wheel makes $5\times 60 = 300$ revolutions (i.e.\ $300$ rpm).
\end{enumerate}
\end{corrige}

\begin{savaistu}[Why radians rule the sciences]
The radian is not an arbitrary convention: it is the only angle unit for which $\dfrac{d}{d\theta}(\sin\theta) = \cos\theta$ holds without extra factors, and the only one for which $s = r\theta$ is exactly true. That is why every formula in this chapter — and later in simple harmonic motion, waves and electromagnetism — silently assumes radians. Your calculator should live in RAD mode during mechanics papers!
\end{savaistu}

We can now describe \emph{how fast} anything rotates, from a clock hand to a grinding wheel. But description is not explanation: the moto-taxi at Carrefour Warda leans because something must \emph{pull} him towards the centre of the curve. In the next section we quantify that something: the centripetal acceleration $a = \omega^{2}r$ and the centripetal force that produces it.

\coursec{Centripetal Acceleration and Centripetal Force}

In the previous section we described uniform circular motion with the angular quantities $\theta$, $\omega$ and the period $T$, and we saw that the linear speed $v = r\omega$ stays \emph{constant}. It is tempting to conclude that ``constant speed'' means ``no acceleration''. That conclusion is wrong, and understanding \emph{why} it is wrong is the heart of this section. We shall show that a body moving in a circle is permanently accelerated towards the centre, derive the famous formula $a = \dfrac{v^{2}}{r}$, and then apply Newton's second law to obtain the \cle{centripetal force} $F = \dfrac{mv^{2}}{r}$.

\courssub{Uniform circular motion is accelerated motion}

\begin{experience}[Swinging a stone on a string]
Tie a small stone to a string about \SI{1}{m} long and whirl it in a horizontal circle above your head. Two observations are immediate:
\begin{enumerate}
\item You must \textbf{keep pulling} on the string the whole time. The instant you let go, the stone flies away.
\item Your hand feels a continuous pull \emph{outwards} from the stone — the stone is ``trying'' to go straight while you force it to turn.
\end{enumerate}
A force is needed to keep the stone on its circle. By Newton's second law, a force means an acceleration. So the stone \emph{is} accelerating, even though its speed never changes.
\end{experience}

The key is the vector nature of velocity. Acceleration is defined by
\[
\vec{a} = \frac{d\vec{v}}{dt},
\]
and the velocity vector $\vec{v}$ has \textbf{both a magnitude and a direction}. In uniform circular motion the magnitude $v$ is constant, but the direction of $\vec{v}$ changes continuously: at every instant $\vec{v}$ is tangent to the circle, and the tangent direction rotates as the body moves. A changing velocity vector — even with constant magnitude — means a non-zero acceleration.

\begin{attention}[Constant speed does not mean zero acceleration]
In circular motion, ``uniform'' refers to the \emph{speed} (the magnitude of $\vec{v}$), not to the velocity vector. Because the \emph{direction} of $\vec{v}$ changes all the time, the acceleration is \textbf{never zero}. Writing ``$v$ constant $\Rightarrow a = 0$'' is one of the most frequent errors at the GCE Advanced Level.
\end{attention}

\courssub{Direction of the acceleration: towards the centre}

Consider a particle moving with constant speed $v$ on a circle of radius $r$ centred at $O$. Let it pass through point $P$ at time $t$ and through a nearby point $Q$ at time $t + \Delta t$, where the radius has turned through a small angle $\Delta\theta$.

\begin{popfigure}
\begin{center}
\begin{tikzpicture}[scale=1.1]
% Circle
\draw[popline, thick] (0,0) circle (2.6);
\fill (0,0) circle (1.5pt) node[below right] {$O$};
% Points P and Q
\coordinate (P) at (2.6,0);
\coordinate (Q) at (50:2.6);
\fill[popblue] (P) circle (2pt) node[right] {$P$};
\fill[popblue] (Q) circle (2pt) node[above right] {$Q$};
% Radii
\draw[popmuted] (0,0) -- (P);
\draw[popmuted] (0,0) -- (Q);
\draw[popmuted] (0.7,0) arc (0:50:0.7) node[midway, right, xshift=2pt] {$\Delta\theta$};
% Velocity vectors (tangents)
\draw[-latex, very thick, poporange] (P) -- ++(0,1.6) node[right] {$\vec{v}_P$};
\draw[-latex, very thick, poporange] (Q) -- ++(140:1.6) node[above left] {$\vec{v}_Q$};
% Velocity triangle (inset)
\begin{scope}[xshift=5.6cm, yshift=-0.4cm]
\coordinate (A) at (0,0);
\coordinate (B) at (0,1.6);
\coordinate (C) at (50:1.6);
\draw[-latex, very thick, poporange] (A) -- (B) node[right] {$\vec{v}_P$};
\draw[-latex, very thick, poporange] (A) -- (C) node[above left] {$\vec{v}_Q$};
\draw[-latex, very thick, popgreen] (B) -- (C) node[midway, above, sloped] {$\Delta\vec{v}$};
\draw[popmuted] (0,0.45) arc (90:50:0.45) node[midway, above] {\scriptsize $\Delta\theta$};
\node[popdark] at (0.5,-0.7) {$\Delta\vec{v} = \vec{v}_Q - \vec{v}_P$};
\end{scope}
\end{tikzpicture}
\end{center}
\legende{Velocity vectors at two nearby points $P$ and $Q$; the change $\Delta\vec{v}$ points towards the interior of the circle, i.e.\ towards the centre as $\Delta\theta \to 0$.}
\end{popfigure}

To find the acceleration we must examine the change of velocity $\Delta\vec{v} = \vec{v}_Q - \vec{v}_P$. Place the two velocity vectors tail to tail (right-hand part of the figure). Since the motion is uniform, $|\vec{v}_P| = |\vec{v}_Q| = v$, so the velocity triangle is isosceles, and the angle between $\vec{v}_P$ and $\vec{v}_Q$ equals the angle $\Delta\theta$ between the radii (each velocity is perpendicular to its radius). As $\Delta\theta \to 0$, the vector $\Delta\vec{v}$ becomes perpendicular to $\vec{v}_P$ and points \emph{inwards}, along the radius, \textbf{towards the centre $O$}. This is why the acceleration is called \cle{centripetal}, from the Latin \emph{petere}: ``seeking the centre''.

\courssub{Derivation of $a = \dfrac{v^{2}}{r}$}

This derivation is instructive and its steps can be asked for at the examination.

\textbf{Step 1: magnitude of $\Delta\vec{v}$.} In the isosceles velocity triangle with sides $v$, $v$ and small angle $\Delta\theta$ between them, the chord gives
\[
|\Delta\vec{v}| = 2v\sin\!\left(\frac{\Delta\theta}{2}\right) \approx v\,\Delta\theta \qquad (\Delta\theta \text{ small, in radians}).
\]

\textbf{Step 2: time taken.} The body moves along the arc from $P$ to $Q$ with speed $v$, and the arc length is $r\,\Delta\theta$, so
\[
\Delta t = \frac{r\,\Delta\theta}{v}.
\]

\textbf{Step 3: magnitude of the acceleration.}
\[
a = \lim_{\Delta t \to 0}\frac{|\Delta\vec{v}|}{\Delta t} = \frac{v\,\Delta\theta}{r\,\Delta\theta / v} = \frac{v^{2}}{r}.
\]

Since $v = r\omega$, we obtain the equivalent form $a = \dfrac{(r\omega)^{2}}{r} = r\omega^{2}$.

\begin{definition}[Centripetal acceleration]
A body in uniform circular motion of radius $r$ with speed $v$ and angular speed $\omega$ has an acceleration of magnitude
\[
\boxed{\,a = \frac{v^{2}}{r} = r\omega^{2}\,}
\]
directed at every instant along the radius \textbf{towards the centre} of the circle. It is called the \cle{centripetal acceleration}. Its SI unit is the $\mathrm{m\,s^{-2}}$.
\end{definition}

\begin{attention}[Two traps with the formula]
\begin{itemize}
\item $a$ is \emph{inversely} proportional to $r$ at fixed $v$, but \emph{directly} proportional to $r$ at fixed $\omega$ (use $a = r\omega^{2}$). Always ask which quantity is held constant before comparing.
\item The acceleration points towards the centre, \textbf{not} along the tangent. The velocity points along the tangent; velocity and acceleration are perpendicular in uniform circular motion.
\end{itemize}
\end{attention}

\courssub{Centripetal force: Newton's second law for circular motion}

Acceleration never appears by magic: it is produced by a resultant force. Applying Newton's second law $\vec{F} = m\vec{a}$ to the centripetal acceleration gives the central law of this chapter.

\begin{propriete}[Centripetal force (Newton's second law for circular motion)]
To keep a body of mass $m$ in uniform circular motion of radius $r$ at speed $v$ (angular speed $\omega$), the \textbf{resultant} of all the forces acting on it must have magnitude
\[
\boxed{\,F = \frac{mv^{2}}{r} = mr\omega^{2}\,}
\]
and must be directed towards the centre of the circle. This resultant is called the \cle{centripetal force}.
\end{propriete}

\begin{aretenir}[Centripetal force is NOT a new force]
The centripetal force is not an extra force to be added to the list of forces. It is the \emph{role played by the resultant} of real, identifiable forces:
\begin{itemize}
\item the \textbf{tension} of a string (stone whirled in a circle);
\item the \textbf{friction} between tyres and road (car rounding a bend);
\item the \textbf{gravitational attraction} of the Earth (satellite or the Moon in orbit);
\item the \textbf{normal reaction} of a wall or banked track (wall of death, velodrome).
\end{itemize}
In each problem, your first task is to identify \emph{which real force(s)} supply the required $\dfrac{mv^{2}}{r}$.
\end{aretenir}

\begin{attention}[There is no ``centrifugal force'' on the body]
You may feel ``thrown outwards'' in a turning car, but no real force pushes the moving body away from the centre. The so-called \emph{centrifugal force} is only the tendency of the body to continue in a straight line (inertia), felt in the rotating frame. \textbf{Never draw a centrifugal force on a free-body diagram} — it is a classic error that costs marks.
\end{attention}

\begin{methode}[Finding the required centripetal force]
\begin{enumerate}
\item Draw the \textbf{free-body diagram}: list all real forces (weight, tension, friction, reaction, \dots).
\item Choose the \textbf{radial direction} (towards the centre) and resolve every force along it.
\item Set the \textbf{resultant towards the centre} equal to $\dfrac{mv^{2}}{r}$ (or $mr\omega^{2}$).
\item Solve for the unknown (tension, friction, speed, \dots) and check the units.
\end{enumerate}
\end{methode}

\courssub{Worked examples}

\begin{exemplebox}[Stone on a string (horizontal circle)]
A stone of mass $m = 0.50\ \mathrm{kg}$ is whirled in a horizontal circle of radius $r = 1.0\ \mathrm{m}$ on a frictionless table, at speed $v = 4.0\ \mathrm{m\,s^{-1}}$. Find the tension in the string.

\textbf{Solution.} On the smooth table the weight is balanced by the normal reaction of the table; the only horizontal force is the tension $T$, which points towards the centre. Hence the tension alone provides the centripetal force:
\[
T = \frac{mv^{2}}{r} = \frac{0.50 \times (4.0)^{2}}{1.0} = 8.0\ \mathrm{N}.
\]
If the string can bear at most $8.0\ \mathrm{N}$, it is exactly at breaking point: any increase of speed snaps it.
\end{exemplebox}

\begin{popfigure}
\begin{center}
\begin{tikzpicture}[scale=1.15]
% Table top hint
\draw[popline] (-3.4,-0.9) -- (3.4,-0.9);
% Circle path
\draw[popline, dashed, thick] (0,0) circle (2.2);
% Centre peg
\fill[popdark] (0,0) circle (2pt) node[below left] {$O$};
% Stone
\coordinate (S) at (2.2,0);
\fill[popblue] (S) circle (5pt);
\node[popblue, right] at (S) {stone $m$};
% String
\draw[popdark, thick] (0,0) -- (S) node[midway, above] {string, radius $r$};
% Tension towards centre
\draw[-latex, very thick, poporange] (S) -- ++(-1.5,0) node[above left] {$\vec{T}$ (centripetal)};
% Velocity tangent
\draw[-latex, very thick, popgreen] (S) -- ++(0,1.4) node[right] {$\vec{v}$};
\end{tikzpicture}
\end{center}
\legende{Free-body diagram (view from above) of a stone on a horizontal string: the tension $\vec{T}$ points towards the centre $O$ and provides the centripetal force; the velocity $\vec{v}$ is tangential.}
\end{popfigure}

\begin{exoresolu}[Car rounding a bend]
A car of mass $m = 1200\ \mathrm{kg}$ rounds a flat, unbanked bend of radius $r = 50\ \mathrm{m}$ at $v = 20\ \mathrm{m\,s^{-1}}$. (a) What centripetal force is required, and what provides it? (b) If the maximum friction available between tyres and road is $F_{\max} = 9600\ \mathrm{N}$, what is the greatest safe speed?
\tcblower
\textbf{(a)} On a flat road the weight and the normal reaction cancel vertically. Horizontally, only the \textbf{friction} $\vec{f}$ of the road on the tyres points towards the centre of the bend, so friction is the centripetal force:
\[
f = \frac{mv^{2}}{r} = \frac{1200 \times 20^{2}}{50} = \frac{1200 \times 400}{50} = 9600\ \mathrm{N}.
\]
\textbf{(b)} The greatest safe speed corresponds to friction at its maximum:
\[
\frac{mv_{\max}^{2}}{r} = F_{\max}
\quad\Longrightarrow\quad
v_{\max} = \sqrt{\frac{F_{\max}\, r}{m}} = \sqrt{\frac{9600 \times 50}{1200}} = \sqrt{400} = 20\ \mathrm{m\,s^{-1}}.
\]
So at $20\ \mathrm{m\,s^{-1}}$ (about $72\ \mathrm{km\,h^{-1}}$) the car is exactly at the limit of grip — on a wet Douala--Yaound\'e road, where $F_{\max}$ is smaller, the safe speed drops accordingly.
\end{exoresolu}

\begin{exemplebox}[Satellite in orbit]
A satellite of mass $m = 800\ \mathrm{kg}$ moves in a circular orbit of radius $r = 7.0\times 10^{6}\ \mathrm{m}$ around the Earth, where the gravitational pull on it is $F = 6.4\times 10^{3}\ \mathrm{N}$. Find its orbital speed.

\textbf{Solution.} The Earth's gravitational attraction is the only force on the satellite, so it plays the role of centripetal force:
\[
\frac{mv^{2}}{r} = F
\quad\Longrightarrow\quad
v = \sqrt{\frac{Fr}{m}} = \sqrt{\frac{6.4\times 10^{3} \times 7.0\times 10^{6}}{800}} = \sqrt{5.6\times 10^{7}} \approx 7.5\times 10^{3}\ \mathrm{m\,s^{-1}}.
\]
The satellite falls towards the Earth continuously — that fall \emph{is} its centripetal acceleration — yet its tangential speed keeps it on the circle.
\end{exemplebox}

\courssub{What happens if the centripetal force is removed?}

Suppose the string of the whirling stone suddenly breaks. From that instant, no horizontal force acts on the stone. By \textbf{Newton's first law}, the stone continues in a straight line at constant velocity — the velocity it had at the moment of release, which is \emph{tangent} to the circle. The stone therefore flies off \textbf{along the tangent}, never radially outwards. This is why sparks from a grinding wheel and mud from a bicycle tyre leave tangentially, and why a hammer thrower releases the hammer so that its tangent line points at the target.

\begin{savaistu}[Why the Moon does not fall on us]
The Moon is permanently ``falling'' towards the Earth: gravity gives it the centripetal acceleration $a = v^{2}/r$ of its (nearly) circular orbit. It never gets closer because its tangential speed continually carries it sideways as it falls. Newton himself used this idea to show that the same gravity that makes an apple fall also holds the Moon in orbit.
\end{savaistu}

\begin{exercice}[Practice]
\begin{enumerate}
\item A particle moves uniformly on a circle of radius $2.0\ \mathrm{m}$ with angular speed $\omega = 3.0\ \mathrm{rad\,s^{-1}}$. Calculate its centripetal acceleration.
\item A child of mass $30\ \mathrm{kg}$ sits $4.0\ \mathrm{m}$ from the axis of a merry-go-round making one revolution every $6.0\ \mathrm{s}$. Find the centripetal force the seat must exert on the child.
\item A bucket of water is whirled in a vertical circle. Explain, using the ideas of this section, why the water does not fall out at the top if the speed is large enough, and identify what provides the centripetal force on the water at the top.
\end{enumerate}
\end{exercice}

\begin{corrige}[Exercise 1]
$a = r\omega^{2} = 2.0 \times (3.0)^{2} = 18\ \mathrm{m\,s^{-2}}$, directed towards the centre.
\end{corrige}

\begin{corrige}[Exercise 2]
$\omega = \dfrac{2\pi}{T} = \dfrac{2\pi}{6.0} = \dfrac{\pi}{3}\ \mathrm{rad\,s^{-1}}$. Then
$F = mr\omega^{2} = 30 \times 4.0 \times \left(\dfrac{\pi}{3}\right)^{2} = \dfrac{120\pi^{2}}{9} \approx 132\ \mathrm{N}$, directed towards the axis.
\end{corrige}

\begin{corrige}[Exercise 3]
At the top of the circle, both the weight $mg$ of the water and the reaction of the bucket bottom point downwards, i.e.\ towards the centre: together they supply the centripetal force, $mg + R = \dfrac{mv^{2}}{r}$. If $v$ is large enough that $\dfrac{mv^{2}}{r} \geq mg$, the required force exceeds the weight, so $R \geq 0$: the bucket keeps pushing on the water and the water stays in. The minimum speed is given by $R = 0$, i.e.\ $v_{\min} = \sqrt{gr}$.
\end{corrige}

\begin{aretenir}[Summary of the section]
\begin{itemize}
\item Uniform circular motion \emph{is} accelerated motion, because the direction of $\vec{v}$ changes continuously.
\item The acceleration is centripetal: $a = \dfrac{v^{2}}{r} = r\omega^{2}$, towards the centre.
\item Newton's second law gives the centripetal force: $F = \dfrac{mv^{2}}{r} = mr\omega^{2}$, supplied by real forces (tension, friction, gravity, reaction).
\item No centrifugal force acts on the body; if the centripetal force vanishes, the body leaves along the tangent (Newton's first law).
\end{itemize}
\end{aretenir}

\coursec{Linear Speed, Angular Speed and Their Relationship}

In the previous section we discovered that a particle moving in a circle, even at constant speed, is permanently accelerated towards the centre, with centripetal acceleration $a = \dfrac{v^{2}}{r}$, and that a \cle{centripetal force} $F = \dfrac{mv^{2}}{r}$ is needed to bend its path. But we left one question open: how exactly are the \emph{linear} speed $v$ (in metres per second) and the \emph{angular} speed $\omega$ (in radians per second) connected? A cyclist in Bamenda, a grinding wheel in a Douala workshop, and the Earth itself all rotate --- yet different points of the same rotating object move at different linear speeds. This section makes the link precise: it is the celebrated relation $v = r\omega$, one of the most frequently used formulas of the whole mechanics syllabus.

\courssub{Linear (Tangential) Speed: Intuition}

Imagine tying a stone to a string and whirling it horizontally above your head. If the string suddenly snaps, the stone does \emph{not} fly radially outwards --- it flies off along the \cle{tangent} to the circle at the point of release. This everyday observation tells us two things:

\begin{itemize}
\item the stone possesses a \textbf{linear speed} $v$, measured as arc length covered per unit time, in $\mathrm{m\,s^{-1}}$;
\item at each instant, the velocity vector $\vec{v}$ is directed along the \textbf{tangent} to the circle, perpendicular to the radius.
\end{itemize}

\begin{definition}[Linear (tangential) speed]
The \cle{linear speed} (or tangential speed) of a particle in uniform circular motion is the distance travelled \emph{along the arc} per unit time:
\[ v = \frac{\text{arc length}}{\text{time}} = \frac{s}{t}. \]
The velocity vector $\vec{v}$ is tangent to the circular path at the position of the particle, and its magnitude $v$ is constant in uniform circular motion (only its \emph{direction} changes).
\end{definition}

Recall also the angular quantities from Section 1: the angular position $\theta$ (in radians), and the angular speed
\[ \omega = \frac{d\theta}{dt}, \qquad \text{in } \mathrm{rad\,s^{-1}}, \]
with $\omega = \dfrac{2\pi}{T} = 2\pi f$, where $T$ is the period and $f$ the frequency.

\courssub{Derivation of $v = r\omega$}

This derivation is short, instructive, and \textbf{examinable} --- learn it line by line.

For a particle moving on a circle of radius $r$, the arc length $s$ corresponding to an angle $\theta$ (in radians) is, by the very definition of the radian:
\[ s = r\theta. \]
Since the radius $r$ is constant, differentiate both sides with respect to time $t$:
\[ \frac{ds}{dt} = \frac{d}{dt}(r\theta) = r\,\frac{d\theta}{dt}. \]
Now $\dfrac{ds}{dt}$ is the rate at which arc length is covered, i.e.\ the linear speed $v$, and $\dfrac{d\theta}{dt}$ is the angular speed $\omega$. Therefore:
\[ \boxed{\,v = r\omega\,} \]

\begin{propriete}[Linking linear and angular speed]
For a particle in uniform circular motion of radius $r$, period $T$ and frequency $f$:
\[ v = r\omega = \frac{2\pi r}{T} = 2\pi r f. \]
The two last forms follow immediately from $\omega = \dfrac{2\pi}{T} = 2\pi f$. \emph{(The derivation of $v = r\omega$ from $s = r\theta$ is examinable.)}
\end{propriete}

\begin{attention}[Units are not optional]
The formula $v = r\omega$ is valid \textbf{only when $\theta$ and $\omega$ are in radians} (radians and radians per second). If a problem gives the rotation in revolutions per minute (rpm) or in degrees, convert first:
\[ \omega\ (\mathrm{rad\,s^{-1}}) = \frac{2\pi \times N\ (\mathrm{rpm})}{60}. \]
Using rpm directly in $v = r\omega$ is one of the most common errors at the GCE.
\end{attention}

\begin{popfigure}
\begin{center}
\begin{tikzpicture}[scale=1.6]
\draw[thick, popblue] (0,0) circle (2);
\draw[->, thick, popdark] (0,0) -- (60:2) node[midway, above left, popdark] {$r$};
\filldraw[poporange] (60:2) circle (0.06) node[right=2pt, popink] {$P$};
\draw[->, very thick, poporange] (60:2) -- ++(150:1.3) node[above left, poporange] {$\vec{v}$};
\draw[->, thick, popgreen] (0.55,0) arc (0:60:0.55) node[midway, right, popgreen] {$\theta$};
\draw[->, thick, poppurple] (2.45,0.35) arc (20:80:0.6) node[midway, right, poppurple] {$\omega$};
\filldraw (0,0) circle (0.05) node[below left, popink] {$O$};
\node[popink, align=center] at (0,-2.45) {$v = r\omega$,\ $\vec{v} \perp$ radius};
\end{tikzpicture}
\legende{Uniform circular motion: the velocity $\vec{v}$ is tangent to the circle, perpendicular to the radius $OP$, with $v = r\omega$.}
\end{center}
\end{popfigure}

\courssub{Same $\omega$, Different $v$: Points of a Rigid Rotating Body}

Here is the key conceptual point of the whole section. Consider a rigid body rotating about a fixed axis --- a grinding wheel, a turntable, a ceiling fan. In a time $t$, \emph{every} radius of the body sweeps through the \emph{same} angle $\theta$: the body rotates ``as one block''. Hence all points of the body share the \textbf{same angular speed} $\omega$.

But their linear speeds differ, because $v = r\omega$ grows in proportion to the distance $r$ from the axis: a point twice as far from the centre moves twice as fast. A point \emph{on} the axis ($r = 0$) does not move at all.

\begin{attention}[The classic trap: $v$ depends on $r$, $\omega$ does not]
Two points $A$ and $B$ on the same rotating wheel, with $r_A < r_B$:
\begin{itemize}
\item have the \textbf{same} angular speed: $\omega_A = \omega_B$;
\item have \textbf{different} linear speeds: $v_B = \dfrac{r_B}{r_A}\,v_A > v_A$;
\item have different centripetal accelerations: $a = r\omega^{2}$ is larger for the outer point.
\end{itemize}
Saying ``all points of the wheel have the same speed'' is false; only the \emph{angular} speed is common.
\end{attention}

\begin{popfigure}
\begin{center}
\begin{tikzpicture}[scale=1.5]
\draw[thick, popblue, fill=popblueL] (0,0) circle (2);
\draw[dashed, popmuted] (0,0) circle (0.9);
\draw[->, thick, popdark] (0,0) -- (30:0.9) node[midway, above, popdark] {$r_1$};
\draw[->, thick, popdark] (0,0) -- (30:2) node[midway, above left, pos=0.8, popdark] {$r_2$};
\filldraw[popgreen] (30:0.9) circle (0.06) node[below right=1pt, popink] {$A$};
\filldraw[poporange] (30:2) circle (0.06) node[right=1pt, popink] {$B$};
\draw[->, very thick, popgreen] (30:0.9) -- ++(120:0.65) node[above left, popgreen] {$\vec{v}_1$};
\draw[->, very thick, poporange] (30:2) -- ++(120:1.4) node[above left, poporange] {$\vec{v}_2$};
\draw[->, thick, poppurple] (-2.4,0.4) arc (160:100:0.65) node[midway, left, poppurple] {$\omega$};
\filldraw (0,0) circle (0.05) node[below, popink] {$O$};
\node[popink, align=center] at (0,-2.5) {Same $\omega$;\ $v_2 = \dfrac{r_2}{r_1}v_1 > v_1$};
\end{tikzpicture}
\legende{A rotating disc: points $A$ and $B$ share the same angular speed $\omega$, but the outer point $B$ has the larger tangential velocity since $v = r\omega$.}
\end{center}
\end{popfigure}

\courssub{Choosing Between $a = \dfrac{v^{2}}{r}$ and $a = r\omega^{2}$}

Since $v = r\omega$, the centripetal acceleration of Section 2 can be written in two equivalent ways:
\[ a = \frac{v^{2}}{r} = r\omega^{2} \qquad \left(\text{also } a = \frac{4\pi^{2}r}{T^{2}} = 4\pi^{2} r f^{2}\right). \]
Both formulas describe the \emph{same} acceleration; the art is to pick the one that uses the data actually given, avoiding unnecessary intermediate calculations.

\begin{methode}[Which centripetal acceleration formula?]
\begin{enumerate}
\item \textbf{Read the data.} List what is given: $v$? $\omega$? $T$ or $f$? $r$?
\item \textbf{Given $v$ and $r$:} use $a = \dfrac{v^{2}}{r}$ directly.
\item \textbf{Given $\omega$ and $r$:} use $a = r\omega^{2}$ directly.
\item \textbf{Given $T$ or $f$:} convert first, $\omega = \dfrac{2\pi}{T} = 2\pi f$, then use $a = r\omega^{2}$.
\item \textbf{Comparing two points of the same rigid body} (same $\omega$): use $a = r\omega^{2}$, so $a \propto r$.
\item \textbf{Comparing motions at the same linear speed} (e.g.\ a belt driving two pulleys): use $a = \dfrac{v^{2}}{r}$, so $a \propto \dfrac{1}{r}$.
\end{enumerate}
\end{methode}

\begin{attention}[A subtle comparison trap]
For two points on the \emph{same wheel}, $a = r\omega^{2}$ shows the outer point has the \emph{larger} acceleration. But for a car taking a \emph{tighter} bend at the \emph{same speed} $v$, $a = \dfrac{v^{2}}{r}$ shows the smaller radius gives the larger acceleration. There is no contradiction: in the first case $\omega$ is fixed, in the second $v$ is fixed. Always ask: \emph{which quantity is common to the two situations?}
\end{attention}

\courssub{Worked Examples}

\begin{exoresolu}[Points on a grinding wheel]
A grinding wheel in a workshop rotates at $1200\ \mathrm{rpm}$. Two grains of the wheel, $A$ and $B$, are at distances $r_A = 4.0\ \mathrm{cm}$ and $r_B = 10.0\ \mathrm{cm}$ from the axis. Find (a) the angular speed of the wheel; (b) the linear speed of each grain; (c) the centripetal acceleration of grain $B$.
\tcblower
\textbf{Solution.}
(a) Convert the rotation rate to $\mathrm{rad\,s^{-1}}$:
\[ \omega = \frac{2\pi \times 1200}{60} = 40\pi \approx 125.7\ \mathrm{rad\,s^{-1}}. \]
This angular speed is the \emph{same} for both grains.

(b) Apply $v = r\omega$ to each grain:
\[ v_A = r_A\omega = 0.040 \times 40\pi = 1.6\pi \approx 5.0\ \mathrm{m\,s^{-1}}, \]
\[ v_B = r_B\omega = 0.100 \times 40\pi = 4\pi \approx 12.6\ \mathrm{m\,s^{-1}}. \]
Note that $\dfrac{v_B}{v_A} = \dfrac{r_B}{r_A} = 2.5$, as expected.

(c) Since $\omega$ and $r$ are known, use $a = r\omega^{2}$:
\[ a_B = 0.100 \times (40\pi)^{2} = 160\pi^{2} \approx 1.58\times 10^{3}\ \mathrm{m\,s^{-2}}. \]
That is more than $160$ times $g$ --- which is why a cracked grinding wheel is so dangerous.
\end{exoresolu}

\begin{exoresolu}[A record turntable]
A turntable rotates at $33\frac{1}{3}\ \mathrm{rpm}$. A speck of dust sits on a record at $r = 12\ \mathrm{cm}$ from the centre. Calculate (a) the period $T$; (b) the linear speed of the speck; (c) its centripetal acceleration.
\tcblower
\textbf{Solution.}
(a) One revolution takes
\[ T = \frac{60}{33\tfrac{1}{3}} = \frac{60 \times 3}{100} = 1.8\ \mathrm{s}. \]
(b) Using $v = \dfrac{2\pi r}{T}$:
\[ v = \frac{2\pi \times 0.12}{1.8} = \frac{0.24\pi}{1.8} \approx 0.42\ \mathrm{m\,s^{-1}}. \]
(c) With $\omega = \dfrac{2\pi}{T} = \dfrac{2\pi}{1.8} \approx 3.49\ \mathrm{rad\,s^{-1}}$:
\[ a = r\omega^{2} = 0.12 \times (3.49)^{2} \approx 1.46\ \mathrm{m\,s^{-2}}. \]
\end{exoresolu}

\begin{exoresolu}[Speed of a point on the Equator]
The Earth completes one rotation about its axis in $24$ hours, and its radius is $R \approx 6.4\times 10^{6}\ \mathrm{m}$. Estimate the linear speed, due to the Earth's rotation, of a person standing on the Equator (e.g.\ near the coast of the South-West Region).
\tcblower
\textbf{Solution.} The period is $T = 24 \times 3600 = 8.64\times 10^{4}\ \mathrm{s}$. Then
\[ v = \frac{2\pi R}{T} = \frac{2\pi \times 6.4\times 10^{6}}{8.64\times 10^{4}} \approx 4.65\times 10^{2}\ \mathrm{m\,s^{-1}}. \]
A person on the Equator is therefore travelling at roughly $465\ \mathrm{m\,s^{-1}}$ --- faster than the speed of sound in air --- without feeling it, because the motion is perfectly uniform and everything around moves together. (Points away from the Equator move on smaller circles of radius $R\cos\lambda$, where $\lambda$ is the latitude, so their speed is smaller while $\omega$ is the same for the whole planet.)
\end{exoresolu}

\courssub{Inverse Problems: Finding $\omega$ or $r$}

The relation $v = r\omega$ contains three quantities; given any two, the third follows:
\[ \omega = \frac{v}{r}, \qquad r = \frac{v}{\omega}, \qquad v = r\omega. \]
Combined with $\omega = \dfrac{2\pi}{T} = 2\pi f$ and $a = \dfrac{v^{2}}{r} = r\omega^{2}$, this solves a wide family of GCE problems.

\begin{exoresolu}[Finding $\omega$ from $v$ and $r$]
The tip of the minute hand of a large clock moves at $v = 2.0\ \mathrm{mm\,s^{-1}}$. The hand is $r = 15\ \mathrm{cm}$ long. (a) Find the angular speed of the tip. (b) Deduce the period of the minute hand and comment.
\tcblower
\textbf{Solution.}
(a) Convert units: $v = 2.0\times 10^{-3}\ \mathrm{m\,s^{-1}}$, $r = 0.15\ \mathrm{m}$.
\[ \omega = \frac{v}{r} = \frac{2.0\times 10^{-3}}{0.15} \approx 1.33\times 10^{-2}\ \mathrm{rad\,s^{-1}}. \]
(b) $T = \dfrac{2\pi}{\omega} = \dfrac{2\pi}{1.33\times 10^{-2}} \approx 4.7\times 10^{2}\ \mathrm{s}$? Let us check: $\dfrac{2\pi}{0.0133} \approx 472\ \mathrm{s}$. Hmm --- a minute hand should take $3600\ \mathrm{s}$! The discrepancy tells us the stated tip speed was too large; a real $15\ \mathrm{cm}$ minute hand has $\omega = \dfrac{2\pi}{3600} \approx 1.75\times 10^{-3}\ \mathrm{rad\,s^{-1}}$ and tip speed $v = r\omega \approx 0.26\ \mathrm{mm\,s^{-1}}$. \textbf{Lesson:} always check the physical plausibility of a numerical answer --- examiners reward this reflex.
\end{exoresolu}

\begin{exoresolu}[Finding $r$ from $v$ and $T$]
A satellite moves in a circular orbit with period $T = 5.6\times 10^{3}\ \mathrm{s}$ at a constant speed $v = 7.6\times 10^{3}\ \mathrm{m\,s^{-1}}$. Find the radius of its orbit.
\tcblower
\textbf{Solution.} From $v = \dfrac{2\pi r}{T}$ we isolate $r$:
\[ r = \frac{vT}{2\pi} = \frac{7.6\times 10^{3} \times 5.6\times 10^{3}}{2\pi} \approx 6.8\times 10^{6}\ \mathrm{m}. \]
The orbit radius is about $6.8\times 10^{6}\ \mathrm{m}$, i.e.\ roughly $400\ \mathrm{km}$ above the Earth's surface.
\end{exoresolu}

\begin{methode}[Exam strategy: period or frequency given]
When a question gives the \textbf{period $T$} or the \textbf{frequency $f$} (or a rotation rate in rpm):
\begin{enumerate}
\item Convert immediately to angular speed: $\omega = \dfrac{2\pi}{T} = 2\pi f = \dfrac{2\pi N}{60}$.
\item Then use $v = r\omega$ and, if required, $a = r\omega^{2}$ or $F = mr\omega^{2}$.
\item Keep $\pi$ symbolic until the final line, then approximate ($\pi \approx 3.1416$, $\pi^{2} \approx 9.87$).
\item Check units: $r$ in metres, $\omega$ in $\mathrm{rad\,s^{-1}}$, $v$ in $\mathrm{m\,s^{-1}}$.
\end{enumerate}
\end{methode}

\begin{aretenir}[Key points of the section]
\begin{itemize}
\item Linear speed: $v = \dfrac{s}{t}$ along the arc; $\vec{v}$ is \textbf{tangent} to the circle.
\item Fundamental link (derivable from $s = r\theta$): $\;v = r\omega = \dfrac{2\pi r}{T} = 2\pi r f$.
\item All points of a rigid rotating body share the same $\omega$; their linear speeds satisfy $v \propto r$.
\item Centripetal acceleration: $a = \dfrac{v^{2}}{r} = r\omega^{2}$ --- choose the form matching the given data.
\item Inverse problems: $\omega = \dfrac{v}{r}$, $r = \dfrac{v}{\omega} = \dfrac{vT}{2\pi}$.
\item Always convert rpm, periods and frequencies to $\omega$ in $\mathrm{rad\,s^{-1}}$ \emph{before} computing.
\end{itemize}
\end{aretenir}

\begin{exercice}[Practice]
\begin{enumerate}
\item A ceiling fan rotates at $300\ \mathrm{rpm}$. Find (a) its angular speed in $\mathrm{rad\,s^{-1}}$; (b) the linear speed of a point on a blade $60\ \mathrm{cm}$ from the axis; (c) the centripetal acceleration of that point.
\item Two points $P$ and $Q$ of a rotating platform are at $0.20\ \mathrm{m}$ and $0.50\ \mathrm{m}$ from the axis. If $v_P = 1.5\ \mathrm{m\,s^{-1}}$, find $\omega$ and $v_Q$.
\item A car wheel of radius $0.30\ \mathrm{m}$ rolls so that its centre moves at $24\ \mathrm{m\,s^{-1}}$. Assuming no slipping, find the angular speed of the wheel and its rotation frequency in rpm.
\end{enumerate}
\end{exercice}

\begin{corrige}[Exercise 1]
\textbf{1.} (a) $\omega = \dfrac{2\pi \times 300}{60} = 10\pi \approx 31.4\ \mathrm{rad\,s^{-1}}$. (b) $v = r\omega = 0.60 \times 10\pi = 6\pi \approx 18.8\ \mathrm{m\,s^{-1}}$. (c) $a = r\omega^{2} = 0.60 \times (10\pi)^{2} = 60\pi^{2} \approx 5.9\times 10^{2}\ \mathrm{m\,s^{-2}}$.

\textbf{2.} $\omega = \dfrac{v_P}{r_P} = \dfrac{1.5}{0.20} = 7.5\ \mathrm{rad\,s^{-1}}$ (same for the whole platform); $v_Q = r_Q\omega = 0.50 \times 7.5 = 3.75\ \mathrm{m\,s^{-1}}$.

\textbf{3.} $\omega = \dfrac{v}{r} = \dfrac{24}{0.30} = 80\ \mathrm{rad\,s^{-1}}$; $N = \dfrac{60\,\omega}{2\pi} = \dfrac{60 \times 80}{2\pi} \approx 764\ \mathrm{rpm}$.
\end{corrige}

\begin{savaistu}[Why this matters beyond the classroom]
The relation $v = r\omega$ explains why the teeth at the rim of a circular saw must be made of harder material than the disc itself: they travel fastest and do all the cutting. It also governs the design of the turbines at the Song Loulou and Memve'ele hydroelectric stations: for a given rotational speed imposed by the alternator, the engineer chooses the rotor radius so that the tangential speed of the blades stays within safe mechanical limits --- too large an $r$ at fixed $\omega$ means destructive centripetal stresses.
\end{savaistu}

\coursec{Applications of Linear and Angular Speed}

In the previous sections we established the fundamental bridge between the two descriptions of uniform circular motion: the \cle{linear speed} $v$ of a particle on its circular path and the \cle{angular speed} $\omega$ of the radius vector, linked by the relation $v = r\omega$. We also proved that any body moving in a circle of radius $r$ at constant speed $v$ experiences a \cle{centripetal acceleration} $a = \dfrac{v^{2}}{r} = r\omega^{2}$ directed towards the centre, and therefore requires a \cle{centripetal force} $F = \dfrac{mv^{2}}{r} = mr\omega^{2}$.

A natural question now arises: \emph{where does this centripetal force come from in real situations?} A force never appears from nowhere — it must be supplied by a concrete physical interaction: friction between tyres and road, tension in a string or belt, the normal reaction of a surface, or the pull of gravity. In this section we tour the most important applications. For each, we follow the same three-step method that GCE examiners expect:
\begin{enumerate}
\item \textbf{Identify} the force (or force component) that provides the centripetal force.
\item \textbf{Write} the equation of motion towards the centre: $F_{\text{centripetal}} = \dfrac{mv^{2}}{r} = mr\omega^{2}$.
\item \textbf{Impose} the physical limit (maximum friction, maximum tension, etc.) to find the condition for the motion to be possible.
\end{enumerate}

\courssub{Vehicles Rounding a Flat (Unbanked) Curve}

\textbf{Intuition.} Every driver in Douala or Bamenda knows that you must slow down before a sharp bend, especially when the road is wet. Why? Because on a \emph{flat} curve, the only horizontal force available to pull the vehicle towards the centre of the circle is the \cle{static friction} between the tyres and the road surface. If the required centripetal force exceeds the maximum available friction, the vehicle skids outwards.

\textbf{Analysis.} Consider a vehicle of mass $m$ rounding a flat curve of radius $r$ at constant speed $v$. The forces acting on it (viewed from behind) are:
\begin{itemize}
\item the weight $W = mg$, vertically downwards;
\item the normal reaction $R$ of the road, vertically upwards;
\item the friction force $f$, horizontal, pointing \emph{towards the centre} of the curve — this is the centripetal force.
\end{itemize}

Vertically there is no acceleration, so $R = mg$. Horizontally, Newton's second law gives
\[ f = \frac{mv^{2}}{r}. \]
But friction is limited by $f \le \mu R = \mu mg$, where $\mu$ is the coefficient of \emph{static} friction between tyres and road (static because the tyre does not slide relative to the road at the contact patch). The motion is therefore possible only if
\[ \frac{mv^{2}}{r} \le \mu mg. \]

\begin{popfigure}
\begin{center}
\begin{tikzpicture}[scale=1.0]
% road surface
\draw[thick, popmuted] (-3.5,0) -- (3.5,0);
% car as a block
\filldraw[popblueL] (-1.1,0) rectangle (1.1,1.1);
\draw[thick, popdark] (-1.1,0) rectangle (1.1,1.1);
% centre of mass
\filldraw[popink] (0,0.55) circle (0.06);
\node[popink, above right] at (0.05,0.6) {\small $G$};
% weight
\draw[-latex, very thick, popdark] (0,0.55) -- (0,-1.5) node[below] {\small $W = mg$};
% normal reaction
\draw[-latex, very thick, popblue] (0,0.55) -- (0,2.3) node[above] {\small $R$};
% friction towards centre (to the left)
\draw[-latex, very thick, poporange] (0,0.55) -- (-2.6,0.55) node[above left, align=center] {\small friction $f$\\ \small (centripetal)};
% centre of circle indication
\node[popmuted, align=center] at (-2.6,-0.5) {\small centre of\\ \small the curve};
\draw[popmuted, dashed] (-2.6,-0.15) -- (-2.6,0.4);
\end{tikzpicture}
\end{center}
\legende{Free-body diagram of a vehicle on a flat curve (viewed from behind). Friction $f$ towards the centre supplies the centripetal force; vertically $R = mg$.}
\end{popfigure}

\begin{propriete}[Maximum safe speed on a flat unbanked curve]
A vehicle of mass $m$ can round a flat curve of radius $r$ without skidding, the coefficient of static friction between tyres and road being $\mu$, only if its speed satisfies
\[ v \le v_{\max} = \sqrt{\mu r g}. \]
The proof is examinable: from $\dfrac{mv^{2}}{r} \le \mu mg$, cancel $m$ and rearrange to $v^{2} \le \mu r g$.
\end{propriete}

\begin{attention}[The mass cancels — and friction is the hero]
Two classic traps: \textbf{(1)} Students keep the mass $m$ in the final answer. It cancels: a heavy lorry and a light car have the \emph{same} maximum safe speed on the same curve (provided $\mu$ is the same). \textbf{(2)} Students say the car is ``thrown outwards by centrifugal force''. There is no outward force on the car in an inertial frame; the car skids because friction is \emph{insufficient} to bend its natural straight-line motion into a circle.
\end{attention}

\begin{exoresolu}[Maximum speed on a wet road]
A car rounds a flat curve of radius $50\ \mathrm{m}$. The coefficient of static friction between the tyres and the wet road is $0.4$. Taking $g = 10\ \mathrm{m\,s^{-2}}$, find the maximum speed at which the car can negotiate the curve without skidding.
\tcblower
\textbf{Solution.} The friction provides the centripetal force. At the limiting speed $\dfrac{mv_{\max}^{2}}{r} = \mu mg$, so
\[ v_{\max} = \sqrt{\mu r g} = \sqrt{0.4 \times 50 \times 10} = \sqrt{200} \approx 14.1\ \mathrm{m\,s^{-1}}. \]
Converting to $\mathrm{km\,h^{-1}}$: $14.1 \times 3.6 \approx 50.8\ \mathrm{km\,h^{-1}}$. The driver must enter the bend below this speed.
\end{exoresolu}

\courssub{Belt and Pulley Systems}

\textbf{Intuition.} In a grinding mill, a sewing machine, or a car's alternator drive, a motor turns a small pulley connected by a belt to a larger pulley. The belt does not stretch and does not slip, so \emph{every point of the belt moves with the same linear speed $v$}. Consequently the rim of each pulley moves at that same speed $v$ — but since the radii differ, the angular speeds differ.

\begin{popfigure}
\begin{center}
\begin{tikzpicture}[scale=1.0]
% pulley 1 (small)
\draw[thick, popblue, fill=popblueL] (0,0) circle (0.7);
\filldraw[popink] (0,0) circle (0.05);
\node[popink, below] at (0,-0.75) {\small $r_1$};
\draw[-latex, popblue, thick] (0,0) -- (0.5,0.5);
% pulley 2 (large)
\draw[thick, poporange, fill=popblueL] (4.5,0) circle (1.3);
\filldraw[popink] (4.5,0) circle (0.05);
\node[popink, below] at (4.5,-1.4) {\small $r_2$};
\draw[-latex, poporange, thick] (4.5,0) -- (5.4,0.9);
% belt (open belt, same sense of rotation)
\draw[very thick, popdark] (0,0.7) -- (4.5,1.3);
\draw[very thick, popdark] (0,-0.7) -- (4.5,-1.3);
% angular speed arrows (both counterclockwise for open belt)
\draw[-latex, thick, popblue] (1.0,0.35) arc (20:110:0.95);
\node[popblue] at (0,1.25) {\small $\omega_1$};
\draw[-latex, thick, poporange] (3.3,0.5) arc (160:250:1.6);
\node[poporange] at (3.1,-0.2) {\small $\omega_2$};
% belt speed label
\draw[-latex, thick, popgreen] (1.6,0.92) -- (2.6,1.08);
\node[popgreen, above] at (2.1,1.15) {\small belt speed $v$};
\end{tikzpicture}
\end{center}
\legende{Two pulleys joined by an open belt. The belt imposes the same linear speed $v$ on both rims, so $\omega_1 r_1 = \omega_2 r_2$: the smaller pulley spins faster. Both rotate in the same sense.}
\end{popfigure}

\begin{methode}[Solving belt/pulley problems]
\begin{enumerate}
\item Identify the radii $r_1$, $r_2$ and the given angular speed (convert rpm to $\mathrm{rad\,s^{-1}}$ using $\omega = \dfrac{2\pi N}{60}$ if needed).
\item Set the belt speeds equal: the rim of each pulley moves at the belt's speed, so
\[ \omega_1 r_1 = \omega_2 r_2. \]
\item Solve for the unknown angular speed, then convert back to rpm if required ($N = \dfrac{60\,\omega}{2\pi}$).
\item If a linear (belt) speed is asked, compute $v = r\omega$ on either pulley — both give the same value; use this as a check.
\end{enumerate}
\end{methode}

\begin{exoresolu}[Grinding mill pulleys]
A motor pulley of radius $5\ \mathrm{cm}$ rotates at $1200\ \mathrm{rpm}$ and drives, by a belt, a mill pulley of radius $20\ \mathrm{cm}$. Find (a) the angular speed of the mill pulley in $\mathrm{rad\,s^{-1}}$, (b) the linear speed of the belt.
\tcblower
\textbf{Solution.} (a) Convert the motor speed: $\omega_1 = \dfrac{2\pi \times 1200}{60} = 40\pi\ \mathrm{rad\,s^{-1}}$. Equal belt speeds give
\[ \omega_2 = \frac{\omega_1 r_1}{r_2} = \frac{40\pi \times 5}{20} = 10\pi \approx 31.4\ \mathrm{rad\,s^{-1}}. \]
(b) Belt speed: $v = r_1\omega_1 = 0.05 \times 40\pi = 2\pi \approx 6.28\ \mathrm{m\,s^{-1}}$. Check with the other pulley: $v = r_2\omega_2 = 0.20 \times 10\pi = 2\pi\ \mathrm{m\,s^{-1}}$. \checkmark
\end{exoresolu}

\courssub{Gears in Contact}

Two gear wheels meshing directly (no belt) behave similarly, with one essential difference: at the point of contact the teeth must not slide over each other, so the \cle{tangential speeds} of the two rims at the contact point are equal in magnitude:
\[ \omega_1 r_1 = \omega_2 r_2 \qquad \text{(gears in contact)}. \]
However, unlike belt-driven pulleys which rotate in the \emph{same} sense (for an open belt), meshing gears rotate in \emph{opposite} senses. Also, since the numbers of teeth are proportional to the radii, one often writes $\omega_1 N_1 = \omega_2 N_2$ where $N_1, N_2$ are the numbers of teeth.

\begin{exemplebox}[Bicycle gears]
On a bicycle, the chain plays the role of the belt: the front chainring and the rear sprocket have the same chain (linear) speed. A large front ring with a small rear sprocket gives a high rear-wheel angular speed — the ``high gear'' used on the flat roads of the Ndop plain.
\end{exemplebox}

\courssub{Fairground Rides: the Rotor (``Wall of Death'')}

\textbf{Intuition.} In the rotor ride, people stand against the inner wall of a large spinning cylinder; the floor is then lowered, yet nobody falls. What holds them up? The wall pushes each person \emph{horizontally inwards} — this normal reaction $R$ is the centripetal force. Friction between the person and the wall then balances the weight.

\textbf{Analysis.} For a person of mass $m$ in a cylinder of radius $r$ spinning at angular speed $\omega$:
\begin{itemize}
\item Horizontally (towards the centre): the normal reaction $R$ provides the centripetal force: $R = mr\omega^{2}$.
\item Vertically (no falling): friction $f$ acts upwards, balancing weight: $f = mg$. The maximum available friction is $f_{\max} = \mu R$, where $\mu$ is the coefficient of static friction between clothing and wall.
\end{itemize}
For the person not to slip down, we need $mg \le \mu R = \mu\, mr\omega^{2}$, giving the condition
\[ \omega \ge \sqrt{\frac{g}{\mu r}}. \]
Note the inequality direction: the ride must spin \emph{fast enough}. The faster it spins, the larger $R$, the larger the available friction.

\begin{exoresolu}[Minimum spin rate of a rotor]
A rotor has radius $2.5\ \mathrm{m}$ and the coefficient of static friction between clothing and wall is $0.4$. Taking $g = 10\ \mathrm{m\,s^{-2}}$, find the minimum angular speed for riders to stay pinned when the floor drops.
\tcblower
\textbf{Solution.} The limiting condition is $\mu R = mg$ with $R = mr\omega^{2}$:
\[ \omega_{\min} = \sqrt{\frac{g}{\mu r}} = \sqrt{\frac{10}{0.4 \times 2.5}} = \sqrt{10} \approx 3.16\ \mathrm{rad\,s^{-1}}. \]
In revolutions per minute: $N = \dfrac{60\omega}{2\pi} = \dfrac{60 \times 3.16}{2\pi} \approx 30.2\ \mathrm{rpm}$.
\end{exoresolu}

\courssub{Objects on a Rotating Turntable}

\textbf{Intuition.} Place a coin on a record player turntable. Near the centre it rotates with the disc; far from the centre it flies off. Why? The coin moves in a circle because \cle{static friction} with the disc pulls it towards the centre. The required centripetal force $mr\omega^{2}$ grows with $r$ and with $\omega^{2}$, but friction is capped at $\mu mg$.

\textbf{Condition.} No slipping requires
\[ mr\omega^{2} \le \mu mg \quad\Longleftrightarrow\quad \omega \le \sqrt{\frac{\mu g}{r}}. \]
Equivalently, at a given $\omega$, the coin stays in place only if $r \le \dfrac{\mu g}{\omega^{2}}$.

\begin{attention}[Static friction, and the direction of the inequality]
In turntable problems the friction is \emph{static} (the coin does not slide on the disc while rotating with it); slipping \emph{starts} exactly when $mr\omega^{2}$ exceeds $\mu mg$. Do not confuse this with the rotor, where a \emph{minimum} $\omega$ is required: on a turntable there is a \emph{maximum} $\omega$. Also note that $\omega_{\max}$ decreases as $r$ increases — objects near the rim slip first.
\end{attention}

\begin{exoresolu}[Coin on a turntable]
A coin is placed $8\ \mathrm{cm}$ from the centre of a turntable; the coefficient of static friction is $0.5$ and $g = 10\ \mathrm{m\,s^{-2}}$. (a) Find the greatest angular speed for which the coin does not slip. (b) The turntable spins at $60\ \mathrm{rpm}$; find the greatest distance from the centre at which the coin can sit without slipping.
\tcblower
\textbf{Solution.} (a) Limiting condition $mr\omega^{2} = \mu mg$:
\[ \omega_{\max} = \sqrt{\frac{\mu g}{r}} = \sqrt{\frac{0.5 \times 10}{0.08}} = \sqrt{62.5} \approx 7.91\ \mathrm{rad\,s^{-1}}. \]
(b) $\omega = \dfrac{2\pi \times 60}{60} = 2\pi\ \mathrm{rad\,s^{-1}}$. Then
\[ r_{\max} = \frac{\mu g}{\omega^{2}} = \frac{0.5 \times 10}{(2\pi)^{2}} = \frac{5}{4\pi^{2}} \approx 0.127\ \mathrm{m} = 12.7\ \mathrm{cm}. \]
\end{exoresolu}

\courssub{Satellites in Circular Orbit (Extension)}

\begin{savaistu}[Orbital motion — the same method]
For a satellite of mass $m$ in a circular orbit of radius $r$ around the Earth, the centripetal force is supplied by \cle{gravity}: $F = \dfrac{GMm}{r^{2}}$, where $M$ is the Earth's mass. Equating with the centripetal requirement,
\[ \frac{GMm}{r^{2}} = \frac{mv^{2}}{r} \quad\Longrightarrow\quad v = \sqrt{\frac{GM}{r}}. \]
The satellite's mass cancels: orbital speed depends only on the orbit radius. Larger orbits mean \emph{slower} satellites. This is why communication satellites that must hover over one point of the equator (geostationary orbit, period $24\ \mathrm{h}$) sit at a very precise, large radius ($\approx 42\,000\ \mathrm{km}$ from Earth's centre). The detailed study belongs to the Gravitation chapter, but the structure — \emph{identify the force provider, equate to $mv^{2}/r$, solve} — is exactly the method of this chapter.
\end{savaistu}

\courssub{Bonus Extension: Banking of Roads}

\begin{savaistu}[Why are some curves tilted?]
On fast highways and velodromes, curves are \emph{banked}: the outer edge is raised. The normal reaction then has a horizontal component pointing towards the centre, which helps (or even replaces) friction. For a vehicle taking a banked curve at the ``design speed'' with \emph{no} friction needed, resolving forces gives
\[ \tan\theta = \frac{v^{2}}{rg}, \]
where $\theta$ is the banking angle. Derivation: vertically $R\cos\theta = mg$, horizontally $R\sin\theta = \dfrac{mv^{2}}{r}$; divide the two equations. This result appears in some GCE papers as an extension — treat it as a bonus. The core syllabus result remains $v_{\max} = \sqrt{\mu r g}$ for flat curves.
\end{savaistu}

\courssub{Consolidation Exercises}

\begin{exercice}[Mixed practice]
\begin{enumerate}
\item A car of mass $900\ \mathrm{kg}$ rounds a flat bend of radius $40\ \mathrm{m}$ at $15\ \mathrm{m\,s^{-1}}$. (a) Calculate the centripetal force required. (b) If $\mu = 0.6$ and $g = 10\ \mathrm{m\,s^{-2}}$, does the car skid?
\item A belt connects a pulley of radius $6\ \mathrm{cm}$ rotating at $900\ \mathrm{rpm}$ to a pulley of radius $15\ \mathrm{cm}$. Find the angular speed of the second pulley in rpm and in $\mathrm{rad\,s^{-1}}$, and the speed of the belt.
\item A small block rests $20\ \mathrm{cm}$ from the axis of a horizontal turntable with $\mu = 0.4$. Find the maximum rotation rate in rpm for which the block does not slip (take $g = 10\ \mathrm{m\,s^{-2}}$).
\item In a rotor of radius $3\ \mathrm{m}$, what minimum angular speed pins a rider to the wall if $\mu = 0.5$? Express your answer in $\mathrm{rad\,s^{-1}}$ and rpm.
\end{enumerate}
\end{exercice}

\begin{corrige}[Mixed practice]
\textbf{1.} (a) $F = \dfrac{mv^{2}}{r} = \dfrac{900 \times 15^{2}}{40} = 5062.5\ \mathrm{N}$. (b) Maximum friction available: $\mu mg = 0.6 \times 900 \times 10 = 5400\ \mathrm{N} > 5062.5\ \mathrm{N}$, so the car does \emph{not} skid (but it is close to the limit; indeed $v_{\max} = \sqrt{0.6 \times 40 \times 10} = \sqrt{240} \approx 15.5\ \mathrm{m\,s^{-1}}$).

\textbf{2.} Equal belt speeds: $\omega_2 = \omega_1 \dfrac{r_1}{r_2} = 900 \times \dfrac{6}{15} = 360\ \mathrm{rpm}$, i.e. $\omega_2 = \dfrac{2\pi \times 360}{60} = 12\pi \approx 37.7\ \mathrm{rad\,s^{-1}}$. Belt speed: $v = r_2\omega_2 = 0.15 \times 12\pi = 1.8\pi \approx 5.65\ \mathrm{m\,s^{-1}}$.

\textbf{3.} $\omega_{\max} = \sqrt{\dfrac{\mu g}{r}} = \sqrt{\dfrac{0.4 \times 10}{0.20}} = \sqrt{20} \approx 4.47\ \mathrm{rad\,s^{-1}}$, i.e. $N = \dfrac{60 \times 4.47}{2\pi} \approx 42.7\ \mathrm{rpm}$.

\textbf{4.} $\omega_{\min} = \sqrt{\dfrac{g}{\mu r}} = \sqrt{\dfrac{10}{0.5 \times 3}} = \sqrt{\dfrac{20}{3}} \approx 2.58\ \mathrm{rad\,s^{-1}}$, i.e. about $24.6\ \mathrm{rpm}$.
\end{corrige}

\begin{aretenir}[Key points of this section]
\begin{itemize}
\item In every application, first \textbf{identify what provides the centripetal force}, then write $F = \dfrac{mv^{2}}{r} = mr\omega^{2}$ and impose the physical limit.
\item Flat curve: friction provides $F$; $v_{\max} = \sqrt{\mu r g}$ (mass cancels).
\item Belt/pulleys and gears: equal tangential speed, $\omega_1 r_1 = \omega_2 r_2$; gears rotate in opposite senses, open-belt pulleys in the same sense.
\item Rotor: normal reaction provides $F$; not slipping requires $\omega \ge \sqrt{g/(\mu r)}$.
\item Turntable: static friction provides $F$; no slipping requires $\omega \le \sqrt{\mu g / r}$.
\item Satellites: gravity provides $F$; $v = \sqrt{GM/r}$ (extension).
\item Master the conversion $\omega = \dfrac{2\pi N}{60}$ between rpm and $\mathrm{rad\,s^{-1}}$ — it appears in nearly every GCE question on this topic.
\end{itemize}
\end{aretenir}

\coursec{The Conical Pendulum: Concept and Forces}

In the previous section we applied the relations $v=r\omega$, $a=\dfrac{v^{2}}{r}=r\omega^{2}$ and $F=\dfrac{mv^{2}}{r}$ to vehicles on roundabouts, satellites and rotating machinery. In every case, the \emph{source} of the centripetal force was different: friction for a car, gravity for a satellite, tension for a stone on a string. We now study one of the most beautiful and most frequently examined situations in circular motion: the \cle{conical pendulum}. It is the perfect illustration of how a single force — the tension in a string — can be split into two components, one balancing the weight and the other supplying the centripetal force. Mastering it is essential, because the GCE Advanced Level examiners return to it again and again.

\courssub{What is a Conical Pendulum?}

\textbf{An observation to begin with.}\par
Take a bunch of keys tied to a cord, hold the free end of the cord above your head and set the keys moving so that they travel round and round in a horizontal circle. You will notice something remarkable: the cord is \emph{not} vertical, and it is \emph{not} horizontal either. It makes a fixed angle with the vertical, and as the keys go round, the cord sweeps out the surface of an upside-down \cle{cone}. The faster you whirl the keys, the wider the angle and the higher the keys rise. This simple toy, which any student in Yaoundé or Bamenda can build in a minute, contains all the physics of this section.

\begin{definition}[Conical pendulum]
A \cle{conical pendulum} consists of a small heavy body (the \cle{bob}) of mass $m$, attached to a fixed support by a light inextensible string of length $l$, and set in motion so that the bob describes a \textbf{horizontal circle with uniform speed} while the string sweeps out the surface of a cone whose apex is the support. The constant angle $\theta$ between the string and the downward vertical is called the \cle{semi-vertical angle} of the cone.
\end{definition}

\begin{attention}[Do not confuse with the simple pendulum]
A \emph{simple} pendulum swings back and forth in a \textbf{vertical plane}; its speed is not constant and the motion is not circular motion with uniform speed. A \emph{conical} pendulum moves in a \textbf{horizontal circle} at \textbf{constant speed}. The two devices look similar but obey completely different physics.
\end{attention}

\begin{savaistu}[Conical pendulums around us]
The conical pendulum is far more than a classroom curiosity. The \emph{swing rides} (``flying chairs'') at amusement fairs are giant conical pendulums: each chair hangs from a chain and rises as the carousel spins faster. Historically, James Watt used a pair of conical-pendulum balls — the \emph{centrifugal governor} — to regulate the speed of steam engines: when the engine ran too fast, the balls rose and closed a valve. The same principle is still used in some engine governors today.
\end{savaistu}

\courssub{Geometry of the Set-up}

Before writing any equation of motion, we must describe the geometry precisely. Let the string have length $l$ and make a constant angle $\theta$ with the vertical. The bob moves on a horizontal circle of radius $r$, at a constant vertical distance $h$ below the support.

\begin{popfigure}
\begin{center}
\begin{tikzpicture}[scale=1.15, line cap=round]
% support
\fill[poppurple] (0,0) circle (0.06);
\draw[popmuted] (-0.6,0) -- (0.6,0);
\foreach \x in {-0.5,-0.25,0,0.25,0.5}{\draw[popmuted] (\x,0) -- (\x+0.15,0.18);}
% cone (dashed)
\draw[dashed, popline] (0,0) -- (2.6,-3.0);
\draw[dashed, popline] (0,0) -- (-2.6,-3.0);
% circular path (ellipse)
\draw[dashed, popblue] (0,-3.0) ellipse (2.6 and 0.55);
% string to bob
\draw[very thick, popdark] (0,0) -- (2.6,-3.0);
% bob
\shade[ball color=poporange] (2.6,-3.0) circle (0.16);
\node[popdark, right] at (2.78,-3.0) {bob $m$};
% vertical axis
\draw[dashed, popmuted] (0,0) -- (0,-3.55);
% angle theta
\draw[popgreen, thick] (0,-1.1) arc (-90:-49.1:1.1);
\node[popgreen] at (0.42,-1.25) {$\theta$};
% radius r
\draw[<->, popblue, thick] (0,-3.62) -- (2.6,-3.62);
\node[popblue, below] at (1.3,-3.62) {$r$};
% height h
\draw[<->, poppink, thick] (-0.35,0) -- (-0.35,-3.0);
\node[poppink, left] at (-0.35,-1.5) {$h$};
% length l
\node[popdark, rotate=-49.1] at (1.05,-1.28) {$l$};
% centre label
\fill[popblue] (0,-3.0) circle (0.05);
\node[popblue, left] at (-0.08,-3.0) {$O$};
% velocity arrow (tangent)
\draw[->, very thick, popgreen] (2.6,-3.0) -- (2.6,-2.0);
\node[popgreen, right] at (2.6,-1.9) {$\vec{v}$};
\end{tikzpicture}
\end{center}
\legende{A conical pendulum: the string of length $l$ sweeps out a cone of semi-vertical angle $\theta$; the bob describes a horizontal circle of radius $r$ at depth $h$ below the support.}
\end{popfigure}

From the right-angled triangle formed by the support, the centre $O$ of the circle and the bob, we read off immediately:

\begin{propriete}[Geometry of the conical pendulum]
If the string has length $l$ and makes an angle $\theta$ with the vertical, then the radius of the circle and the depth of the bob below the support are
\[ r = l\sin\theta \qquad\text{and}\qquad h = l\cos\theta. \]
\end{propriete}

\begin{attention}[The radius is $r = l\sin\theta$, not $l$]
This is the single most common error at the GCE. The bob does \emph{not} move on a circle of radius $l$: it moves on a \textbf{horizontal} circle whose radius is only $l\sin\theta$. Whenever you write $a=\dfrac{v^{2}}{r}$ or $F=mr\omega^{2}$, the $r$ must be $l\sin\theta$. Writing $l$ instead of $l\sin\theta$ destroys the whole problem.
\end{attention}

\courssub{Forces Acting on the Bob: the Free-Body Diagram}

\textbf{How many forces act on the bob?}\par
Students often want to add a third force — an ``outward'' or ``centrifugal'' force, or a separate ``centripetal force'' drawn alongside the others. Both are wrong. Look at the actual situation: the bob is in contact with nothing but the string, and it is attracted by the Earth. That is all.

\begin{propriete}[The two forces on the bob]
At every instant, exactly \textbf{two} forces act on the bob:
\begin{itemize}
\item its \cle{weight} $\vec{W}=m\vec{g}$, acting vertically downwards;
\item the \cle{tension} $\vec{T}$ of the string, acting along the string towards the support.
\end{itemize}
There is no additional ``centripetal force'': the centripetal force is not a new force, it is the \emph{role played} by the horizontal component of the tension.
\end{propriete}

\begin{popfigure}
\begin{center}
\begin{tikzpicture}[scale=1.3, line cap=round]
% bob
\shade[ball color=poporange] (0,0) circle (0.18);
% tension along string (up-left)
\draw[->, very thick, popblue] (0,0) -- (-1.9,1.55);
\node[popblue, above left] at (-1.55,1.25) {$\vec{T}$};
% string direction dashed continuation
\draw[dashed, popmuted] (0,0) -- (-2.6,2.12);
% angle
\draw[popgreen, thick] (-0.75,0.61) arc (140.8:180:0.97);
\node[popgreen] at (-1.15,0.32) {$\theta$};
% vertical component
\draw[->, very thick, popgreen] (0,0) -- (0,1.55);
\node[popgreen, right] at (0.05,1.15) {$T\cos\theta$};
% horizontal component
\draw[->, very thick, poppink] (0,0) -- (-1.9,0);
\node[poppink, below] at (-1.0,-0.06) {$T\sin\theta$ (towards centre)};
% weight
\draw[->, very thick, popdark] (0,0) -- (0,-1.55);
\node[popdark, right] at (0.05,-1.2) {$m\vec{g}$};
% construction dashed
\draw[dashed, popmuted] (-1.9,0) -- (-1.9,1.55);
\draw[dashed, popmuted] (0,1.55) -- (-1.9,1.55);
% centre indication
\node[popmuted, align=center, text width=3.2cm] at (1.9,0.4) {centre $O$ of the circle is to the left};
\end{tikzpicture}
\end{center}
\legende{Free-body diagram of the bob: the tension is resolved into a vertical component $T\cos\theta$ (balancing the weight) and a horizontal component $T\sin\theta$ (the centripetal force).}
\end{popfigure}

\courssub{Resolving the Tension and Writing the Equations of Motion}

Since the bob moves in a \emph{horizontal} circle, its acceleration is \emph{horizontal}, directed towards the centre $O$, of magnitude $\dfrac{v^{2}}{r}=r\omega^{2}$. Vertically, the bob neither rises nor falls: its vertical acceleration is zero. We resolve the tension $\vec{T}$ along these two natural directions:

\begin{itemize}
\item \textbf{vertical component:} $T\cos\theta$ (upwards);
\item \textbf{horizontal component:} $T\sin\theta$ (towards the centre).
\end{itemize}

Applying Newton's second law in each direction:

\textbf{Vertically} (no acceleration, forces balance):
\[ T\cos\theta = mg. \]

\textbf{Horizontally} (the resultant provides the centripetal force):
\[ T\sin\theta = \frac{mv^{2}}{r} = mr\omega^{2}. \]

\begin{propriete}[Equations of the conical pendulum — examinable]
For a conical pendulum of semi-vertical angle $\theta$, with a bob of mass $m$ moving at speed $v$ on a circle of radius $r=l\sin\theta$:
\[ T\cos\theta = mg \qquad\text{and}\qquad T\sin\theta = \frac{mv^{2}}{r}=mr\omega^{2}. \]
Dividing the second equation by the first eliminates $T$ and $m$, giving the \textbf{key relation}
\[ \tan\theta = \frac{v^{2}}{rg} = \frac{r\omega^{2}}{g}. \]
This derivation is examinable: you must be able to produce it from a clear free-body diagram and Newton's second law.
\end{propriete}

\begin{methode}[Solving any conical pendulum problem]
\begin{enumerate}
\item Draw the free-body diagram: \textbf{only} the weight $mg$ (down) and the tension $T$ (along the string).
\item Write the geometry: $r=l\sin\theta$.
\item Resolve the tension: vertical component $T\cos\theta$, horizontal component $T\sin\theta$.
\item Write $T\cos\theta = mg$ (vertical equilibrium) and $T\sin\theta = \dfrac{mv^{2}}{r}$ (centripetal force).
\item Divide to get $\tan\theta=\dfrac{v^{2}}{rg}$, or solve the system for the unknown requested ($T$, $v$, $\theta$, $\omega$).
\end{enumerate}
\end{methode}

\begin{attention}[The bob is NOT in equilibrium]
The vertical forces balance, so there is no \emph{vertical} acceleration — but the bob is certainly not in equilibrium! The horizontal resultant $T\sin\theta$ is non-zero and permanently pulls the bob towards the centre: this is precisely the centripetal force that bends the path into a circle. If all forces balanced, the bob would move in a straight line (Newton's first law). Likewise, never draw an outward ``centrifugal force'': in the inertial frame of the ground, no such force exists.
\end{attention}

\begin{exoresolu}[A bunch of keys as a conical pendulum]
A student ties a bunch of keys of mass $m=0.20\ \mathrm{kg}$ to a cord of length $l=0.80\ \mathrm{m}$ and whirls it so that the keys describe a horizontal circle, the cord making an angle $\theta=30^{\circ}$ with the vertical. Take $g=9.81\ \mathrm{m\,s^{-2}}$.
\begin{enumerate}
\item Find the tension in the cord.
\item Find the speed of the keys.
\item Find the angular speed of the keys.
\end{enumerate}
\tcblower
\textbf{Solution.} The free-body diagram gives the two standard equations, with $r=l\sin\theta$.
\[ r = l\sin\theta = 0.80\sin 30^{\circ} = 0.80\times 0.5 = 0.40\ \mathrm{m}. \]
\textbf{1. Tension.} From the vertical equation:
\[ T\cos\theta = mg \;\Longrightarrow\; T = \frac{mg}{\cos\theta} = \frac{0.20\times 9.81}{\cos 30^{\circ}} = \frac{1.962}{0.8660} \approx 2.27\ \mathrm{N}. \]
\textbf{2. Speed.} From the key relation:
\[ \tan\theta = \frac{v^{2}}{rg} \;\Longrightarrow\; v^{2} = rg\tan\theta = 0.40\times 9.81\times \tan 30^{\circ} = 0.40\times 9.81\times 0.5774 \approx 2.266, \]
\[ v \approx 1.51\ \mathrm{m\,s^{-1}}. \]
\textbf{3. Angular speed.} Using $v=r\omega$:
\[ \omega = \frac{v}{r} = \frac{1.51}{0.40} \approx 3.77\ \mathrm{rad\,s^{-1}}. \]
\emph{Check:} $T\sin\theta = 2.27\times 0.5 = 1.13\ \mathrm{N}$ and $\dfrac{mv^{2}}{r}=\dfrac{0.20\times 2.266}{0.40}=1.13\ \mathrm{N}$. \checkmark\ The horizontal equation is satisfied.
\end{exoresolu}

\begin{exemplebox}[Reading the physics from the key relation]
The relation $\tan\theta=\dfrac{r\omega^{2}}{g}$ tells us something physically important: \textbf{the faster the bob rotates, the larger the angle $\theta$}. Indeed, with $r=l\sin\theta$, the relation becomes
\[ \tan\theta = \frac{l\sin\theta\,\omega^{2}}{g} \;\Longrightarrow\; \cos\theta = \frac{g}{l\omega^{2}}, \]
so as $\omega$ increases, $\cos\theta$ decreases and $\theta$ grows towards $90^{\circ}$. This is exactly why the chairs of a swing ride rise when the carousel speeds up, and why the balls of Watt's governor lifted when the engine ran too fast. Note also that $\cos\theta\le 1$ requires $\omega^{2}\ge \dfrac{g}{l}$: below this angular speed, no conical motion is possible and the string simply hangs vertically.
\end{exemplebox}

\begin{exercice}[Practice]
A small ball of mass $0.15\ \mathrm{kg}$ is attached to a string of length $1.2\ \mathrm{m}$ fixed at its upper end. The ball moves in a horizontal circle with angular speed $\omega = 3.0\ \mathrm{rad\,s^{-1}}$. Take $g = 9.81\ \mathrm{m\,s^{-2}}$.
\begin{enumerate}
\item Show that $\cos\theta = \dfrac{g}{l\omega^{2}}$ and deduce the semi-vertical angle $\theta$.
\item Calculate the radius of the circle and the tension in the string.
\item Calculate the linear speed of the ball.
\end{enumerate}
\end{exercice}

\begin{corrige}[Exercise 1]
\textbf{1.} From $T\cos\theta=mg$ and $T\sin\theta = mr\omega^{2}$ with $r=l\sin\theta$:
\[ \frac{T\sin\theta}{T\cos\theta}=\tan\theta=\frac{l\sin\theta\,\omega^{2}}{g}\;\Longrightarrow\;\cos\theta=\frac{g}{l\omega^{2}}=\frac{9.81}{1.2\times 9.0}=0.908, \]
so $\theta \approx 24.7^{\circ}$.
\textbf{2.} $r=l\sin\theta = 1.2\times\sin 24.7^{\circ}\approx 0.50\ \mathrm{m}$; \; $T=\dfrac{mg}{\cos\theta}=\dfrac{0.15\times 9.81}{0.908}\approx 1.62\ \mathrm{N}$.
\textbf{3.} $v=r\omega = 0.50\times 3.0 = 1.50\ \mathrm{m\,s^{-1}}$.
\end{corrige}

\begin{aretenir}[Key points of this section]
\begin{itemize}
\item A conical pendulum: a bob describing a \textbf{horizontal circle at constant speed} while the string sweeps a cone of semi-vertical angle $\theta$.
\item Geometry: $r=l\sin\theta$, $h=l\cos\theta$ — never confuse $r$ with $l$.
\item Only \textbf{two} forces act on the bob: weight $mg$ and tension $T$.
\item Resolve the tension: $T\cos\theta=mg$ (vertical balance) and $T\sin\theta=\dfrac{mv^{2}}{r}=mr\omega^{2}$ (centripetal force).
\item Key relation: $\tan\theta=\dfrac{v^{2}}{rg}=\dfrac{r\omega^{2}}{g}$, equivalently $\cos\theta=\dfrac{g}{l\omega^{2}}$.
\item The bob is \textbf{not} in equilibrium: it accelerates towards the centre. There is no outward force.
\end{itemize}
\end{aretenir}

Now that the forces on the conical pendulum are fully understood, a natural question arises: how long does the bob take to complete one revolution? In the next section we use the equations established here to \emph{derive the formula for the period} of the conical pendulum and to compare it with the period of the simple pendulum — a classic GCE examination question.

\coursec{Period of the Conical Pendulum: Derivation and Applications}

In the previous section we analysed the forces acting on a conical pendulum and established the two fundamental equations of motion: the vertical equilibrium $T\cos\theta = mg$ and the horizontal centripetal equation $T\sin\theta = m\omega^{2}r$. We now harvest the fruit of that work. By combining these equations we shall derive one of the most elegant results of circular motion: the period of a conical pendulum depends \emph{only} on the vertical height of the support above the plane of the circle. This result is a favourite of GCE examiners, both for its derivation and for its applications.

\courssub{From the Equations of Motion to the Angular Speed}

Recall the situation: a bob of mass $m$ is attached to a light inextensible string of length $l$ fixed at a point $S$. The bob describes a horizontal circle of radius $r$ with constant angular speed $\omega$, the string making a constant angle $\theta$ with the vertical.

\begin{popfigure}
\begin{center}
\begin{tikzpicture}[scale=2.2]
\coordinate (S) at (0,0);
\coordinate (B) at (1.3,-1.5);
\coordinate (C) at (0,-1.5);
% ceiling
\draw[thick] (-1.2,0) -- (1.2,0);
\foreach \x in {-1.1,-0.9,...,1.1}{\draw (\x,0) -- (\x+0.12,0.15);}
% string and bob
\draw[thick,popblue] (S) -- (B);
\fill[poporange] (B) circle (0.07);
\node[below right,popdark] at (B) {bob, mass $m$};
% circle of motion (ellipse)
\draw[dashed,popmuted] (C) ellipse (1.3 and 0.28);
% vertical line and h
\draw[dashed,popdark] (S) -- (C);
\draw[<->,popgreen,thick] (-0.35,0) -- (-0.35,-1.5) node[midway,left]{$h = l\cos\theta$};
% radius
\draw[<->,poppurple,thick] (C) -- (B) node[midway,below]{$r = l\sin\theta$};
% angle
\draw[poporange,thick] (0,-0.55) arc (-90:-60:0.55);
\node[poporange] at (0.28,-0.62) {$\theta$};
% length label
\node[popblue,above left] at (0.5,-0.85) {$l$};
% angular speed arrow
\draw[->,popgreen,thick] (1.55,-1.35) arc (20:-40:0.35);
\node[popgreen,right] at (1.85,-1.5) {$\omega$};
\node[above,popdark] at (S) {$S$};
\end{tikzpicture}
\end{center}
\legende{The geometry of the conical pendulum: the vertical distance from the support $S$ to the plane of the circle is $h = l\cos\theta$, and the radius of the circle is $r = l\sin\theta$.}
\end{popfigure}

Dividing the horizontal equation by the vertical equation eliminates the tension:
\[
\frac{T\sin\theta}{T\cos\theta} = \frac{m\omega^{2}r}{mg}
\quad\Longrightarrow\quad
\tan\theta = \frac{\omega^{2}r}{g}.
\]
Now substitute the geometric relation $r = l\sin\theta$:
\[
\frac{\sin\theta}{\cos\theta} = \frac{\omega^{2}l\sin\theta}{g}.
\]
Since $\theta \neq 0$ during the motion (the bob genuinely moves in a circle), we have $\sin\theta \neq 0$ and we may divide both sides by $\sin\theta$:
\[
\frac{1}{\cos\theta} = \frac{\omega^{2}l}{g}
\quad\Longrightarrow\quad
\omega^{2} = \frac{g}{l\cos\theta}.
\]
But from the geometry of the figure, $l\cos\theta = h$, the vertical distance from the support to the plane of the circle. Hence
\[
\boxed{\ \omega^{2} = \frac{g}{l\cos\theta} = \frac{g}{h}\ }.
\]

\begin{attention}[Why we may divide by $\sin\theta$]
The step ``divide by $\sin\theta$'' is legitimate only because the pendulum is actually revolving, so $\theta > 0$. If $\theta = 0$ the bob simply hangs at rest and there is no circular motion to analyse. Mentioning this in an examination shows the examiner that you understand the mathematics, not just the algebra.
\end{attention}

\courssub{The Period of the Conical Pendulum}

\begin{propriete}[Period of a conical pendulum (proof examinable)]
A conical pendulum whose bob moves in a horizontal circle at a vertical distance $h$ below the support has period
\[
T = \frac{2\pi}{\omega} = 2\pi\sqrt{\frac{l\cos\theta}{g}} = 2\pi\sqrt{\frac{h}{g}},
\]
where $l$ is the length of the string, $\theta$ the angle the string makes with the vertical, and $h = l\cos\theta$. The derivation of this result is examinable at the GCE Advanced Level.
\end{propriete}

\textbf{Proof.} Since the bob completes one revolution ($2\pi$ radians) in one period $T$, the angular speed is $\omega = \dfrac{2\pi}{T}$. Substituting into $\omega^{2} = \dfrac{g}{h}$:
\[
\frac{4\pi^{2}}{T^{2}} = \frac{g}{h}
\quad\Longrightarrow\quad
T^{2} = \frac{4\pi^{2}h}{g}
\quad\Longrightarrow\quad
T = 2\pi\sqrt{\frac{h}{g}}.
\]
Since $h = l\cos\theta$, the equivalent form $T = 2\pi\sqrt{\dfrac{l\cos\theta}{g}}$ follows immediately. \hfill$\blacksquare$

\begin{aretenir}[The remarkable feature of the formula]
The period $T = 2\pi\sqrt{h/g}$ depends \emph{only} on:
\begin{itemize}
\item the vertical height $h$ of the support above the plane of the circle;
\item the acceleration due to gravity $g$.
\end{itemize}
It does \textbf{not} depend on the mass $m$ of the bob, nor directly on the length $l$ or the angle $\theta$ separately. Two conical pendulums of different lengths, different masses and different angles have exactly the same period provided their bobs revolve in the same horizontal plane (same $h$).
\end{aretenir}

\begin{methode}[Derivation checklist — learn each step]
\begin{enumerate}
\item Draw the force diagram: weight $mg$ vertically down, tension $T$ along the string.
\item Resolve vertically (no vertical acceleration): $T\cos\theta = mg$.
\item Resolve horizontally towards the centre (Newton's second law, centripetal force): $T\sin\theta = m\omega^{2}r$.
\item Divide the two equations to eliminate $T$: $\tan\theta = \dfrac{\omega^{2}r}{g}$.
\item Substitute the geometry $r = l\sin\theta$ and simplify: $\omega^{2} = \dfrac{g}{l\cos\theta} = \dfrac{g}{h}$.
\item Use $\omega = \dfrac{2\pi}{T}$ and solve for the period: $T = 2\pi\sqrt{\dfrac{h}{g}}$.
\end{enumerate}
\end{methode}

\begin{attention}[Period $T$ versus tension $T$ — a classic trap]
The same letter $T$ is traditionally used for both the \emph{period} (a time, in seconds) and the \emph{tension} (a force, in newtons). In your examination script, never write ``$T = 2\pi\sqrt{h/g}$'' and ``$T = mg/\cos\theta$'' in the same piece of work without comment — the examiner cannot tell which is which. Use $\tau$ (or write ``period'' and ``tension'' in words) to distinguish them. Confusing the two costs marks every year.
\end{attention}

\courssub{The Tension in the String}

\begin{propriete}[Tension in a conical pendulum]
The tension in the string of a conical pendulum is
\[
T = \frac{mg}{\cos\theta} = m\omega^{2}l.
\]
Consequently, the tension increases as $\theta$ increases (equivalently, as the rotation gets faster).
\end{propriete}

\textbf{Justification.} From vertical equilibrium, $T\cos\theta = mg$, so $T = \dfrac{mg}{\cos\theta}$. For the second form, the horizontal equation gives $T\sin\theta = m\omega^{2}r = m\omega^{2}l\sin\theta$, and dividing by $\sin\theta \neq 0$ yields $T = m\omega^{2}l$. The two expressions are consistent: substituting $\omega^{2} = \dfrac{g}{l\cos\theta}$ into $m\omega^{2}l$ gives $\dfrac{mg}{\cos\theta}$.

Since $\cos\theta$ \emph{decreases} as $\theta$ grows from $0$ to $90^{\circ}$, the tension $\dfrac{mg}{\cos\theta}$ \emph{increases} without bound. This is why the string of a fast-spinning toy or the chains of a swing carousel are under great stress.

\courssub{Limiting Cases}

\textbf{Case $\theta \to 0$.} Then $\cos\theta \to 1$, so $h \to l$ and
\[
T \to 2\pi\sqrt{\frac{l}{g}},
\]
which is exactly the period of a \emph{simple pendulum} of length $l$ oscillating with small amplitude. This is physically reasonable: when the cone is very narrow, the bob barely leaves the vertical and the motion resembles a gentle swing. (This is a qualitative remark; the simple pendulum formula itself is established properly in the study of simple harmonic motion.)

\textbf{Case $\theta \to 90^{\circ}$.} Then $\cos\theta \to 0$, so the tension $T = \dfrac{mg}{\cos\theta} \to \infty$ and the angular speed $\omega = \sqrt{\dfrac{g}{l\cos\theta}} \to \infty$. A perfectly horizontal string is therefore \emph{impossible in practice}: no finite tension can make the string horizontal, because some vertical component of tension is always needed to balance the weight.

\begin{savaistu}[Why the swing carousel tilts outwards]
At amusement parks, the chairs of a swing carousel rise and tilt outwards as the ride speeds up. Each chair and rider behaves like a conical pendulum: from $\tan\theta = \dfrac{\omega^{2}r}{g}$, a larger $\omega$ requires a larger $\theta$. The chains can never become horizontal, however fast the ride spins — exactly as the theory predicts.
\end{savaistu}

\courssub{Graphical Interpretation: $T$ against $\sqrt{h}$}

Since $T = \dfrac{2\pi}{\sqrt{g}}\sqrt{h}$, a graph of the period $T$ against $\sqrt{h}$ is a straight line through the origin with gradient $\dfrac{2\pi}{\sqrt{g}}$. Equivalently, $T^{2} = \dfrac{4\pi^{2}}{g}h$, so a graph of $T^{2}$ against $h$ is a straight line through the origin of gradient $\dfrac{4\pi^{2}}{g}$.

\begin{popfigure}
\begin{center}
\begin{tikzpicture}
\begin{axis}[
    width=0.75\linewidth, height=6.5cm,
    xlabel={$\sqrt{h}\;(\mathrm{m^{1/2}})$},
    ylabel={$T\;(\mathrm{s})$},
    xmin=0, xmax=1.15, ymin=0, ymax=2.4,
    xtick={0,0.2,0.4,0.6,0.8,1.0},
    ytick={0,0.5,1.0,1.5,2.0},
    grid=both, grid style={popline, very thin},
    axis lines=left, thick,
]
\addplot[popblue, very thick, domain=0:1.1, samples=2] {2.006*x};
\addplot[only marks, mark=*, poporange, mark size=2pt] coordinates {(0.316,0.634) (0.548,1.099) (0.707,1.418) (0.837,1.678) (1.0,2.006)};
\node[popblue, anchor=south west, font=\small] at (axis cs:0.5,1.35) {gradient $= \dfrac{2\pi}{\sqrt{g}}$};
\end{axis}
\end{tikzpicture}
\end{center}
\legende{Period $T$ plotted against $\sqrt{h}$ for a conical pendulum: a straight line through the origin of gradient $2\pi/\sqrt{g} \approx 2.006\ \mathrm{s\,m^{-1/2}}$ (taking $g = 9.81\ \mathrm{m\,s^{-2}}$).}
\end{popfigure}

\begin{experience}[Experimental determination of $g$ with a conical pendulum]
This provides an excellent practical method of measuring $g$, in the spirit of GCE practical-style questions.
\begin{enumerate}
\item Suspend a bob on a string of measured length $l$ and set it revolving in a horizontal circle (a small horizontal push starts the motion).
\item Time $n = 20$ complete revolutions with a stopwatch; the period is $T = t/n$. Repeat and average.
\item Measure the vertical height $h$ from the support to the plane of the circle (or compute $h = l\cos\theta$ from a measured $\theta$).
\item Repeat for several values of $h$ (vary $l$ or the rotation speed).
\item Plot $T^{2}$ against $h$: since $T^{2} = \dfrac{4\pi^{2}}{g}h$, the graph is a straight line through the origin of gradient $\dfrac{4\pi^{2}}{g}$. Hence $g = \dfrac{4\pi^{2}}{\text{gradient}}$.
\end{enumerate}
Precautions: keep the circle truly horizontal, time many revolutions to reduce reaction-time error, and ensure the angle $\theta$ stays constant during timing.
\end{experience}

\courssub{Applications and Worked Examples}

The formulas derived above answer all the standard questions:
\begin{itemize}
\item \textbf{Finding $\theta$ given the rotation rate:} from $\cos\theta = \dfrac{g}{\omega^{2}l}$, compute $\theta$.
\item \textbf{Finding the speed of the bob:} $v = \omega r = \omega l\sin\theta$.
\item \textbf{Finding the tension:} $T = \dfrac{mg}{\cos\theta} = m\omega^{2}l$.
\item \textbf{Finding the period:} $T = 2\pi\sqrt{\dfrac{h}{g}}$.
\end{itemize}

\begin{exoresolu}[GCE-style worked example]
A small bob of mass $0.200\ \mathrm{kg}$ is attached to a light inextensible string of length $l = 1.5\ \mathrm{m}$ fixed at a point $S$. The bob moves in a horizontal circle with constant speed, the string making an angle $\theta = 30^{\circ}$ with the vertical. Taking $g = 9.81\ \mathrm{m\,s^{-2}}$, find:
\begin{enumerate}
\item the tension in the string;
\item the angular speed $\omega$ of the bob;
\item the linear speed $v$ of the bob;
\item the period of revolution.
\end{enumerate}
\tcblower
\textbf{Solution.}

\textbf{1. Tension.} Vertical equilibrium gives $T\cos\theta = mg$, so
\[
T = \frac{mg}{\cos\theta} = \frac{0.200 \times 9.81}{\cos 30^{\circ}} = \frac{1.962}{0.8660} \approx 2.27\ \mathrm{N}.
\]

\textbf{2. Angular speed.} Using $\omega^{2} = \dfrac{g}{l\cos\theta}$:
\[
\omega^{2} = \frac{9.81}{1.5 \times \cos 30^{\circ}} = \frac{9.81}{1.299} = 7.552
\quad\Longrightarrow\quad
\omega \approx 2.75\ \mathrm{rad\,s^{-1}}.
\]

\textbf{3. Linear speed.} The radius of the circle is $r = l\sin\theta = 1.5 \times \sin 30^{\circ} = 0.75\ \mathrm{m}$. Hence
\[
v = \omega r = 2.75 \times 0.75 \approx 2.06\ \mathrm{m\,s^{-1}}.
\]

\textbf{4. Period.} The height is $h = l\cos\theta = 1.5 \times 0.8660 = 1.299\ \mathrm{m}$, so
\[
T = 2\pi\sqrt{\frac{h}{g}} = 2\pi\sqrt{\frac{1.299}{9.81}} = 2\pi\sqrt{0.1324} \approx 2.29\ \mathrm{s}.
\]
\textbf{Check:} $T = \dfrac{2\pi}{\omega} = \dfrac{2\pi}{2.75} \approx 2.29\ \mathrm{s}$. Consistent. \hfill$\blacksquare$
\end{exoresolu}

\begin{exoresolu}[Finding the angle from the rotation rate]
A conical pendulum has a string of length $0.80\ \mathrm{m}$. The bob makes $60$ revolutions per minute. Find the angle the string makes with the vertical. (Take $g = 9.81\ \mathrm{m\,s^{-2}}$.)
\tcblower
\textbf{Solution.} Convert the rotation rate: $60\ \mathrm{rev\,min^{-1}} = 1\ \mathrm{rev\,s^{-1}}$, so
\[
\omega = 2\pi \times 1 = 2\pi\ \mathrm{rad\,s^{-1}}.
\]
From $\omega^{2} = \dfrac{g}{l\cos\theta}$:
\[
\cos\theta = \frac{g}{\omega^{2}l} = \frac{9.81}{(2\pi)^{2}\times 0.80} = \frac{9.81}{31.58} = 0.3106.
\]
Hence $\theta = \cos^{-1}(0.3106) \approx 71.9^{\circ}$. \hfill$\blacksquare$
\end{exoresolu}

\begin{attention}[Common mistakes in applications]
\begin{itemize}
\item Forgetting to convert revolutions per minute into $\mathrm{rad\,s^{-1}}$: multiply by $\dfrac{2\pi}{60}$.
\item Using $r = l$ instead of $r = l\sin\theta$ for the radius of the circle.
\item Using degrees instead of radians when applying $\omega = \dfrac{2\pi}{T}$ — angular speed formulas require radians.
\item Writing $\omega^{2} = \dfrac{g}{l\cos\theta}$ as $\omega = \dfrac{g}{l\cos\theta}$: do not drop the square root.
\item Quoting $T = 2\pi\sqrt{l/g}$ (the simple pendulum formula) instead of $T = 2\pi\sqrt{l\cos\theta/g}$: the factor $\cos\theta$ is essential.
\end{itemize}
\end{attention}

\begin{savaistu}[Exam tip — derivations earn full marks]
The derivation of $T = 2\pi\sqrt{h/g}$ is one of the most frequently set ``show that'' questions in the mechanics papers. Examiners award marks for \emph{each justified step}: the two resolutions, the division eliminating the tension, the substitution $r = l\sin\theta$, and the use of $\omega = 2\pi/T$. Reproduce the six-step checklist above with a clear force diagram, and state each physical principle you use (vertical equilibrium; Newton's second law towards the centre). A bare formula with no justification earns almost nothing.
\end{savaistu}

\begin{exercice}[Practice]
\begin{enumerate}
\item A conical pendulum consists of a bob of mass $0.5\ \mathrm{kg}$ on a string of length $2.0\ \mathrm{m}$. The bob revolves so that the string makes $40^{\circ}$ with the vertical. Taking $g = 9.81\ \mathrm{m\,s^{-2}}$, find the tension, the angular speed, the speed of the bob and the period.
\item Show that if a conical pendulum completes $n$ revolutions per second, then $\cos\theta = \dfrac{g}{4\pi^{2}n^{2}l}$. Hence find the least possible value of $n$ for a string of length $1.0\ \mathrm{m}$, and explain what limits the rotation rate in practice.
\item In an experiment, a conical pendulum gives the following results: $h = 0.25\ \mathrm{m}$, $T = 1.00\ \mathrm{s}$. Use $T = 2\pi\sqrt{h/g}$ to estimate $g$.
\item Two conical pendulums have strings of lengths $1.0\ \mathrm{m}$ and $4.0\ \mathrm{m}$. They revolve in the same horizontal plane, $0.50\ \mathrm{m}$ below their supports. Show that they have the same period, and find the angle each string makes with the vertical.
\end{enumerate}
\end{exercice}

\begin{corrige}[Exercise 1]
\textbf{1.} Tension: $T = \dfrac{mg}{\cos\theta} = \dfrac{0.5 \times 9.81}{\cos 40^{\circ}} = \dfrac{4.905}{0.766} \approx 6.40\ \mathrm{N}$.
Angular speed: $\omega = \sqrt{\dfrac{g}{l\cos\theta}} = \sqrt{\dfrac{9.81}{2.0 \times 0.766}} = \sqrt{6.403} \approx 2.53\ \mathrm{rad\,s^{-1}}$.
Speed: $r = l\sin\theta = 2.0\sin 40^{\circ} = 1.286\ \mathrm{m}$, so $v = \omega r = 2.53 \times 1.286 \approx 3.25\ \mathrm{m\,s^{-1}}$.
Period: $h = l\cos\theta = 1.532\ \mathrm{m}$, so $T = 2\pi\sqrt{\dfrac{1.532}{9.81}} \approx 2.48\ \mathrm{s}$.

\textbf{2.} With $\omega = 2\pi n$, the relation $\omega^{2} = \dfrac{g}{l\cos\theta}$ gives $4\pi^{2}n^{2} = \dfrac{g}{l\cos\theta}$, hence $\cos\theta = \dfrac{g}{4\pi^{2}n^{2}l}$. Since $\cos\theta \leq 1$, we need $n^{2} \geq \dfrac{g}{4\pi^{2}l}$, i.e.\ a \emph{minimum} rate $n_{\min} = \dfrac{1}{2\pi}\sqrt{\dfrac{g}{l}} = \dfrac{1}{2\pi}\sqrt{9.81} \approx 0.50\ \mathrm{rev\,s^{-1}}$ for $l = 1.0\ \mathrm{m}$ (below this, the bob simply hangs vertically). There is no theoretical maximum, but as $n$ increases, $\cos\theta \to 0$, $\theta \to 90^{\circ}$ and the tension $T = m\omega^{2}l$ grows without bound until the string breaks — this is the practical limit.

\textbf{3.} From $T^{2} = \dfrac{4\pi^{2}h}{g}$: $g = \dfrac{4\pi^{2}h}{T^{2}} = \dfrac{4\pi^{2}\times 0.25}{1.00^{2}} = \pi^{2} \approx 9.87\ \mathrm{m\,s^{-2}}$.

\textbf{4.} The period depends only on $h$: $T = 2\pi\sqrt{\dfrac{0.50}{9.81}} \approx 1.42\ \mathrm{s}$ for both pendulums, since both have $h = 0.50\ \mathrm{m}$. From $\cos\theta = \dfrac{h}{l}$: for $l = 1.0\ \mathrm{m}$, $\cos\theta = 0.50$, so $\theta = 60^{\circ}$; for $l = 4.0\ \mathrm{m}$, $\cos\theta = 0.125$, so $\theta \approx 82.8^{\circ}$. The longer string hangs much closer to the horizontal, yet both periods are equal — a striking confirmation of the theory.
\end{corrige}

\begin{aretenir}[Summary of the section]
\begin{itemize}
\item Dividing the equations of motion gives $\tan\theta = \dfrac{\omega^{2}r}{g}$; with $r = l\sin\theta$ this yields $\omega^{2} = \dfrac{g}{l\cos\theta} = \dfrac{g}{h}$.
\item Period: $T = 2\pi\sqrt{\dfrac{l\cos\theta}{g}} = 2\pi\sqrt{\dfrac{h}{g}}$ — independent of the mass, depending only on the height $h$.
\item Tension: $T = \dfrac{mg}{\cos\theta} = m\omega^{2}l$, increasing with $\theta$.
\item Limits: $\theta \to 0$ recovers the simple-pendulum period $2\pi\sqrt{l/g}$; $\theta = 90^{\circ}$ is impossible.
\item A graph of $T$ against $\sqrt{h}$ (or $T^{2}$ against $h$) is a straight line through the origin — the basis of an experimental measurement of $g$.
\end{itemize}
\end{aretenir}

\coursec{Angular and Linear Speed in the Conical Pendulum; Synthesis and Exam Practice}

In the previous section we derived the period of the conical pendulum, $T = 2\pi\sqrt{h/g}$, where $h = l\cos\theta$ is the vertical distance from the point of suspension to the plane of the circle. That formula hides a rich structure: the bob has \emph{two} speeds (angular and linear), and both are completely determined by the geometry of the cone. In this final section we collect everything, express $v$ and $\omega$ in terms of $l$, $\theta$ and $g$, learn how the cone responds when it is spun faster, and finish with a complete problem-solving strategy and timed GCE-style practice.

\courssub{The Two Speeds of the Bob: A Recap}

Imagine the bob of a conical pendulum, like a small stone on a string whirled above your head. Two different observers describe its motion:

\begin{itemize}
\item An observer looking \emph{down from above} sees the bob sweeping out a circle of radius $r = l\sin\theta$. The radius vector rotates about the vertical axis with \cle{angular speed} $\omega$ (in $\mathrm{rad\,s^{-1}}$).
\item An observer standing \emph{in the plane of the circle} sees the bob moving along the circle with \cle{linear speed} $v$ (in $\mathrm{m\,s^{-1}}$), tangential to the circle at every instant.
\end{itemize}

The two are linked by the fundamental relation of circular motion:

\begin{aretenir}[The two speeds]
For the bob of a conical pendulum moving on a circle of radius $r = l\sin\theta$:
\[ v = r\omega = l\omega\sin\theta, \qquad \omega = \frac{2\pi}{T} = 2\pi f. \]
The angular speed $\omega$ is the same for every point of the string; the linear speed $v$ concerns the bob alone.
\end{aretenir}

\courssub{Expressing $v$ and $\omega$ in Terms of the Geometry}

Recall the two equations of motion obtained in Section 5 by resolving the tension $S$ in the string (vertical equilibrium and horizontal centripetal force):
\begin{align*}
S\cos\theta &= mg \quad \text{(vertical, no acceleration)}\\
S\sin\theta &= \frac{mv^{2}}{r} = mr\omega^{2} \quad \text{(horizontal, towards the centre).}
\end{align*}
Dividing the second equation by the first eliminates $S$ and $m$:
\[ \tan\theta = \frac{v^{2}}{rg} = \frac{r\omega^{2}}{g}. \]

This single relation is the key to the whole section. Solving it for $v$ and for $\omega$ gives:

\begin{propriete}[Speeds of the conical pendulum in terms of the geometry]
For a conical pendulum of length $l$, semi-vertical angle $\theta$, radius $r = l\sin\theta$ and height $h = l\cos\theta$:
\[ v = \sqrt{rg\tan\theta} = \sqrt{gl\sin\theta\tan\theta} \qquad \omega = \sqrt{\frac{g}{l\cos\theta}} = \sqrt{\frac{g}{h}} \]
\emph{Derivation (examinable).} From $\tan\theta = v^{2}/(rg)$ we get $v^{2} = rg\tan\theta$, hence $v = \sqrt{rg\tan\theta}$; substituting $r = l\sin\theta$ gives the second form. From $\tan\theta = r\omega^{2}/g$ with $r = l\sin\theta$: $\omega^{2} = \dfrac{g\tan\theta}{l\sin\theta} = \dfrac{g\sin\theta}{l\sin\theta\cos\theta} = \dfrac{g}{l\cos\theta}$, so $\omega = \sqrt{g/(l\cos\theta)} = \sqrt{g/h}$.
\end{propriete}

Notice something remarkable: $\omega$ depends only on the \emph{vertical height} $h$ of the suspension point above the plane of the circle, not on the mass, nor directly on the radius. Two conical pendulums whose bobs circle in the same horizontal plane have the \emph{same} angular speed, whatever their lengths.

A second consequence: since $\cos\theta \leq 1$, we have $\omega^{2} = g/(l\cos\theta) \geq g/l$. A conical pendulum of length $l$ can only exist for
\[ \omega \geq \sqrt{\frac{g}{l}}. \]
Below this minimum angular speed the string simply hangs vertically ($\theta = 0$) and there is no circular motion at all.

\begin{attention}[Angle measured from the vertical]
Throughout this chapter $\theta$ is the angle between the string and the \textbf{vertical}. If a problem gives the angle $\alpha$ between the string and the \emph{horizontal}, convert first: $\theta = 90^{\circ} - \alpha$. Substituting $\alpha$ directly into $\tan\theta$ or $\cos\theta$ is one of the most frequent errors at the GCE.
\end{attention}

\courssub{Spinning Faster: How the Cone Widens}

From $\omega = \sqrt{g/(l\cos\theta)}$ we read off the qualitative behaviour:
\begin{itemize}
\item If $\omega$ increases, then $\cos\theta$ must decrease, so $\theta$ \emph{increases}: the string rises towards the horizontal.
\item Then $r = l\sin\theta$ increases and $v = \sqrt{rg\tan\theta}$ increases even faster (both $r$ and $\tan\theta$ grow).
\item As $\omega \to \infty$, $\cos\theta \to 0$, i.e. $\theta \to 90^{\circ}$: the string can never become exactly horizontal at finite speed, because a horizontal string would have no vertical component of tension to support the weight.
\end{itemize}

\begin{popfigure}
\begin{center}
\begin{tikzpicture}[scale=0.85]
% Slow pendulum (left)
\coordinate (O1) at (0,0);
\fill[popdark] (O1) circle (0.05);
\draw[dashed, popmuted] (O1) -- (0,-3.2);
\draw[thick, popblue] (O1) -- ({2.6*sin(25)},{-2.6*cos(25)}) node[midway, above left, popblue] {$l$};
\fill[popblue] ({2.6*sin(25)},{-2.6*cos(25)}) circle (0.09);
\draw[popblue, dashed] (0,{-2.6*cos(25)}) ellipse ({2.6*sin(25)} and {0.35});
\draw[<->, popdark] (0.55,0) arc (0:-65:0.55) node[midway, right] {$\theta_1$};
\draw[<->, popmuted] (-0.6,0) -- (-0.6,{-2.6*cos(25)}) node[midway, left] {$h_1$};
\node[popblue] at (0,-3.7) {slow: $\omega_1$ small};
% Fast pendulum (right)
\begin{scope}[xshift=7.5cm]
\coordinate (O2) at (0,0);
\fill[popdark] (O2) circle (0.05);
\draw[dashed, popmuted] (O2) -- (0,-3.2);
\draw[thick, poporange] (O2) -- ({2.6*sin(60)},{-2.6*cos(60)}) node[midway, above left, poporange] {$l$};
\fill[poporange] ({2.6*sin(60)},{-2.6*cos(60)}) circle (0.09);
\draw[poporange, dashed] (0,{-2.6*cos(60)}) ellipse ({2.6*sin(60)} and {0.5});
\draw[<->, popdark] (0.55,0) arc (0:-30:0.55) node[midway, right] {$\theta_2$};
\draw[<->, popmuted] (3.1,0) -- (3.1,{-2.6*cos(60)}) node[midway, right] {$h_2$};
\node[poporange] at (0,-3.7) {fast: $\omega_2 > \omega_1$};
\end{scope}
\end{tikzpicture}
\end{center}
\legende{Same string, two rotation rates: as $\omega$ increases, $\theta$ increases ($\theta_2>\theta_1$), the circle widens ($r$ increases) and the plane of the circle rises ($h$ decreases).}
\end{popfigure}

\courssub{Strategy for Mixed Problems}

Typical GCE questions give you any two of $(l, \theta, \omega, v, T)$ and ask for the others. The toolbox contains exactly three independent relations:
\[ \tan\theta = \frac{r\omega^{2}}{g}, \qquad v = r\omega, \qquad r = l\sin\theta \ \ (\text{plus } \omega = 2\pi/T). \]

\begin{methode}[Solving any conical pendulum problem]
\begin{enumerate}
\item \textbf{Draw} the free-body diagram: weight $mg$ down, tension $S$ along the string; mark $\theta$ from the vertical.
\item \textbf{Write} the two resolved equations: $S\cos\theta = mg$ and $S\sin\theta = mr\omega^{2}$ (or $mv^{2}/r$).
\item \textbf{Eliminate} $S$ by dividing: $\tan\theta = r\omega^{2}/g$.
\item \textbf{Substitute} the geometry $r = l\sin\theta$ and the data; solve for the unknown.
\item \textbf{Convert} if needed: $v = r\omega$, $T = 2\pi/\omega$.
\item \textbf{Check} units and that $\theta < 90^{\circ}$, $\omega$ in $\mathrm{rad\,s^{-1}}$, and $\omega \geq \sqrt{g/l}$.
\end{enumerate}
\end{methode}

\begin{exoresolu}[Mixed problem: given $l$ and $T$]
A conical pendulum consists of a bob of mass $0.4\ \mathrm{kg}$ on a light inextensible string of length $1.5\ \mathrm{m}$. It describes a horizontal circle with period $T = 2.0\ \mathrm{s}$. Take $g = 9.81\ \mathrm{m\,s^{-2}}$. Find (a) the angular speed, (b) the semi-vertical angle, (c) the radius and linear speed of the bob, (d) the tension in the string.
\tcblower
\textbf{(a)} $\omega = \dfrac{2\pi}{T} = \dfrac{2\pi}{2.0} = \pi \approx 3.14\ \mathrm{rad\,s^{-1}}$.

\textbf{(b)} From $\omega^{2} = \dfrac{g}{l\cos\theta}$: $\cos\theta = \dfrac{g}{l\omega^{2}} = \dfrac{9.81}{1.5\pi^{2}} = \dfrac{9.81}{14.80} = 0.663$, so $\theta = \cos^{-1}(0.663) \approx 48.5^{\circ}$.

\textbf{(c)} $r = l\sin\theta = 1.5\sin 48.5^{\circ} \approx 1.12\ \mathrm{m}$; then $v = r\omega = 1.12 \times 3.14 \approx 3.53\ \mathrm{m\,s^{-1}}$.
\emph{Check with the other formula:} $v = \sqrt{rg\tan\theta} = \sqrt{1.12 \times 9.81 \times 1.131} \approx 3.53\ \mathrm{m\,s^{-1}}$. \checkmark

\textbf{(d)} $S = \dfrac{mg}{\cos\theta} = \dfrac{0.4 \times 9.81}{0.663} \approx 5.92\ \mathrm{N}$.
\end{exoresolu}

\courssub{Comparison Problems: Two Pendulums Rotating Together}

A classic examination situation: two conical pendulums of \emph{different} lengths $l_1 > l_2$ are attached to the same point (or driven with the same angular speed $\omega$). Which makes the larger angle with the vertical?

Since $\omega$ is common and $\omega^{2} = \dfrac{g}{l\cos\theta}$, the product $l\cos\theta = \dfrac{g}{\omega^{2}}$ must be the \emph{same} for both. Hence the longer string must have the \emph{smaller} $\cos\theta$, i.e. the \textbf{larger} angle $\theta$:
\[ l_1\cos\theta_1 = l_2\cos\theta_2 = \frac{g}{\omega^{2}}. \]
Geometrically: both bobs rotate in the \emph{same horizontal plane} at depth $h = g/\omega^{2}$ below the support. The longer string must lean further out to reach that plane.

\begin{exemplebox}[Numerical comparison]
Two strings of lengths $1.0\ \mathrm{m}$ and $2.0\ \mathrm{m}$ rotate as conical pendulums with the same $\omega = 4.0\ \mathrm{rad\,s^{-1}}$. Then $h = g/\omega^{2} = 9.81/16 = 0.613\ \mathrm{m}$ for both. So $\cos\theta_1 = 0.613/1.0 = 0.613 \Rightarrow \theta_1 \approx 52.2^{\circ}$, while $\cos\theta_2 = 0.613/2.0 = 0.307 \Rightarrow \theta_2 \approx 72.1^{\circ}$. The longer pendulum indeed makes the larger angle. (Both values of $\omega$ exceed the minima $\sqrt{g/l}$, namely $3.13$ and $2.21\ \mathrm{rad\,s^{-1}}$, so both motions are possible.)
\end{exemplebox}

\courssub{Chapter Summary Map}

\begin{popfigure}
\begin{center}
\begin{tikzpicture}[node distance=0.55cm and 1.1cm,
box/.style={draw=popblue, thick, rounded corners, fill=popblueL, text width=3.6cm, align=center, minimum height=0.9cm, font=\small},
arr/.style={->, thick, popdark}]
\node[box, align=center] (ang){Angular quantities\\ $\omega = \dfrac{2\pi}{T} = 2\pi f$};
\node[box, right=of ang, align=center] (speed){Linear speed\\ $v = r\omega$};
\node[box, right=of speed, align=center] (acc){Centripetal acceleration\\ $a = \dfrac{v^{2}}{r} = r\omega^{2}$};
\node[box, below=of acc, align=center] (force){Centripetal force\\ $F = \dfrac{mv^{2}}{r} = mr\omega^{2}$};
\node[box, below=of force, align=center] (cone){Conical pendulum\\ $\tan\theta = \dfrac{r\omega^{2}}{g}$};
\node[box, below=of cone, align=center] (period){Period and speeds\\ $T = 2\pi\sqrt{\dfrac{h}{g}}$, $\omega = \sqrt{\dfrac{g}{h}}$, $v=\sqrt{rg\tan\theta}$};
\draw[arr] (ang) -- (speed);
\draw[arr] (speed) -- (acc);
\draw[arr] (acc) -- (force);
\draw[arr] (force) -- (cone);
\draw[arr] (cone) -- (period);
\end{tikzpicture}
\end{center}
\legende{Concept map of the chapter: each formula follows from the previous one; the conical pendulum is the centripetal-force law applied to a bob on a string.}
\end{popfigure}

\begin{aretenir}[The five key formulae of the chapter]
\begin{enumerate}
\item $\omega = \dfrac{2\pi}{T} = 2\pi f$ \quad (angular speed)
\item $v = r\omega$ \quad (linear speed on the circle)
\item $a = \dfrac{v^{2}}{r} = r\omega^{2}$ \quad (centripetal acceleration, towards the centre)
\item $F = \dfrac{mv^{2}}{r} = mr\omega^{2}$ \quad (centripetal force)
\item Conical pendulum: $\tan\theta = \dfrac{r\omega^{2}}{g}$, giving $T = 2\pi\sqrt{\dfrac{h}{g}}$, $\omega = \sqrt{\dfrac{g}{l\cos\theta}}$, $v = \sqrt{rg\tan\theta}$, with $\omega \geq \sqrt{g/l}$.
\end{enumerate}
\end{aretenir}

\courssub{Timed GCE-Style Practice}

\begin{attention}[Exam technique]
GCE marking schemes award \textbf{method marks} for the free-body diagram and the two resolved equations $S\cos\theta = mg$, $S\sin\theta = mr\omega^{2}$ \emph{before} any numerical work. Always show them, even if the final calculation goes wrong. Quote $\theta$ from the vertical, give $\omega$ in $\mathrm{rad\,s^{-1}}$, and carry at least 3 significant figures.
\end{attention}

\begin{exercice}[Structured question — 25 minutes]
A small bob of mass $0.25\ \mathrm{kg}$ is attached to a light inextensible string of length $0.80\ \mathrm{m}$ fixed at a point $O$. The bob moves in a horizontal circle with constant angular speed $\omega = 5.0\ \mathrm{rad\,s^{-1}}$. Take $g = 9.81\ \mathrm{m\,s^{-2}}$.
\begin{enumerate}
\item Define the terms \emph{angular speed} and \emph{centripetal acceleration}.
\item Draw the free-body diagram of the bob and derive the relation $\tan\theta = \dfrac{r\omega^{2}}{g}$, where $\theta$ is the angle between the string and the vertical.
\item Calculate (i) the angle $\theta$, (ii) the radius $r$ of the circle, (iii) the linear speed $v$ of the bob, (iv) the tension in the string.
\item The angular speed is now doubled to $10\ \mathrm{rad\,s^{-1}}$. Without full calculation, state and justify what happens to $\theta$, $r$, $v$ and the height $h$ of $O$ above the plane of the circle.
\end{enumerate}
\end{exercice}

\begin{corrige}[Exercise 1]
\textbf{(a)} Angular speed is the angle swept out per unit time, $\omega = 2\pi/T$ ($\mathrm{rad\,s^{-1}}$). Centripetal acceleration is the acceleration directed towards the centre of the circle, of magnitude $v^{2}/r = r\omega^{2}$, responsible for changing the direction of the velocity.

\textbf{(b)} Forces on the bob: weight $mg$ vertically down, tension $S$ along the string. Vertical equilibrium: $S\cos\theta = mg$. Horizontal motion (towards centre): $S\sin\theta = mr\omega^{2}$. Dividing: $\tan\theta = r\omega^{2}/g$.

\textbf{(c)} (i) $\cos\theta = \dfrac{g}{l\omega^{2}} = \dfrac{9.81}{0.80 \times 25} = 0.4905$, so $\theta \approx 60.6^{\circ}$. (ii) $r = l\sin\theta = 0.80 \times 0.8715 \approx 0.697\ \mathrm{m}$. (iii) $v = r\omega = 0.697 \times 5.0 \approx 3.49\ \mathrm{m\,s^{-1}}$. (iv) $S = \dfrac{mg}{\cos\theta} = \dfrac{0.25 \times 9.81}{0.4905} \approx 5.0\ \mathrm{N}$.

\textbf{(d)} Since $\cos\theta = g/(l\omega^{2})$, doubling $\omega$ divides $\cos\theta$ by $4$: $\theta$ increases towards $90^{\circ}$. Hence $r = l\sin\theta$ increases, $h = l\cos\theta$ decreases (the circle rises), and $v = \sqrt{rg\tan\theta}$ increases sharply. Numerically $\cos\theta = 0.1226$, $\theta \approx 83.0^{\circ}$, $r \approx 0.794\ \mathrm{m}$, $v \approx 7.9\ \mathrm{m\,s^{-1}}$.
\end{corrige}

\begin{exercice}[Comparison question — 15 minutes]
Two conical pendulums are suspended from the same point. Pendulum A has length $1.2\ \mathrm{m}$ and pendulum B has length $0.6\ \mathrm{m}$. They are set rotating with the \emph{same} angular speed $\omega = 3.5\ \mathrm{rad\,s^{-1}}$.
\begin{enumerate}
\item Show that any bob rotating at this angular speed must move in a horizontal plane at depth $h = g/\omega^{2}$ below the point of suspension, and find this depth.
\item Determine the semi-vertical angle of each pendulum, commenting on any difficulty you meet.
\item Which bob has the greater linear speed? Justify.
\end{enumerate}
\end{exercice}

\begin{corrige}[Exercise 2]
\textbf{(a)} From $\omega^{2} = g/h$ we get $h = g/\omega^{2} = 9.81/12.25 = 0.80\ \mathrm{m}$, which depends only on $\omega$ and $g$, not on the length of the string. Any bob rotating at $3.5\ \mathrm{rad\,s^{-1}}$ circles in the horizontal plane $0.80\ \mathrm{m}$ below the support.

\textbf{(b)} A: $\cos\theta_A = h/l_A = 0.80/1.2 = 0.667 \Rightarrow \theta_A \approx 48.2^{\circ}$. B: $\cos\theta_B = 0.80/0.6 = 1.33 > 1$ — impossible! Pendulum B \emph{cannot} rotate at $\omega = 3.5\ \mathrm{rad\,s^{-1}}$: its minimum possible angular speed is $\omega_{\min} = \sqrt{g/l} = \sqrt{9.81/0.6} \approx 4.04\ \mathrm{rad\,s^{-1}}$. Below this value the string simply hangs vertically.

\textbf{(c)} Only pendulum A rotates: $v_A = \sqrt{rg\tan\theta_A}$ with $r = 1.2\sin 48.2^{\circ} = 0.894\ \mathrm{m}$ and $\tan 48.2^{\circ} = 1.118$, giving $v_A \approx \sqrt{0.894 \times 9.81 \times 1.118} \approx 3.13\ \mathrm{m\,s^{-1}}$. The question illustrates that a conical pendulum of length $l$ requires $\omega \geq \sqrt{g/l}$.
\end{corrige}

\begin{savaistu}[Why the conical pendulum matters]
The conical pendulum is not just an exam topic: it was at the heart of Huygens' work on pendulum clocks in the seventeenth century, and the same physics governs the ``flying chairs'' fairground ride, the centrifugal governor of a steam engine, and the banking of aircraft and of cyclists in a velodrome — in each case the horizontal component of an inclined force supplies the centripetal force $mr\omega^{2}$.
\end{savaistu}

\newpage

\coursec{Integration activity}

\courssub{Aims of the activity}

This integration activity brings together, in one authentic engineering task, all the notions developed in this chapter:

\begin{itemize}
\item \cle{Angular quantities}: converting a rotation rate from revolutions per minute to angular speed in $\mathrm{rad\,s^{-1}}$, using $\omega = \dfrac{2\pi N}{60}$.
\item \cle{Uniform circular motion}: identifying the radius of the circle described by a seat and computing its \cle{linear speed} using $v = r\omega$.
\item \cle{Centripetal acceleration and force}: recognising that the horizontal component of the chain's tension supplies the centripetal force $F = m r \omega^2$.
\item \cle{The conical pendulum model}: resolving forces vertically and horizontally, and using the results
\[ T\cos\theta = mg, \qquad T\sin\theta = m r \omega^2, \qquad \tan\theta = \frac{r\omega^2}{g}, \qquad T = \frac{mg}{\cos\theta}. \]
\item \cle{Scientific communication}: presenting a calculation with free-body diagrams, checking a safety constraint, and formulating a justified recommendation, as a real engineer would.
\item \cle{Critical thinking}: discussing the limits of a model (air resistance, mass of the chains, gusts of wind) and the notion of a safety margin.
\end{itemize}

By the end of the activity, you should be able to pass fluently from a real rotating system to its conical-pendulum model, to extract numerical predictions from it, and to judge whether those predictions can be trusted.

\courssub{Situation}

\textbf{Designing a Safe Swing Carousel for a Yaoundé Amusement Park.}\par\medskip

A new amusement park is about to open on the outskirts of Yaoundé, not far from the Carrefour Nsam. Among the attractions ordered by the park management is a \emph{swing carousel} (the famous ``chaises volantes''): a vertical mast topped by a rotating crown from which chains hang. Each chain carries a seat, and when the crown rotates, the seats swing outwards and describe horizontal circles, to the delight of the riders.

The technical data supplied by the manufacturer are the following:

\begin{center}
\begin{tabularx}{\linewidth}{|l|X|}
\hline
\rowcolor{popblueL} \textbf{Quantity} & \textbf{Value} \\
\hline
Length of each chain (from crown to seat) & $l = 4.0\ \mathrm{m}$ \\
\hline
Horizontal distance from the axis of rotation to the attachment point of a chain on the crown & $a = 2.0\ \mathrm{m}$ \\
\hline
Maximum authorised mass per seat (child + seat) & $m = 50\ \mathrm{kg}$ \\
\hline
Breaking safety limit of one chain (tension must stay below this) & $T_{\max} = 800\ \mathrm{N}$ \\
\hline
Maximum angle allowed between a chain and the vertical (comfort and clearance) & $\theta_{\max} = 45^{\circ}$ \\
\hline
Normal operating rotation rate & $N = 12$ revolutions per minute (rpm) \\
\hline
Acceleration due to gravity & $g = 9.81\ \mathrm{m\,s^{-2}}$ \\
\hline
\end{tabularx}
\end{center}

The park manager is worried. During a test run, a journalist claimed that ``the chains will snap at full speed''. The manager hires your group of Upper Sixth students as consultants. Your mission: model each seat (with its rider) as a \cle{conical pendulum}, compute the relevant quantities at the normal operating rate, verify the safety of the chains, determine the fastest rotation rate that keeps the chains within the allowed angle, and write a short recommendation.

\begin{attention}[Modelling assumptions]
Throughout the activity, assume that: the rotation is uniform (constant angular speed); each seat moves in a \emph{horizontal} circle at constant height; the chains are light (negligible mass) and inextensible; air resistance is neglected; the rider plus seat is treated as a single particle of mass $m = 50\ \mathrm{kg}$ fixed at the lower end of the chain. The radius of the circle is $r = a + l\sin\theta$, because the chain is attached at a distance $a = 2.0\ \mathrm{m}$ from the axis. Forgetting the offset $a$ is the most common error in this type of problem.
\end{attention}

\courssub{Tasks}

\textbf{Part A — Setting up the model}

\begin{enumerate}
\item Draw a clear diagram of one seat in operation: the axis of rotation, the crown, the chain of length $l$ making an angle $\theta$ with the vertical, and the circle of radius $r = a + l\sin\theta$ described by the seat.
\item Draw the free-body diagram of the seat + rider system. Name each force and state which body exerts it.
\item Explain carefully why the vertical forces must balance, and why the resultant force must be horizontal, pointing towards the axis. What is the name of this resultant?
\item By resolving the tension vertically and horizontally, establish the two equations
\[ T\cos\theta = mg \qquad \text{and} \qquad T\sin\theta = m r \omega^2, \]
and deduce that $\tan\theta = \dfrac{r\omega^2}{g}$.
\end{enumerate}

\textbf{Part B — Normal operation at 12 rpm}

\begin{enumerate}
\setcounter{enumi}{4}
\item Convert the rotation rate $N = 12$ rpm into an angular speed $\omega$ in $\mathrm{rad\,s^{-1}}$.
\item Using $\tan\theta = \dfrac{(a + l\sin\theta)\omega^2}{g}$, show that $\theta$ satisfies the equation
\[ 9.81\tan\theta = 1.579\left(2.0 + 4.0\sin\theta\right). \]
Solve this equation (by trial with your calculator, or graphically) to find $\theta$ to the nearest degree.
\item Hence compute: the radius $r$ of the circle; the linear speed $v = r\omega$ of a seat; the period $T_0$ of one revolution (check that it agrees with $N = 12$ rpm); the tension $T$ in a chain.
\item \textbf{Safety check.} Compare the tension with the chain's limit of $800\ \mathrm{N}$. Compute the \emph{safety factor} $\dfrac{T_{\max}}{T}$. Is the journalist's claim justified at 12 rpm?
\end{enumerate}

\textbf{Part C — Maximum rotation rate}

\begin{enumerate}
\setcounter{enumi}{8}
\item Using the full relation $\tan\theta = \dfrac{(a + l\sin\theta)\omega^2}{g}$, determine the angular speed $\omega_{\max}$ for which $\theta = 45^{\circ}$. Give $\omega_{\max}$ in $\mathrm{rad\,s^{-1}}$.
\item Convert $\omega_{\max}$ into rpm and compute the corresponding linear speed of a seat and the tension in a chain. Verify that the chain is still safe at this limiting rate, and state which constraint (angle or tension) is the binding one.
\end{enumerate}

\textbf{Part D — Communication and critical discussion}

\begin{enumerate}
\setcounter{enumi}{11}
\item Write a short recommendation (5 to 8 lines) to the park manager, stating whether the carousel is safe at 12 rpm, what rotation rate must never be exceeded, and what instruction should be given to the operator.
\item Discuss two or three reasons why the real situation differs from the model (air resistance on the riders, mass of the chains, wind gusts, a rider who leans out). For each, say whether it \emph{increases} or \emph{decreases} the tension compared with your calculation, and explain why the manufacturer's safety limit of $800\ \mathrm{N}$ (well below the true breaking load) protects against these effects.
\end{enumerate}

\courssub{Model answer}

\textbf{Part A.}

\textbf{A1.} Diagram of one seat:

\begin{popfigure}
\begin{tikzpicture}[scale=0.9]
% axis
\draw[thick,popmuted] (0,-0.4) -- (0,3.6) node[above]{\small axis};
% crown
\draw[very thick,popdark] (-0.15,3.2) -- (2.0,3.2) node[midway,above]{\small crown};
\fill[popdark] (2.0,3.2) circle (2pt);
% chain
\draw[very thick,popblue] (2.0,3.2) -- (4.83,0.37) node[midway,above right=-2pt]{\small chain $l$};
% seat
\fill[poporange] (4.83,0.37) circle (4pt) node[right]{\small seat $m$};
% vertical dashed
\draw[dashed,popmuted] (2.0,3.2) -- (2.0,0.37);
% angle
\draw[popdark] (2.0,2.2) arc (-90:-35:1.0) node[midway,below right=-4pt]{\small $\theta$};
% radius
\draw[<->,popgreen] (0,-0.25) -- (4.83,-0.25) node[midway,below]{\small $r = a + l\sin\theta$};
\draw[<->,popmuted] (0,2.9) -- (2.0,2.9) node[midway,above]{\small $a$};
% circle hint
\draw[dashed,popgreen] (4.83,0.37) ellipse (0.55 and 0.18);
\end{tikzpicture}
\legende{One seat of the carousel modelled as a conical pendulum.}
\end{popfigure}

\textbf{A2.} Free-body diagram: two forces act on the system seat + rider:
\begin{itemize}
\item its \cle{weight} $\vec{P} = m\vec{g}$, exerted by the Earth, vertically downwards;
\item the \cle{tension} $\vec{T}$ of the chain, exerted by the chain, directed along the chain towards the attachment point.
\end{itemize}

\begin{popfigure}
\begin{tikzpicture}[scale=1.0]
\fill[poporange] (0,0) circle (4pt) node[below left]{\small seat};
\draw[->,very thick,popdark] (0,0) -- (0,-2.2) node[below]{\small $\vec{P}=m\vec{g}$};
\draw[->,very thick,popblue] (0,0) -- (-2.0,1.9) node[above left]{\small $\vec{T}$};
\draw[->,thick,popblue,dashed] (0,0) -- (-1.9,0) node[midway,below]{\small $T\sin\theta$};
\draw[->,thick,popblue,dashed] (-1.9,0) -- (-2.0,1.9) node[midway,left]{\small $T\cos\theta$};
\draw[->,very thick,popgreen] (0,0) -- (-1.3,0) node[midway,above=-2pt]{\small $\vec{F}_c$};
\end{tikzpicture}
\legende{Free-body diagram: the horizontal component $T\sin\theta$ is the centripetal force.}
\end{popfigure}

\textbf{A3.} The seat moves in a \emph{horizontal} circle at constant height: its vertical acceleration is zero, so the vertical components of the forces must balance. But uniform circular motion requires a \cle{centripetal acceleration} $r\omega^2$ directed towards the centre of the circle; by Newton's second law, the resultant force must therefore be horizontal, of magnitude $mr\omega^2$, pointing towards the axis. This resultant is the \cle{centripetal force}; here it is provided entirely by the horizontal component of the tension. Note that the centripetal force is not an extra force to add to the diagram: it is the \emph{resultant} of the real forces (weight and tension).

\textbf{A4.} Resolving $\vec{T}$ into a vertical component $T\cos\theta$ and a horizontal component $T\sin\theta$:
\begin{align*}
\text{vertically (no vertical acceleration):}\quad & T\cos\theta = mg,\\
\text{horizontally (Newton's second law):}\quad & T\sin\theta = mr\omega^2.
\end{align*}
Dividing the second equation by the first eliminates $T$ and $m$:
\[ \tan\theta = \frac{r\omega^2}{g}. \]

\textbf{Part B.}

\textbf{B5.} One revolution is $2\pi$ rad and one minute is $60\ \mathrm{s}$, so
\[ \omega = \frac{2\pi N}{60} = \frac{12 \times 2\pi}{60} = \frac{24\pi}{60} \approx 1.257\ \mathrm{rad\,s^{-1}}, \qquad \omega^2 \approx 1.579\ \mathrm{rad^2\,s^{-2}}. \]

\textbf{B6.} With $r = 2.0 + 4.0\sin\theta$ and $\omega^2 \approx 1.579$:
\[ \tan\theta = \frac{1.579\left(2.0 + 4.0\sin\theta\right)}{9.81} \quad\Longleftrightarrow\quad 9.81\tan\theta = 1.579\left(2.0 + 4.0\sin\theta\right). \]
This equation cannot be solved directly for $\theta$ (it contains both $\tan\theta$ and $\sin\theta$), so we proceed by trial:
\begin{itemize}
\item $\theta = 30^{\circ}$: LHS $= 9.81 \times 0.5774 \approx 5.66$; RHS $= 1.579 \times (2.0 + 4.0 \times 0.500) = 1.579 \times 4.000 \approx 6.32$. LHS $<$ RHS.
\item $\theta = 35^{\circ}$: LHS $= 9.81 \times 0.7002 \approx 6.87$; RHS $= 1.579 \times (2.0 + 4.0 \times 0.5736) = 1.579 \times 4.294 \approx 6.78$. LHS $>$ RHS.
\item $\theta = 34^{\circ}$: LHS $= 9.81 \times 0.6745 \approx 6.62$; RHS $= 1.579 \times (2.0 + 4.0 \times 0.5592) = 1.579 \times 4.237 \approx 6.69$. LHS $<$ RHS.
\end{itemize}
The solution lies between $34^{\circ}$ and $35^{\circ}$. Refining at $\theta = 34.4^{\circ}$: LHS $\approx 9.81 \times 0.6836 \approx 6.71$; RHS $\approx 1.579 \times (2.0 + 4.0 \times 0.5652) \approx 1.579 \times 4.261 \approx 6.73$. The two sides agree to within $0.3\%$, so
\[ \boxed{\theta \approx 34^{\circ}\ \text{(more precisely } 34.4^{\circ}\text{)}}. \]

\textbf{B7.} With $\theta \approx 34.4^{\circ}$, $\sin\theta \approx 0.565$ and $\cos\theta \approx 0.825$:
\begin{align*}
r &= 2.0 + 4.0 \times 0.565 \approx 4.26\ \mathrm{m},\\
v &= r\omega \approx 4.26 \times 1.257 \approx 5.35\ \mathrm{m\,s^{-1}} \ (\approx 19\ \mathrm{km\,h^{-1}}),\\
T_0 &= \frac{2\pi}{\omega} = \frac{2\pi}{1.257} \approx 5.0\ \mathrm{s} \quad \left(\text{indeed } \frac{60\ \mathrm{s}}{12} = 5.0\ \mathrm{s\ per\ revolution}\right),\\
T &= \frac{mg}{\cos\theta} = \frac{50 \times 9.81}{0.825} \approx \frac{490.5}{0.825} \approx 594\ \mathrm{N}.
\end{align*}
\emph{Consistency check:} $T\sin\theta \approx 594 \times 0.565 \approx 336\ \mathrm{N}$ and $mr\omega^2 \approx 50 \times 4.26 \times 1.579 \approx 336\ \mathrm{N}$. The two expressions of the centripetal force agree.

\textbf{B8.} Safety check: $T \approx 594\ \mathrm{N} < 800\ \mathrm{N} = T_{\max}$. The safety factor is
\[ \frac{T_{\max}}{T} \approx \frac{800}{594} \approx 1.35. \]
The chain operates at about $74\%$ of its authorised limit: at 12 rpm the journalist's claim is \textbf{not} justified.

\textbf{Part C.}

\textbf{C9.} Set $\theta = 45^{\circ}$ in the full relation, with
\[ r = 2.0 + 4.0\sin 45^{\circ} = 2.0 + 4.0 \times 0.7071 = 4.828\ \mathrm{m}: \]
\[ \omega_{\max}^2 = \frac{g\tan 45^{\circ}}{r} = \frac{9.81 \times 1}{4.828} \approx 2.032 \quad\Longrightarrow\quad \omega_{\max} \approx 1.43\ \mathrm{rad\,s^{-1}}. \]

\textbf{C10.} Conversion and consequences:
\begin{align*}
N_{\max} &= \frac{60\,\omega_{\max}}{2\pi} \approx \frac{60 \times 1.43}{6.283} \approx 13.6\ \mathrm{rpm},\\
v_{\max} &= r\omega_{\max} \approx 4.828 \times 1.43 \approx 6.9\ \mathrm{m\,s^{-1}},\\
T &= \frac{mg}{\cos 45^{\circ}} = \frac{490.5}{0.7071} \approx 694\ \mathrm{N} < 800\ \mathrm{N}.
\end{align*}
Even at the limiting angle, the chain remains within its safety limit ($694\ \mathrm{N} < 800\ \mathrm{N}$, safety factor $\approx 1.15$). The \emph{angle} constraint ($45^{\circ}$) is therefore the binding one: it is reached at about $13.6$ rpm, only slightly above the normal rate of $12$ rpm.

\textbf{Part D.}

\textbf{D11.} \emph{Recommendation to the park manager.} Our calculations, based on the standard conical-pendulum model with the heaviest authorised load ($50\ \mathrm{kg}$), show that at the normal rate of 12 rpm each chain carries a tension of about $594\ \mathrm{N}$, i.e. only $74\%$ of the $800\ \mathrm{N}$ safety limit, and the chains incline at about $34^{\circ}$, below the $45^{\circ}$ clearance limit. The carousel is therefore safe at 12 rpm. However, the angle limit is reached at about 13.6 rpm, so the margin is small: we recommend that the motor be fitted with a speed limiter set at 12 rpm, that the operator never override it, and that chains be inspected weekly for wear.

\textbf{D12.} \emph{Limits of the model.}
\begin{itemize}
\item \textbf{Air resistance} on the rider opposes the motion; it adds a horizontal drag that the chain must also pull along, slightly \emph{increasing} the tension and making the seat trail backwards (the plane of the circle lags behind the crown arm).
\item \textbf{Mass of the chains} (a real $4\ \mathrm{m}$ chain may weigh several kilograms): the upper part of each chain supports the weight of the lower part, so the tension at the crown is \emph{higher} than our value $T = mg/\cos\theta$ computed for a light chain.
\item \textbf{Wind gusts and a rider leaning outwards} shift the centre of mass and add lateral forces, again \emph{increasing} the tension and the effective angle.
\end{itemize}
All three effects increase the true tension above the model's prediction. This is precisely why the manufacturer fixes the \emph{authorised} limit ($800\ \mathrm{N}$) far below the actual breaking load of the chain: the gap absorbs the unmodelled effects. Engineers call this gap a \cle{safety margin}, and our role as consultants is to check that the modelled tension stays comfortably inside it — which it does at 12 rpm.

\begin{aretenir}[What this activity consolidates]
\begin{itemize}
\item Convert rotation rates: $\omega = \dfrac{2\pi N}{60}$ (with $N$ in rpm, $\omega$ in $\mathrm{rad\,s^{-1}}$).
\item In any uniform circular motion, the resultant force is the centripetal force $mr\omega^2 = \dfrac{mv^2}{r}$, directed towards the centre; it is not an extra force.
\item For a conical pendulum: $T\cos\theta = mg$, $T\sin\theta = mr\omega^2$, hence $\tan\theta = \dfrac{r\omega^2}{g}$ and $T = \dfrac{mg}{\cos\theta}$.
\item The radius must be measured from the \emph{axis of rotation}: here $r = a + l\sin\theta$, a frequent source of error.
\item When an equation mixes $\sin\theta$ and $\tan\theta$, solve it numerically by bracketing the root with a calculator.
\item A good scientific answer ends with a check (comparison with a limit, consistency between formulas) and a critical discussion of the model.
\end{itemize}
\end{aretenir}

\newpage

\coursec{Methods and worked examples}

\courssub{Method 1: Converting and Using Angular Quantities}

In uniform circular motion, the position of the body is most naturally described by the \emph{angle} it has swept, not by a coordinate along a straight line. The first skill is therefore to move fluently between revolutions, degrees and radians, and between frequency, period and angular speed.

\begin{methode}[Handling angular quantities: revolutions, degrees, radians]
To work with angular quantities in circular motion:
\begin{enumerate}
\item \textbf{Choose the SI unit.} Always convert angles to \cle{radians} and angular speed to $\mathrm{rad\,s^{-1}}$ before using any formula. Recall: $1$ revolution $= 2\pi\ \mathrm{rad} = 360^{\circ}$, so $1\ \mathrm{rad} = \dfrac{180^{\circ}}{\pi} \approx 57.3^{\circ}$.
\item \textbf{Convert frequencies.} If a body makes $N$ revolutions per minute (rpm), then the frequency is $f = \dfrac{N}{60}\ \mathrm{Hz}$ and the angular speed is
\[ \omega = 2\pi f = \frac{2\pi N}{60}\ \mathrm{rad\,s^{-1}}. \]
\item \textbf{Use the period.} The period $T$ (time for one complete revolution) satisfies $T = \dfrac{1}{f}$ and $\omega = \dfrac{2\pi}{T}$.
\item \textbf{Angle swept.} In time $t$ at constant angular speed, the angle swept is $\theta = \omega t$ (with $\theta$ in radians).
\item \textbf{Arc length.} The distance travelled along the circle is $s = r\theta = r\omega t$, where $r$ is the radius of the path.
\end{enumerate}
\textbf{Reflex:} as soon as you see ``rpm'' or ``revolutions per second'', convert immediately to $\mathrm{rad\,s^{-1}}$ before any calculation.
\end{methode}

\begin{exoresolu}[Angular speed of a ceiling fan]
The blade of a ceiling fan in a classroom in Buea rotates at $150$ revolutions per minute. The tip of a blade is $0.60\ \mathrm{m}$ from the axis of rotation.
\begin{enumerate}
\item[(a)] Calculate the angular speed of the blade in $\mathrm{rad\,s^{-1}}$.
\item[(b)] Calculate the period of the motion.
\item[(c)] Find the angle, in radians, swept by a blade in $0.50\ \mathrm{s}$, and the distance travelled by the tip in that time.
\end{enumerate}
\tcblower
\textbf{(a) Angular speed.} The frequency is
\[ f = \frac{150}{60} = 2.5\ \mathrm{Hz}. \]
Hence
\[ \omega = 2\pi f = 2\pi \times 2.5 = 5\pi \approx 15.7\ \mathrm{rad\,s^{-1}}. \]

\textbf{(b) Period.}
\[ T = \frac{1}{f} = \frac{1}{2.5} = 0.40\ \mathrm{s}. \]
Check: $T = \dfrac{2\pi}{\omega} = \dfrac{2\pi}{5\pi} = 0.40\ \mathrm{s}$. \checkmark

\textbf{(c) Angle swept and distance.} In $t = 0.50\ \mathrm{s}$:
\[ \theta = \omega t = 5\pi \times 0.50 = 2.5\pi \approx 7.85\ \mathrm{rad}. \]
The distance travelled by the tip along its circular path is
\[ s = r\theta = 0.60 \times 2.5\pi = 1.5\pi \approx 4.71\ \mathrm{m}. \]
\textbf{Remark:} $2.5\pi\ \mathrm{rad} = 1.25$ revolutions, consistent with $f \times t = 2.5 \times 0.5 = 1.25$ revolutions.
\end{exoresolu}

\courssub{Method 2: Linking Linear Speed and Angular Speed}

Every point of a rotating rigid body sweeps the \emph{same angle} in the same time, so all points share the same $\omega$; but points further from the axis travel longer arcs, hence move \emph{faster}. The bridge between the two descriptions is $v = r\omega$.

\begin{methode}[Using $v = r\omega$]
To relate linear and angular quantities for a point on a rotating body:
\begin{enumerate}
\item \textbf{Identify the radius} $r$ of the circular path of the point considered (different points of the same rigid body have different $r$ but the \emph{same} $\omega$).
\item \textbf{Apply the bridge formula} $v = r\omega$, with $v$ in $\mathrm{m\,s^{-1}}$, $r$ in metres, $\omega$ in $\mathrm{rad\,s^{-1}}$.
\item \textbf{Combine with period/frequency} if needed: $v = r\omega = \dfrac{2\pi r}{T} = 2\pi r f$.
\item \textbf{Compare points} on the same rotating body: since $\omega$ is common, $v \propto r$ — the further from the axis, the faster the point moves.
\end{enumerate}
\textbf{Reflex:} $v = r\omega$ only works with $\omega$ in \emph{radians} per second. Never use revolutions or degrees in this formula.
\end{methode}

\begin{exoresolu}[The grinding wheel]
A grinding wheel in a workshop in Douala has diameter $30\ \mathrm{cm}$ and spins at $1200\ \mathrm{rpm}$.
\begin{enumerate}
\item[(a)] Find the angular speed of the wheel.
\item[(b)] Find the linear speed of a point on the rim of the wheel.
\item[(c)] A spark leaves the rim tangentially. Explain why it flies off along the tangent, and state its speed at the instant it leaves.
\end{enumerate}
\tcblower
\textbf{(a)} Radius: $r = \dfrac{30}{2} = 15\ \mathrm{cm} = 0.15\ \mathrm{m}$. Frequency:
\[ f = \frac{1200}{60} = 20\ \mathrm{Hz}, \qquad \omega = 2\pi f = 40\pi \approx 125.7\ \mathrm{rad\,s^{-1}}. \]

\textbf{(b)} Linear speed of a rim point:
\[ v = r\omega = 0.15 \times 40\pi = 6\pi \approx 18.8\ \mathrm{m\,s^{-1}}. \]

\textbf{(c)} In circular motion the velocity vector is always \emph{tangent} to the circle. When the particle of material detaches, the centripetal force on it vanishes, so by Newton's first law it continues in a straight line with the velocity it had at the instant of release: along the tangent, at speed $v \approx 18.8\ \mathrm{m\,s^{-1}}$. This is why grinding wheels are always fitted with guards.
\end{exoresolu}

\courssub{Method 3: Computing Centripetal Acceleration and Centripetal Force}

Even at \emph{constant speed}, a body in circular motion is accelerated, because the \emph{direction} of its velocity changes continuously. The acceleration points towards the centre of the circle; by Newton's second law, the resultant force must also point towards the centre.

\begin{methode}[Applying $a = \dfrac{v^2}{r} = r\omega^2$ and $F = ma$]
To find the centripetal acceleration or force in uniform circular motion:
\begin{enumerate}
\item \textbf{Draw a diagram} and identify the centre of the circular path; the acceleration is directed \emph{towards the centre} at every instant.
\item \textbf{Choose the convenient form:}
\[ a = \frac{v^2}{r} = r\omega^2 = \frac{4\pi^2 r}{T^2}. \]
Use the form matching the data given ($v$, $\omega$ or $T$).
\item \textbf{Apply Newton's second law} along the radius: the resultant force towards the centre is
\[ F = ma = \frac{mv^2}{r} = mr\omega^2. \]
\item \textbf{Identify the physical agent} providing $F$ (tension, friction, normal reaction, gravity, or a component of one of these). The ``centripetal force'' is not a new force: it is the \emph{role} played by the resultant of real forces.
\end{enumerate}
\begin{attention}[Common traps]
Do not add a ``centrifugal force'' acting on the body in an inertial frame. Do not forget to convert $r$ to metres and $\omega$ to $\mathrm{rad\,s^{-1}}$. The acceleration is \emph{not} zero even though the \emph{speed} is constant, because the \emph{velocity direction} changes.
\end{attention}
\end{methode}

\begin{exoresolu}[Stone whirled on a string]
A stone of mass $0.50\ \mathrm{kg}$ is tied to a string of length $1.2\ \mathrm{m}$ and whirled in a horizontal circle at constant speed, making $2$ revolutions per second. Take $g = 9.81\ \mathrm{m\,s^{-2}}$.
\begin{enumerate}
\item[(a)] Calculate the angular speed and the linear speed of the stone.
\item[(b)] Calculate the centripetal acceleration of the stone.
\item[(c)] Find the tension in the string (assume the string is horizontal).
\item[(d)] The string can withstand a maximum tension of $400\ \mathrm{N}$. Find the greatest number of revolutions per second before it breaks.
\end{enumerate}
\tcblower
\textbf{(a)} $f = 2\ \mathrm{Hz}$, so
\[ \omega = 2\pi f = 4\pi \approx 12.57\ \mathrm{rad\,s^{-1}}, \qquad v = r\omega = 1.2 \times 4\pi = 4.8\pi \approx 15.1\ \mathrm{m\,s^{-1}}. \]

\textbf{(b)} Using $a = r\omega^2$:
\[ a = 1.2 \times (4\pi)^2 = 1.2 \times 16\pi^2 = 19.2\pi^2 \approx 189.5\ \mathrm{m\,s^{-2}}. \]
(That is about $19$ times $g$!)

\textbf{(c)} The tension provides the centripetal force:
\[ T = mr\omega^2 = 0.50 \times 189.5 \approx 94.7\ \mathrm{N}. \]

\textbf{(d)} At breaking point, $T_{\max} = mr\omega_{\max}^2$, so
\[ \omega_{\max} = \sqrt{\frac{T_{\max}}{mr}} = \sqrt{\frac{400}{0.50 \times 1.2}} = \sqrt{666.7} \approx 25.8\ \mathrm{rad\,s^{-1}}. \]
Hence
\[ f_{\max} = \frac{\omega_{\max}}{2\pi} = \frac{25.8}{2\pi} \approx 4.1\ \mathrm{rev\,s^{-1}}. \]
The string breaks if the stone is whirled faster than about $4.1$ revolutions per second.
\end{exoresolu}

\courssub{Method 4: Circular Motion of Vehicles — Flat and Banked Tracks}

A vehicle rounding a bend is in (approximately) uniform circular motion: something must push it towards the centre of the bend. On a flat road that ``something'' is friction; on a banked track, a component of the normal reaction helps.

\begin{methode}[Analysing a vehicle rounding a bend]
\begin{enumerate}
\item \textbf{Flat road:} the centripetal force is supplied entirely by friction $F_f$ between tyres and road:
\[ F_f = \frac{mv^2}{r} \le \mu mg \quad\Longrightarrow\quad v_{\max} = \sqrt{\mu g r}. \]
\item \textbf{Banked track (angle $\theta$, no friction needed):} resolve the normal reaction $R$. Vertically: $R\cos\theta = mg$. Horizontally (towards centre): $R\sin\theta = \dfrac{mv^2}{r}$. Dividing:
\[ \tan\theta = \frac{v^2}{gr} \quad\Longrightarrow\quad v = \sqrt{gr\tan\theta}. \]
\item \textbf{Always} identify which real force (or component) points towards the centre before writing Newton's second law.
\item \textbf{Check units and reasonableness:} a larger radius or a steeper banking allows a larger speed.
\end{enumerate}
\end{methode}

\begin{popfigure}
\begin{center}
\begin{tikzpicture}[scale=1.0]
% Banked road cross-section
\fill[popblueL] (0,0) -- (5,1.4) -- (5,1.9) -- (0,0.5) -- cycle;
\draw[thick, popink] (0,0.5) -- (5,1.9);
\draw[dashed, popmuted] (0,0.5) -- (5,0.5);
\draw[popmuted] (3.6,0.5) arc (0:15.6:3.1);
\node[popmuted] at (3.35,0.72) {$\theta$};
% Car (simple box)
\draw[thick, popdark, fill=white, rotate around={15.6:(2.5,1.2)}] (1.9,1.2) rectangle (3.1,1.7);
\draw[fill=popink, rotate around={15.6:(2.5,1.2)}] (2.15,1.2) circle (0.09);
\draw[fill=popink, rotate around={15.6:(2.5,1.2)}] (2.85,1.2) circle (0.09);
% Forces
\draw[-latex, very thick, poporange] (2.5,1.55) -- ++(-0.42,1.55) node[above left] {$R$};
\draw[-latex, very thick, popgreen] (2.5,1.55) -- ++(0,-1.5) node[below right] {$mg$};
\draw[-latex, very thick, poppurple] (2.5,1.55) -- ++(-1.6,0) node[above] {$R\sin\theta$};
\end{tikzpicture}
\end{center}
\legende{Car on a banked bend (cross-section, centre of the circle to the left). The horizontal component $R\sin\theta$ of the normal reaction provides the centripetal force; the vertical component balances the weight.}
\end{popfigure}

\begin{exoresolu}[Car on a bend near Yaound\'e]
A car of mass $1200\ \mathrm{kg}$ rounds a bend of radius $50\ \mathrm{m}$.
\begin{enumerate}
\item[(a)] The road is flat and the coefficient of friction between tyres and road is $\mu = 0.8$. Find the maximum safe speed.
\item[(b)] The same bend is now banked at $12^{\circ}$ to the horizontal. Find the speed at which the car can take the bend with \emph{no} tendency to skid (no friction required). Take $g = 9.81\ \mathrm{m\,s^{-2}}$, $\tan 12^{\circ} = 0.2126$.
\end{enumerate}
\tcblower
\textbf{(a) Flat road.} Friction provides the centripetal force, and at the limit $F_f = \mu R = \mu mg$:
\[ \frac{mv_{\max}^2}{r} = \mu mg \quad\Longrightarrow\quad v_{\max} = \sqrt{\mu g r} = \sqrt{0.8 \times 9.81 \times 50} = \sqrt{392.4} \approx 19.8\ \mathrm{m\,s^{-1}}. \]
That is about $71\ \mathrm{km\,h^{-1}}$. Note that the mass cancels: the limit is the same for a small car or a loaded lorry.

\textbf{(b) Banked road.} Let $R$ be the normal reaction. Resolving:
\begin{align*}
\text{Vertically:} \quad & R\cos\theta = mg,\\
\text{Horizontally:} \quad & R\sin\theta = \frac{mv^2}{r}.
\end{align*}
Dividing the second equation by the first:
\[ \tan\theta = \frac{v^2}{gr} \quad\Longrightarrow\quad v = \sqrt{gr\tan\theta} = \sqrt{9.81 \times 50 \times 0.2126} = \sqrt{104.3} \approx 10.2\ \mathrm{m\,s^{-1}}. \]
At about $10.2\ \mathrm{m\,s^{-1}}$ ($\approx 37\ \mathrm{km\,h^{-1}}$) the car needs no friction at all; above this speed friction must act down the slope to supplement $R\sin\theta$.
\end{exoresolu}

\courssub{Method 5: Analysing the Conical Pendulum}

\begin{methode}[Setting up the conical pendulum equations]
A \cle{conical pendulum} is a small bob of mass $m$ on a light inextensible string of length $l$, moving in a \emph{horizontal} circle while the string traces a cone of semi-vertical angle $\theta$.
\begin{enumerate}
\item \textbf{Draw the forces:} only two — the tension $T$ along the string and the weight $mg$ vertically downwards.
\item \textbf{Identify the geometry:} the radius of the circle is $r = l\sin\theta$; the centre of the circle is vertically below the point of suspension, at depth $h = l\cos\theta$.
\item \textbf{Resolve the tension} into $T\cos\theta$ (vertical) and $T\sin\theta$ (horizontal, towards the centre).
\item \textbf{Write the two equations:}
\begin{align*}
\text{Vertical equilibrium:} \quad & T\cos\theta = mg,\\
\text{Horizontal (Newton II):} \quad & T\sin\theta = mr\omega^2 = \frac{mv^2}{r}.
\end{align*}
\item \textbf{Divide} to eliminate $T$:
\[ \tan\theta = \frac{r\omega^2}{g} = \frac{v^2}{gr}. \]
\item \textbf{Deduce the period:} substituting $r = l\sin\theta$ gives
\[ \omega^2 = \frac{g}{l\cos\theta} = \frac{g}{h}, \qquad T_{\text{period}} = 2\pi\sqrt{\frac{l\cos\theta}{g}} = 2\pi\sqrt{\frac{h}{g}}. \]
\end{enumerate}
\textbf{Key remark:} the period depends only on the \emph{vertical depth} $h = l\cos\theta$ of the bob below the support, not on the mass.
\end{methode}

\begin{popfigure}
\begin{center}
\begin{tikzpicture}[scale=1.1]
% Support
\fill[popink] (0,0) circle (0.05);
\draw[thick, popink] (-0.8,0) -- (0.8,0);
% String and bob
\draw[thick, popblue] (0,0) -- (2.12,-2.12);
\fill[poporange] (2.12,-2.12) circle (0.14);
\node[popdark] at (2.5,-2.0) {$m$};
% Vertical dashed line and angle
\draw[dashed, popmuted] (0,0) -- (0,-2.12);
\draw[popmuted] (0,-0.9) arc (-90:-45:0.9);
\node[popmuted] at (0.33,-0.95) {$\theta$};
% Depth h
\draw[latex-latex, popgreen] (-0.35,0) -- (-0.35,-2.12) node[midway, left] {$h$};
% Circle of motion (ellipse)
\draw[dashed, poppurple] (0,-2.12) ellipse (2.12 and 0.55);
% Radius r
\draw[latex-latex, poppurple] (0,-2.12) -- (2.12,-2.12) node[midway, above] {$r$};
% Forces on bob
\draw[-latex, very thick, popblue] (2.12,-2.12) -- ++(-1.06,1.06) node[above right, pos=0.85] {$T$};
\draw[-latex, very thick, popgreen] (2.12,-2.12) -- ++(0,-1.1) node[below] {$mg$};
\draw[-latex, very thick, poppurple] (2.12,-2.12) -- ++(-1.2,0) node[above left, pos=0.9] {$T\sin\theta$};
% Length label
\node[popblue] at (0.85,-0.8) {$l$};
\end{tikzpicture}
\end{center}
\legende{The conical pendulum: the bob moves in a horizontal circle of radius $r = l\sin\theta$ at depth $h = l\cos\theta$ below the support. Only two forces act: tension $T$ and weight $mg$; the horizontal component $T\sin\theta$ is the centripetal force.}
\end{popfigure}

\begin{exoresolu}[Period of a conical pendulum]
A small bob of mass $0.20\ \mathrm{kg}$ is attached to a light inextensible string of length $1.5\ \mathrm{m}$ fixed at a point $O$. The bob moves in a horizontal circle with constant angular speed, the string making an angle of $30^{\circ}$ with the vertical. Take $g = 9.81\ \mathrm{m\,s^{-2}}$.
\begin{enumerate}
\item[(a)] Find the tension in the string.
\item[(b)] Find the angular speed of the bob.
\item[(c)] Deduce the period of the motion.
\item[(d)] Find the linear speed of the bob.
\end{enumerate}
\tcblower
\textbf{(a) Tension.} Vertical equilibrium: $T\cos\theta = mg$, so
\[ T = \frac{mg}{\cos 30^{\circ}} = \frac{0.20 \times 9.81}{0.8660} = \frac{1.962}{0.8660} \approx 2.27\ \mathrm{N}. \]

\textbf{(b) Angular speed.} The radius is $r = l\sin\theta = 1.5 \times \sin 30^{\circ} = 1.5 \times 0.5 = 0.75\ \mathrm{m}$. Horizontal equation: $T\sin\theta = mr\omega^2$, hence
\[ \omega^2 = \frac{T\sin\theta}{mr} = \frac{2.27 \times 0.5}{0.20 \times 0.75} = \frac{1.133}{0.15} = 7.556, \qquad \omega \approx 2.75\ \mathrm{rad\,s^{-1}}. \]
\emph{Check with the shortcut:} $\omega = \sqrt{\dfrac{g}{l\cos\theta}} = \sqrt{\dfrac{9.81}{1.5 \times 0.8660}} = \sqrt{7.553} \approx 2.75\ \mathrm{rad\,s^{-1}}$. \checkmark

\textbf{(c) Period.}
\[ T_{\text{period}} = \frac{2\pi}{\omega} = \frac{2\pi}{2.75} \approx 2.29\ \mathrm{s}. \]
Equivalently, $T_{\text{period}} = 2\pi\sqrt{\dfrac{l\cos\theta}{g}} = 2\pi\sqrt{\dfrac{1.299}{9.81}} \approx 2.29\ \mathrm{s}$. \checkmark

\textbf{(d) Linear speed.}
\[ v = r\omega = 0.75 \times 2.75 \approx 2.06\ \mathrm{m\,s^{-1}}. \]
\emph{Check:} $\tan\theta = \dfrac{v^2}{gr}$ gives $v = \sqrt{gr\tan\theta} = \sqrt{9.81 \times 0.75 \times 0.5774} = \sqrt{4.248} \approx 2.06\ \mathrm{m\,s^{-1}}$. \checkmark
\end{exoresolu}

\courssub{Method 6: Comparing and Designing Conical Pendulums}

\begin{methode}[Using the relationships $\omega^2 = g/h$ and $v = \sqrt{gr\tan\theta}$]
For problems comparing or designing conical pendulums:
\begin{enumerate}
\item \textbf{Translate the geometry first:} $r = l\sin\theta$, $h = l\cos\theta$.
\item \textbf{Same depth, same period:} since $\omega^2 = \dfrac{g}{h}$, all conical pendulums whose bobs rotate in the \emph{same horizontal plane} have the same angular speed and period, whatever their lengths or masses.
\item \textbf{Increasing the speed} makes $\theta$ increase (since $\cos\theta = \dfrac{g}{l\omega^2}$): the bob rises, but $\theta$ can never reach $90^{\circ}$ in practice, because that would require infinite $\omega$.
\item \textbf{To find one quantity from another,} chain the relations:
\[ \cos\theta = \frac{g}{l\omega^2}, \qquad r = l\sin\theta, \qquad v = r\omega, \qquad T = \frac{mg}{\cos\theta} = ml\omega^2. \]
\item \textbf{Check consistency} with $\tan\theta = \dfrac{v^2}{gr}$.
\end{enumerate}
\end{methode}

\begin{exoresolu}[Designing a rotating swing ride]
At a funfair in Limbe, a ``swing carousel'' carries children in seats attached by chains of length $4.0\ \mathrm{m}$ to a rotating frame. A child with seat has total mass $40\ \mathrm{kg}$. The ride is intended to rotate so that each seat completes one revolution in $5.0\ \mathrm{s}$. Take $g = 9.81\ \mathrm{m\,s^{-2}}$.
\begin{enumerate}
\item[(a)] Calculate the angular speed corresponding to this intended period.
\item[(b)] Show that the angle $\theta$ of the chains to the vertical is independent of the mass of the child, and investigate whether the intended period is attainable.
\item[(c)] For an actual period of $3.5\ \mathrm{s}$, find the radius of the circle described by a seat and the linear speed of the seat.
\item[(d)] Calculate the tension in one chain for the $3.5\ \mathrm{s}$ period. What happens to $\theta$ if the rotation speeds up?
\end{enumerate}
\tcblower
\textbf{(a)} $\omega = \dfrac{2\pi}{T_{\text{period}}} = \dfrac{2\pi}{5.0} = 1.257\ \mathrm{rad\,s^{-1}}$.

\textbf{(b)} For a conical pendulum, $\cos\theta = \dfrac{g}{l\omega^2}$: the mass does not appear, so $\theta$ is the same for a light or a heavy child. Numerically, for the intended period:
\[ \cos\theta = \frac{9.81}{4.0 \times (1.257)^2} = \frac{9.81}{4.0 \times 1.579} = \frac{9.81}{6.317} = 1.553. \]
Since $\cos\theta > 1$ is impossible, the intended period is \emph{not attainable} with $4.0\ \mathrm{m}$ chains: the seats would simply hang vertically. The smallest possible period corresponds to $\theta \to 90^{\circ}$, i.e. $\omega \to \sqrt{g/l} = \sqrt{9.81/4.0} = 1.566\ \mathrm{rad\,s^{-1}}$, giving
\[ T_{\min} = \frac{2\pi}{1.566} \approx 4.01\ \mathrm{s}. \]
Any working period must therefore be \emph{shorter} than about $4.0\ \mathrm{s}$. Take instead a period of $3.5\ \mathrm{s}$: then $\omega = \dfrac{2\pi}{3.5} = 1.795\ \mathrm{rad\,s^{-1}}$ and
\[ \cos\theta = \frac{9.81}{4.0 \times (1.795)^2} = \frac{9.81}{12.89} = 0.761 \quad\Longrightarrow\quad \theta \approx 40.4^{\circ}. \]

\textbf{(c)} With $\theta = 40.4^{\circ}$:
\[ r = l\sin\theta = 4.0 \times \sin 40.4^{\circ} = 4.0 \times 0.648 = 2.59\ \mathrm{m}, \]
\[ v = r\omega = 2.59 \times 1.795 \approx 4.65\ \mathrm{m\,s^{-1}}. \]

\textbf{(d)} Tension:
\[ T = \frac{mg}{\cos\theta} = \frac{40 \times 9.81}{0.761} \approx 516\ \mathrm{N}. \]
(Equivalently $T = ml\omega^2 = 40 \times 4.0 \times 3.22 \approx 516\ \mathrm{N}$. \checkmark) If the ride speeds up, $\omega^2$ increases, so $\cos\theta = \dfrac{g}{l\omega^2}$ decreases and $\theta$ \emph{increases}: the seats rise higher. The tension $T = ml\omega^2$ also increases, which is why the chains must be strong and regularly inspected.
\end{exoresolu}

\courssub{Method 7: Synthesis — Multi-Step GCE Problems on Circular Motion}

\begin{methode}[Strategy for exam questions mixing several notions]
\begin{enumerate}
\item \textbf{Read twice and list the data} with symbols and SI units; convert rpm, km\,h$^{-1}$, cm.
\item \textbf{Draw a clear diagram:} circle, centre, radius, forces. Decide which force(s) point towards the centre.
\item \textbf{Write the two structural equations} (vertical balance; radial Newton's second law) before substituting numbers.
\item \textbf{Choose the right bridge:} $v = r\omega$, $\omega = \dfrac{2\pi}{T}$, $a = \dfrac{v^2}{r} = r\omega^2$.
\item \textbf{Conclude with a sentence} answering the question asked, with units and a sensible number of significant figures.
\end{enumerate}
\end{methode}

\begin{exoresolu}[GCE-style synthesis: the bucket of water]
A bucket of water of total mass $2.0\ \mathrm{kg}$ is whirled in a \emph{vertical} circle of radius $0.90\ \mathrm{m}$ at constant angular speed on the end of a rope. Take $g = 9.81\ \mathrm{m\,s^{-2}}$.
\begin{enumerate}
\item[(a)] At the top of the circle, write down the equation of motion of the bucket along the radius.
\item[(b)] Show that the minimum angular speed for the rope to remain taut at the top is $\omega_{\min} = \sqrt{g/r}$, and calculate its value.
\item[(c)] Find the corresponding minimum linear speed and the period of rotation at this limiting speed.
\item[(d)] Explain why the water does not fall out of the bucket at the top when the bucket moves faster than this minimum speed.
\end{enumerate}
\tcblower
\textbf{(a)} At the top of the circle, the centre is \emph{below} the bucket, so the centripetal direction is vertically downwards. The forces towards the centre are the weight $mg$ and the tension $T$ of the rope (both point down at the top). Newton's second law along the radius:
\[ T + mg = mr\omega^2. \]

\textbf{(b)} The rope remains taut as long as $T \ge 0$. The limiting case is $T = 0$:
\[ mg = mr\omega_{\min}^2 \quad\Longrightarrow\quad \omega_{\min}^2 = \frac{g}{r} \quad\Longrightarrow\quad \omega_{\min} = \sqrt{\frac{g}{r}}. \]
Numerically:
\[ \omega_{\min} = \sqrt{\frac{9.81}{0.90}} = \sqrt{10.9} \approx 3.30\ \mathrm{rad\,s^{-1}}. \]

\textbf{(c)} Minimum linear speed:
\[ v_{\min} = r\omega_{\min} = \sqrt{gr} = \sqrt{9.81 \times 0.90} = \sqrt{8.829} \approx 2.97\ \mathrm{m\,s^{-1}}. \]
Period at this speed:
\[ T_{\text{period}} = \frac{2\pi}{\omega_{\min}} = \frac{2\pi}{3.30} \approx 1.90\ \mathrm{s}. \]

\textbf{(d)} The water and the bucket both need a centripetal force $mr\omega^2$ to follow the circle. When $\omega > \omega_{\min}$, the required centripetal force exceeds the weight alone, so the bottom of the bucket must push on the water (normal reaction) to supply the difference: the water is pressed against the bottom of the bucket instead of falling. At exactly $\omega = \omega_{\min}$, gravity alone provides precisely the required centripetal force and the reaction is zero — the water is in ``free fall'' along the circle. Below $\omega_{\min}$, gravity provides \emph{more} than needed, the circular path cannot be maintained, and the water spills. This is the same principle used by pilots looping an aircraft and by the design of roller-coaster loops.
\end{exoresolu}

\begin{attention}[Final checklist before submitting your answer]
\begin{itemize}
\item Angles in radians, $\omega$ in $\mathrm{rad\,s^{-1}}$, distances in metres.
\item Centripetal acceleration is $a = \dfrac{v^2}{r} = r\omega^2$, directed towards the centre — never outwards.
\item The centripetal force is the \emph{resultant} towards the centre; identify the real forces providing it.
\item For the conical pendulum: $T\cos\theta = mg$, $T\sin\theta = mr\omega^2$, period $= 2\pi\sqrt{\dfrac{l\cos\theta}{g}}$.
\item Give final answers with units and $2$–$3$ significant figures.
\end{itemize}
\end{attention}

\newpage

\coursec{Exercises}

\courssub{I check my knowledge}

\begin{exercice}[Multiple choice: angular quantities]
A particle moves with uniform speed on a circle of radius $r$ with angular speed $\omega$. Choose the correct statement and justify briefly:
\begin{enumerate}
\item[(a)] The acceleration of the particle is zero because its speed is constant.
\item[(b)] The acceleration is directed along the tangent to the circle.
\item[(c)] The acceleration has magnitude $\omega^{2} r$ and is directed towards the centre of the circle.
\item[(d)] The acceleration has magnitude $\dfrac{v^{2}}{r}$ and is directed away from the centre.
\end{enumerate}
\end{exercice}

\begin{exercice}[True or false]
State whether each of the following statements is true or false, justifying your answer in one or two sentences.
\begin{enumerate}
\item In uniform circular motion, the velocity of the particle is constant.
\item The centripetal force is an extra force that must be added to the other forces acting on the body.
\item If the angular speed of a body moving on a circle is doubled while the radius is unchanged, the centripetal acceleration is multiplied by $4$.
\item For a conical pendulum, the period depends on the mass of the bob.
\end{enumerate}
\end{exercice}

\begin{exercice}[Definitions]
Define each of the following terms, giving its symbol and SI unit:
\begin{enumerate}
\item angular speed;
\item period of a circular motion;
\item centripetal acceleration;
\item conical pendulum.
\end{enumerate}
\end{exercice}

\begin{exercice}[Direct calculations]
A grinding wheel in a Douala workshop rotates at $1500\ \mathrm{rpm}$. Calculate:
\begin{enumerate}
\item its angular speed in $\mathrm{rad\,s^{-1}}$;
\item its period of rotation;
\item the linear speed of a point on the wheel situated $12\ \mathrm{cm}$ from the axis.
\end{enumerate}
\end{exercice}

\courssub{I practise}

\begin{exercice}[Stone whirled on a table]
A stone of mass $0.5\ \mathrm{kg}$ is whirled in a horizontal circle of radius $0.8\ \mathrm{m}$ on a frictionless table, at a rate of $2$ revolutions per second. The stone is attached to the centre by a light inextensible string. Find:
\begin{enumerate}
\item the angular speed of the stone;
\item its linear speed;
\item its centripetal acceleration;
\item the tension in the string.
\end{enumerate}
\end{exercice}

\begin{exercice}[Car rounding a bend]
A car of mass $900\ \mathrm{kg}$ rounds a flat (unbanked) bend of radius $50\ \mathrm{m}$. The coefficient of friction between the tyres and the road is $0.6$. Taking $g = 10\ \mathrm{m\,s^{-2}}$:
\begin{enumerate}
\item identify the force that provides the centripetal force;
\item find the maximum safe speed of the car round the bend;
\item state what happens if the driver exceeds this speed.
\end{enumerate}
\end{exercice}

\begin{exercice}[Belt-driven pulleys]
Two pulleys of radii $10\ \mathrm{cm}$ and $25\ \mathrm{cm}$ are connected by a belt that does not slip. The smaller pulley rotates at $300\ \mathrm{rpm}$.
\begin{enumerate}
\item Explain why the linear speed of the belt is the same at both pulleys.
\item Find the angular speed of the larger pulley in $\mathrm{rpm}$ and in $\mathrm{rad\,s^{-1}}$.
\item Calculate the linear speed of the belt.
\end{enumerate}
\end{exercice}

\begin{exercice}[Coin on a turntable]
A small coin is placed $20\ \mathrm{cm}$ from the centre of a horizontal turntable. The coefficient of static friction between the coin and the turntable is $0.5$. Taking $g = 10\ \mathrm{m\,s^{-2}}$, find:
\begin{enumerate}
\item the greatest angular speed, in $\mathrm{rad\,s^{-1}}$, for which the coin does not slip;
\item this angular speed expressed in revolutions per minute.
\end{enumerate}
\end{exercice}

\courssub{I prepare for the GCE}

\begin{exercice}[The conical pendulum: structured question]
A conical pendulum consists of a bob of mass $0.4\ \mathrm{kg}$ attached to a light inextensible string of length $1.2\ \mathrm{m}$. The bob describes a horizontal circle with the string making a constant angle of $37^{\circ}$ with the vertical. Take $\sin 37^{\circ} = 0.6$, $\cos 37^{\circ} = 0.8$ and $g = 10\ \mathrm{m\,s^{-2}}$.
\begin{enumerate}
\item Draw a clear diagram showing all the forces acting on the bob. \hfill (2 marks)
\item By resolving forces vertically, calculate the tension in the string. \hfill (3 marks)
\item Write down the equation of motion in the horizontal (radial) direction. \hfill (2 marks)
\item Hence calculate the angular speed of the bob. \hfill (3 marks)
\item Deduce the linear speed of the bob and the period of the motion. \hfill (4 marks)
\item State one assumption made about the string and one about air resistance. \hfill (2 marks)
\end{enumerate}
\end{exercice}

\begin{exercice}[Derivation of the period formula]
\begin{enumerate}
\item Starting from a diagram of a conical pendulum of string length $l$, with the bob of mass $m$ describing a horizontal circle while the string makes an angle $\theta$ with the vertical, derive from first principles the relation
\[ \omega^{2} = \frac{g}{l\cos\theta}, \]
stating clearly each step and the forces considered. \hfill (6 marks)
\item Deduce that the period of the conical pendulum is
\[ T = 2\pi\sqrt{\frac{h}{g}}, \]
where $h$ is the vertical distance from the point of suspension to the plane of the circle. \hfill (3 marks)
\item Two conical pendulums of lengths $1.0\ \mathrm{m}$ and $1.5\ \mathrm{m}$ rotate with the same angular speed. Using the relation of part (a), determine which pendulum makes the larger angle with the vertical, justifying your answer carefully. \hfill (4 marks)
\item Show that the linear speed of the bob can be written $v = \sqrt{gl\sin\theta\tan\theta}$, and hence find $v$ when $l = 2.0\ \mathrm{m}$ and $\theta = 30^{\circ}$, taking $g = 9.81\ \mathrm{m\,s^{-2}}$. \hfill (5 marks)
\end{enumerate}
\end{exercice}

\begin{exercice}[Circular orbits: extension problem]
A satellite of mass $m$ moves with uniform speed $v$ in a circular orbit of radius $r$ about the Earth.
\begin{enumerate}
\item Write down an expression for the centripetal force acting on the satellite, and explain what physical agent provides this force. \hfill (3 marks)
\item At the altitude of the orbit, the gravitational force per unit mass is $8.7\ \mathrm{N\,kg^{-1}}$ and $r = 6.9\times 10^{6}\ \mathrm{m}$. Show that the orbital speed is approximately $7.7\ \mathrm{km\,s^{-1}}$. \hfill (4 marks)
\item Calculate the period of the orbit, giving your answer in minutes. \hfill (4 marks)
\item Explain why the satellite can remain in orbit although it is constantly accelerating towards the Earth. \hfill (3 marks)
\item The satellite is transferred to an orbit of larger radius. Without calculation, state and justify whether its angular speed increases or decreases. \hfill (3 marks)
\end{enumerate}
\end{exercice}

\newpage

\coursec{Solutions to the exercises}

\begin{corrige}[Exercise 1 — Multiple choice: angular quantities]
The correct statement is \textbf{(c)}.

In uniform circular motion the \emph{speed} is constant but the \emph{velocity} is not, because its direction changes continuously. A changing velocity means a non-zero acceleration, so (a) is false. This acceleration is the \cle{centripetal acceleration}: it is directed along the radius \emph{towards the centre} (not along the tangent, not away from the centre), so (b) and (d) are false. Its magnitude is
\[ a = \omega^{2} r \qquad \Bigl(\text{equivalently } a = \frac{v^{2}}{r}\Bigr), \]
which is exactly statement (c). Note that $\dfrac{v^{2}}{r}$ in (d) has the correct magnitude but the wrong direction, so (d) is rejected.
\end{corrige}

\begin{corrige}[Exercise 2 — True or false]
\begin{enumerate}
\item \textbf{False.} The speed is constant, but the velocity vector changes direction at every instant (it is always tangent to the circle). Hence the velocity is not constant — this is precisely why an acceleration exists.
\item \textbf{False.} The centripetal force is not an additional force; it is the \emph{resultant} of the real forces (tension, friction, weight, reaction\ldots) directed towards the centre. One must never add a ``centripetal force'' to a force diagram.
\item \textbf{True.} Since $a = \omega^{2}r$ with $r$ fixed, $a$ is proportional to $\omega^{2}$. Doubling $\omega$ multiplies $a$ by $2^{2} = 4$.
\item \textbf{False.} For a conical pendulum $T = 2\pi\sqrt{\dfrac{l\cos\theta}{g}} = 2\pi\sqrt{\dfrac{h}{g}}$: the mass $m$ cancels when the vertical equilibrium ($T\cos\theta = mg$) is combined with the radial equation. The period depends only on $l$, $\theta$ (or $h$) and $g$.
\end{enumerate}
\end{corrige}

\begin{corrige}[Exercise 3 — Definitions]
\begin{enumerate}
\item \textbf{Angular speed} $\omega$: the angle swept out per unit time by the radius joining the particle to the centre, $\omega = \dfrac{\theta}{t}$ (constant for uniform circular motion). SI unit: $\mathrm{rad\,s^{-1}}$.
\item \textbf{Period} $T$: the time taken by the particle to complete exactly one full revolution. SI unit: the second ($\mathrm{s}$). It satisfies $\omega = \dfrac{2\pi}{T}$.
\item \textbf{Centripetal acceleration} $\vec{a}$: the acceleration of a particle in circular motion, directed along the radius towards the centre, of magnitude $a = \dfrac{v^{2}}{r} = \omega^{2}r$. SI unit: $\mathrm{m\,s^{-2}}$.
\item \textbf{Conical pendulum}: a system consisting of a small bob of mass $m$ attached to a fixed point by a light inextensible string of length $l$, the bob describing a horizontal circle with uniform speed so that the string sweeps out the surface of a cone, making a constant angle $\theta$ with the vertical.
\end{enumerate}
\end{corrige}

\begin{corrige}[Exercise 4 — Direct calculations]
Given: $1500\ \mathrm{rpm}$, i.e. frequency $f = \dfrac{1500}{60} = 25\ \mathrm{Hz}$; radius $r = 12\ \mathrm{cm} = 0.12\ \mathrm{m}$.
\begin{enumerate}
\item Angular speed:
\[ \omega = 2\pi f = 2\pi \times 25 = 50\pi \approx 157.1\ \mathrm{rad\,s^{-1}}. \]
\item Period:
\[ T = \frac{1}{f} = \frac{1}{25} = 0.04\ \mathrm{s}. \]
\item Linear speed of a point at $r = 0.12\ \mathrm{m}$:
\[ v = \omega r = 50\pi \times 0.12 = 6\pi \approx 18.8\ \mathrm{m\,s^{-1}}. \]
\end{enumerate}
\end{corrige}

\begin{corrige}[Exercise 5 — Stone whirled on a table]
Given: $m = 0.5\ \mathrm{kg}$, $r = 0.8\ \mathrm{m}$, $f = 2\ \mathrm{rev\,s^{-1}}$.
\begin{enumerate}
\item Angular speed:
\[ \omega = 2\pi f = 2\pi \times 2 = 4\pi \approx 12.6\ \mathrm{rad\,s^{-1}}. \]
\item Linear speed:
\[ v = \omega r = 4\pi \times 0.8 = 3.2\pi \approx 10.1\ \mathrm{m\,s^{-1}}. \]
\item Centripetal acceleration:
\[ a = \omega^{2} r = (4\pi)^{2}\times 0.8 = 12.8\pi^{2} \approx 126.3\ \mathrm{m\,s^{-2}}. \]
\item On the frictionless table, the weight and the normal reaction cancel vertically; the \cle{tension} alone provides the centripetal force:
\[ T = m\omega^{2}r = 0.5 \times 12.8\pi^{2} = 6.4\pi^{2} \approx 63.2\ \mathrm{N}. \]
\end{enumerate}
\end{corrige}

\begin{corrige}[Exercise 6 — Car rounding a bend]
Given: $m = 900\ \mathrm{kg}$, $r = 50\ \mathrm{m}$, $\mu = 0.6$, $g = 10\ \mathrm{m\,s^{-2}}$.
\begin{enumerate}
\item The bend is flat, so the normal reaction is vertical and balances the weight. The only horizontal force available is the \cle{friction} between the tyres and the road: it is friction that provides the centripetal force, directed towards the centre of the bend.
\item Vertically: $R = mg$. The maximum friction is $F_{\max} = \mu R = \mu mg$. At the limiting speed $v_{\max}$:
\[ \frac{mv_{\max}^{2}}{r} = \mu mg \quad\Longrightarrow\quad v_{\max} = \sqrt{\mu g r} = \sqrt{0.6\times 10\times 50} = \sqrt{300} \approx 17.3\ \mathrm{m\,s^{-1}}. \]
(About $62\ \mathrm{km\,h^{-1}}$.) Note that the mass cancels: the limiting speed does not depend on $m$.
\item If the driver exceeds $v_{\max}$, the required centripetal force $\dfrac{mv^{2}}{r}$ exceeds the maximum available friction: the tyres lose grip and the car \emph{skids outwards}, leaving the circular path (it moves off along a curve of larger radius, tangentially outwards).
\end{enumerate}
\end{corrige}

\begin{corrige}[Exercise 7 — Belt-driven pulleys]
Given: $r_{1} = 10\ \mathrm{cm} = 0.10\ \mathrm{m}$, $r_{2} = 25\ \mathrm{cm} = 0.25\ \mathrm{m}$, small pulley at $300\ \mathrm{rpm}$.
\begin{enumerate}
\item The belt does not slip, so every point of the belt has the same linear speed $v$. At the rim of each pulley, the point of contact must move with the belt (no relative motion), hence the rim speeds of both pulleys equal the belt speed: $v_{1} = v_{2} = v$.
\item Equality of rim speeds gives $\omega_{1}r_{1} = \omega_{2}r_{2}$, so
\[ \omega_{2} = \omega_{1}\frac{r_{1}}{r_{2}} = 300\times\frac{10}{25} = 120\ \mathrm{rpm}. \]
In $\mathrm{rad\,s^{-1}}$: $\omega_{2} = 120\times\dfrac{2\pi}{60} = 4\pi \approx 12.6\ \mathrm{rad\,s^{-1}}$.
\item Belt speed (using the small pulley, $\omega_{1} = 300\times\frac{2\pi}{60} = 10\pi\ \mathrm{rad\,s^{-1}}$):
\[ v = \omega_{1}r_{1} = 10\pi\times 0.10 = \pi \approx 3.14\ \mathrm{m\,s^{-1}}. \]
(Check with the large pulley: $4\pi\times 0.25 = \pi\ \mathrm{m\,s^{-1}}$ — consistent.)
\end{enumerate}
\end{corrige}

\begin{corrige}[Exercise 8 — Coin on a turntable]
Given: $r = 20\ \mathrm{cm} = 0.20\ \mathrm{m}$, $\mu = 0.5$, $g = 10\ \mathrm{m\,s^{-2}}$.
\begin{enumerate}
\item Vertically the coin is in equilibrium: $R = mg$. The centripetal force is supplied by static friction, $F = m\omega^{2}r$, and slipping begins when $F$ reaches $\mu R = \mu mg$:
\[ m\omega_{\max}^{2}r = \mu mg \quad\Longrightarrow\quad \omega_{\max} = \sqrt{\frac{\mu g}{r}} = \sqrt{\frac{0.5\times 10}{0.20}} = \sqrt{25} = 5\ \mathrm{rad\,s^{-1}}. \]
\item Converting to revolutions per minute:
\[ N = \omega_{\max}\times\frac{60}{2\pi} = 5\times\frac{60}{2\pi} = \frac{150}{\pi} \approx 47.7\ \mathrm{rpm}. \]
\end{enumerate}
\end{corrige}

\begin{corrige}[Exercise 9 — The conical pendulum: structured question]
Given: $m = 0.4\ \mathrm{kg}$, $l = 1.2\ \mathrm{m}$, $\theta = 37^{\circ}$, $\sin 37^{\circ} = 0.6$, $\cos 37^{\circ} = 0.8$, $g = 10\ \mathrm{m\,s^{-2}}$.

\textbf{Indicative mark scheme} (total 16 marks) is given within each part.
\begin{enumerate}
\item \textbf{Diagram (2 marks).} 1 mark for the two forces only: tension $T$ along the string towards the support and weight $mg$ vertically downwards; 1 mark for showing the horizontal circle, the angle $\theta$ with the vertical and the centre $O$ of the circle. No ``centripetal force'' may be drawn as a separate force.
\begin{popfigure}
\begin{tikzpicture}[scale=1.1]
\coordinate (P) at (0,0);
\coordinate (B) at (2.4,-3.2);
\draw[dashed,popmuted] (P) -- (0,-3.2);
\draw[dashed,popmuted] (0,-3.2) ellipse (2.4 and 0.55);
\draw[thick,popblue] (P) -- (B);
\fill[poppink] (B) circle (0.14);
\draw[->,thick,popdark] (B) -- ++(0,-1.3) node[below] {$mg$};
\draw[->,thick,popgreen] (B) -- ++(-1.0,1.33) node[above left] {$T$};
\draw (0,-0.9) arc (-90:-53:0.9) node[midway,below right] {$\theta$};
\fill (0,-3.2) circle (0.05) node[below] {$O$};
\node at (0.25,0.1) {$S$};
\end{tikzpicture}
\legende{Forces on the bob of a conical pendulum.}
\end{popfigure}
\item \textbf{Vertical resolution (3 marks).} The bob has no vertical acceleration, so (M1 for resolving):
\[ T\cos\theta = mg \quad\Longrightarrow\quad T = \frac{mg}{\cos\theta} = \frac{0.4\times 10}{0.8} = 5\ \mathrm{N}. \]
(A1 for the equation, A1 for the value.)
\item \textbf{Radial equation (2 marks).} The horizontal component of the tension provides the centripetal force, with $r = l\sin\theta$:
\[ T\sin\theta = m\omega^{2}r = m\omega^{2}l\sin\theta. \]
(M1 for identifying $T\sin\theta$ as the radial resultant, A1 for $r = l\sin\theta$.)
\item \textbf{Angular speed (3 marks).} Dividing by $m\sin\theta$ (M1):
\[ \omega^{2} = \frac{T}{ml} = \frac{5}{0.4\times 1.2} = \frac{5}{0.48} \approx 10.42 \quad\Longrightarrow\quad \omega \approx 3.23\ \mathrm{rad\,s^{-1}}. \]
(M1 for substitution, A1 for the value. Equivalently $\omega = \sqrt{g/(l\cos\theta)} = \sqrt{10/0.96}$.)
\item \textbf{Linear speed and period (4 marks).} Radius: $r = l\sin\theta = 1.2\times 0.6 = 0.72\ \mathrm{m}$ (B1).
\[ v = \omega r = 3.23\times 0.72 \approx 2.32\ \mathrm{m\,s^{-1}} \quad (\text{A1}); \qquad T_{\text{period}} = \frac{2\pi}{\omega} = \frac{2\pi}{3.23} \approx 1.95\ \mathrm{s} \quad (\text{M1 A1}). \]
\item \textbf{Assumptions (2 marks, 1 each).} The string is \emph{light} (negligible mass) and \emph{inextensible} (constant length); \emph{air resistance is neglected} (the motion continues at constant angular speed).
\end{enumerate}
\end{corrige}

\begin{corrige}[Exercise 10 — Derivation of the period formula]
\begin{enumerate}
\item \textbf{Derivation (6 marks).} Consider the bob of mass $m$ describing a horizontal circle of radius $r = l\sin\theta$ with uniform angular speed $\omega$. The forces on the bob are \emph{only}: its weight $mg$ (vertically downwards) and the tension $T$ of the string (along the string towards the point of suspension). \textbf{(B1: correct force diagram; B1: statement of the two forces.)}

Since the bob stays in a horizontal plane, its vertical acceleration is zero. Resolving vertically:
\[ T\cos\theta = mg. \tag{1} \quad \textbf{(M1)} \]
The resultant in the horizontal (radial) direction is the component $T\sin\theta$, directed towards the centre; by Newton's second law with centripetal acceleration $\omega^{2}r$:
\[ T\sin\theta = m\omega^{2}r = m\omega^{2}l\sin\theta. \tag{2} \quad \textbf{(M1 A1)} \]
Dividing (2) by (1):
\[ \tan\theta = \frac{\omega^{2}l\sin\theta}{g} \quad\Longrightarrow\quad \frac{\sin\theta}{\cos\theta} = \frac{\omega^{2}l\sin\theta}{g}, \]
and since $\sin\theta \neq 0$ (the pendulum actually rotates), we cancel $\sin\theta$:
\[ \boxed{\ \omega^{2} = \frac{g}{l\cos\theta}\ } \quad \textbf{(A1)}. \]
\item \textbf{Period (3 marks).} The vertical distance from the suspension point to the plane of the circle is $h = l\cos\theta$ \textbf{(B1)}. Hence $\omega^{2} = \dfrac{g}{h}$, and since $\omega = \dfrac{2\pi}{T}$:
\[ \frac{4\pi^{2}}{T^{2}} = \frac{g}{h} \quad\Longrightarrow\quad \boxed{\ T = 2\pi\sqrt{\frac{h}{g}}\ } \quad \textbf{(M1 A1)}. \]
The period depends only on the \emph{height} $h$, not on $m$, $l$ or $\theta$ separately.
\item \textbf{Comparison (4 marks).} From $\omega^{2} = \dfrac{g}{l\cos\theta}$ with the same $\omega$ for both pendulums:
\[ l_{1}\cos\theta_{1} = l_{2}\cos\theta_{2} = \frac{g}{\omega^{2}} \quad\Longrightarrow\quad \cos\theta = \frac{g}{\omega^{2}l}. \quad \textbf{(M1)} \]
Thus $\cos\theta$ is a \emph{decreasing} function of $l$: the longer pendulum ($l = 1.5\ \mathrm{m}$) has the \emph{smaller} $\cos\theta$. Since cosine decreases on $[0, 90^{\circ}]$, a smaller $\cos\theta$ means a \emph{larger} angle $\theta$ \textbf{(A1, A1 for the reasoning)}. Conclusion: the pendulum of length $1.5\ \mathrm{m}$ makes the larger angle with the vertical \textbf{(A1)}.
\item \textbf{Linear speed (5 marks).} From (2): $m\omega^{2}l\sin\theta = T\sin\theta$, and with (1), $T = \dfrac{mg}{\cos\theta}$, so $\omega^{2} = \dfrac{g}{l\cos\theta}$ \textbf{(M1)}. Then
\[ v = \omega r = \omega l\sin\theta \quad\Longrightarrow\quad v^{2} = \omega^{2}l^{2}\sin^{2}\theta = \frac{g\,l^{2}\sin^{2}\theta}{l\cos\theta} = gl\,\frac{\sin^{2}\theta}{\cos\theta} = gl\sin\theta\tan\theta. \quad \textbf{(M1 A1)} \]
Hence $\boxed{v = \sqrt{gl\sin\theta\tan\theta}}$.

Numerically, with $l = 2.0\ \mathrm{m}$, $\theta = 30^{\circ}$, $\sin 30^{\circ} = 0.5$, $\tan 30^{\circ} = \dfrac{1}{\sqrt{3}} \approx 0.5774$, $g = 9.81\ \mathrm{m\,s^{-2}}$:
\[ v = \sqrt{9.81\times 2.0\times 0.5\times 0.5774} = \sqrt{5.664} \approx 2.38\ \mathrm{m\,s^{-1}}. \quad \textbf{(M1 A1)} \]
\end{enumerate}
\end{corrige}

\begin{corrige}[Exercise 11 — Circular orbits: extension problem]
\begin{enumerate}
\item \textbf{Centripetal force (3 marks).} $F = \dfrac{mv^{2}}{r}$ \textbf{(B1)}, directed towards the centre of the Earth \textbf{(B1)}. It is provided entirely by the \cle{gravitational attraction} of the Earth on the satellite \textbf{(B1)} — there is no engine or string.
\item \textbf{Orbital speed (4 marks).} The gravitational force per unit mass is $g' = 8.7\ \mathrm{N\,kg^{-1}}$, so the gravitational force on the satellite is $mg'$. Equating this to the centripetal force \textbf{(M1)}:
\[ mg' = \frac{mv^{2}}{r} \quad\Longrightarrow\quad v = \sqrt{g'r} \quad \textbf{(M1)}. \]
\[ v = \sqrt{8.7\times 6.9\times 10^{6}} = \sqrt{6.003\times 10^{7}} \approx 7.75\times 10^{3}\ \mathrm{m\,s^{-1}} \approx 7.7\ \mathrm{km\,s^{-1}}. \quad \textbf{(A1 A1)} \]
\item \textbf{Period (4 marks).} 
\[ T = \frac{2\pi r}{v} \quad \textbf{(M1)} = \frac{2\pi\times 6.9\times 10^{6}}{7.75\times 10^{3}} \approx 5.59\times 10^{3}\ \mathrm{s} \quad \textbf{(A1)}. \]
Converting: $T \approx \dfrac{5590}{60} \approx 93\ \mathrm{min}$ \textbf{(M1 A1)} — about one and a half hours, typical of a low Earth orbit.
\item \textbf{Why it stays in orbit (3 marks).} The acceleration is \emph{perpendicular} to the velocity at every instant, so it changes only the \emph{direction} of the velocity, not its magnitude \textbf{(B1)}. The satellite is in permanent ``free fall'' towards the Earth, but its large tangential speed means that as it falls, the Earth's surface curves away beneath it \textbf{(B1)}: it continually falls \emph{around} the Earth, staying at the same radius \textbf{(B1)}.
\item \textbf{Larger orbit (3 marks).} From part (b), $v^{2} = g'r$ with $g'$ the gravitational field at radius $r$. More fundamentally, gravity supplies $m\omega^{2}r$, and the gravitational force \emph{decreases} as $r$ increases while the required centripetal force per unit mass $\omega^{2}r$ would increase if $\omega$ stayed constant \textbf{(M1)}. Equating gravitational force to $m\omega^{2}r$ shows $\omega^{2} = \dfrac{g'}{r}$ decreases (both because $g'$ decreases and $r$ increases) \textbf{(M1)}. Hence the angular speed \emph{decreases} in a larger orbit \textbf{(A1)} — distant satellites revolve more slowly.
\end{enumerate}
\end{corrige}

\newpage

\coursec{Summary sheet}

\begin{popfigure}
\begin{tikzpicture}[mindmap, grow cyclic,
  every node/.style={concept, font=\small},
  root concept/.append style={concept color=popblue, font=\bfseries\small, minimum size=2.9cm},
  level 1 concept/.append style={sibling angle=60, level distance=3.4cm, minimum size=2.2cm, font=\scriptsize}]
\node[root concept]{Uniform Circular Motion}
  child[concept color=poporange]{ node {Angular quantities: $\theta$, $\omega$, $T$, $f$} }
  child[concept color=popgreen]{ node {Centripetal acceleration $a=\dfrac{v^{2}}{r}=r\omega^{2}$} }
  child[concept color=poppurple]{ node {Centripetal force $F=\dfrac{mv^{2}}{r}=mr\omega^{2}$} }
  child[concept color=popgold]{ node {Linear vs angular speed $v=r\omega$} }
  child[concept color=poppink]{ node {Applications: vehicles, banking, satellites} }
  child[concept color=popmuted]{ node {Conical pendulum: $T=2\pi\sqrt{\dfrac{h}{g}}$} };
\end{tikzpicture}
\legende{Mind map of the chapter: uniform circular motion and the conical pendulum.}
\end{popfigure}

\begin{bilan}
\textbf{1. Key definitions}
\begin{itemize}
\item \cle{Circular motion}: motion of a particle along a circular path of radius $r$; it is \emph{uniform} when the speed $v$ is constant.
\item \cle{Angular displacement} $\theta$: angle swept by the radius, measured in \textbf{radians} ($2\pi\ \mathrm{rad}=360^{\circ}$).
\item \cle{Angular speed} $\omega=\dfrac{\theta}{t}=\dfrac{2\pi}{T}=2\pi f$ (in $\mathrm{rad\,s^{-1}}$), where $T$ is the \cle{period} and $f$ the \cle{frequency} (in Hz).
\item \cle{Linear speed} $v$: distance travelled per unit time along the arc; $v=\dfrac{2\pi r}{T}$.
\item \cle{Conical pendulum}: a particle attached to a string whose other end is fixed, moving in a \emph{horizontal} circle so that the string sweeps out a cone of semi-vertical angle $\alpha$.
\end{itemize}

\textbf{2. Laws, formulas and theorems to know}
\begin{itemize}
\item Arc length: $s=r\theta$ ($\theta$ in radians).
\item Relationship between linear and angular speed:
\[ v=r\omega \qquad\text{hence}\qquad v=\frac{2\pi r}{T}=2\pi r f. \]
\item \cle{Centripetal acceleration} (directed towards the centre):
\[ a=\frac{v^{2}}{r}=r\omega^{2}=v\omega=\frac{4\pi^{2}r}{T^{2}}. \]
\item \cle{Centripetal force} (Newton's second law applied radially):
\[ F=ma=\frac{mv^{2}}{r}=mr\omega^{2}. \]
It is \emph{not} a new force: it is the \textbf{resultant} of real forces (tension, friction, weight component, reaction) directed towards the centre.
\item Conical pendulum (string of length $l$, semi-vertical angle $\alpha$, height of cone $h=l\cos\alpha$):
\begin{itemize}
\item Vertical equilibrium: $T\cos\alpha=mg$ (here $T$ = tension).
\item Radial equation: $T\sin\alpha=\dfrac{mv^{2}}{r}=mr\omega^{2}$, with $r=l\sin\alpha$.
\item Angular speed: $\omega=\sqrt{\dfrac{g}{h}}=\sqrt{\dfrac{g}{l\cos\alpha}}$.
\item \textbf{Period}: $\displaystyle T=2\pi\sqrt{\frac{h}{g}}=2\pi\sqrt{\frac{l\cos\alpha}{g}}$; linear speed: $v=r\omega=\sqrt{gr\tan\alpha}$.
\end{itemize}
\item Banking of a track (no friction): $\tan\alpha=\dfrac{v^{2}}{rg}$.
\end{itemize}

\textbf{3. Methods}
\begin{itemize}
\item \textbf{To solve any circular motion problem:} (1) draw the force diagram; (2) resolve vertically (equilibrium if no vertical motion) and horizontally towards the centre (Newton's second law $=mr\omega^{2}$); (3) eliminate unknowns.
\item \textbf{To convert units:} rev/min $\rightarrow$ rad/s: multiply by $\dfrac{2\pi}{60}$; km/h $\rightarrow$ m/s: divide by $3.6$.
\item \textbf{Conical pendulum:} identify $h$ (vertical distance from support to the plane of the circle), then apply $\omega=\sqrt{g/h}$.
\end{itemize}

\textbf{4. Common mistakes}
\begin{itemize}
\item Mixing degrees and radians: \textbf{all} formulas ($s=r\theta$, $v=r\omega$) require \textbf{radians}.
\item Believing a ``centrifugal force'' pushes the particle outwards: in an inertial frame only the inward centripetal resultant acts.
\item Writing $F=mv^{2}/r$ as an \emph{extra} force on the diagram: it is the \emph{resultant}, never an additional force.
\item Confusing $T$ (period, in seconds) with $T$ (tension, in newtons) -- use different symbols.
\item For the conical pendulum, using $l$ instead of $h=l\cos\alpha$ in $T=2\pi\sqrt{h/g}$, or $l$ instead of $r=l\sin\alpha$ in $v=r\omega$.
\item Forgetting that constant \emph{speed} does not mean zero acceleration: the \emph{velocity} changes direction continuously.
\end{itemize}

\textbf{5. GCE tips}
\begin{itemize}
\item Quote $v=r\omega$ and $a=v^{2}/r$ before substituting; method marks are awarded for correct formulas.
\item Show the force diagram clearly: examiners award marks for correct resolution (vertical / radial).
\item Keep $g=\SI{9.81}{m.s^{-2}}$ (or $9.8$ as instructed) and give the final answer to 3 significant figures with units.
\item Typical GCE items: car on a banked bend, stone on a string, satellite in circular orbit, conical pendulum period -- practise deriving $T=2\pi\sqrt{l\cos\alpha/g}$ from scratch; the derivation itself is examinable.
\item Check coherence: larger $r$ at fixed $\omega$ $\Rightarrow$ larger $v$ and larger $a$; larger $\alpha$ in a conical pendulum $\Rightarrow$ shorter period.
\end{itemize}
\end{bilan}

\findecours

\end{document}
